% Lesson 18 — Vocabulary: Words and expressions related to banking services — Anglais Tle A — Cours signé OnBuch+
% Programme officiel MINESEC. Compiler avec : tectonic EN18-vocabulary-words-and-expressions-related-to-banking-ser.tex
\newcommand{\DOCMATIERE}{Anglais}
\newcommand{\DOCNIVEAU}{Tle A}
\newcommand{\DOCMODULE}{Chapter 2 : Economic Life and Occupations}
\newcommand{\DOCLECON}{EN18}
\newcommand{\DOCTITRE}{Lesson 18 — Vocabulary: Words and expressions related to banking services}
% OnBuch+ — préambule « cours signés » ENRICHI (compilé avec tectonic / XeLaTeX)
% Système visuel « Pop » d'OnBuch+ : fond crème (jamais blanc), bordures encre
% épaisses, ombres « dures » décalées, titres Archivo Black en capitales.
% Version MATHÉMATIQUES : graphes (pgfplots/tikz), boîtes pédagogiques
% (définition, propriété, méthode, attention, le savais-tu, exercice résolu,
% exercice, TP, bilan), page de garde et signature OnBuch+ ancrée en bas.
%
% Macros attendues dans main.tex AVANT \input{preamble} :
%   \DOCMATIERE  \DOCNIVEAU  \DOCMODULE  \DOCLECON (numéro)  \DOCTITRE
\documentclass[11pt]{article}

\usepackage{fontspec}
\usepackage[english,french]{babel} % français = langue principale ; l'anglais via \eng{..} / textebox
\usepackage[a4paper,margin=2cm,top=3.4cm,bottom=2.8cm,headheight=1.4cm,footskip=1.3cm]{geometry}
\usepackage{amsmath,amssymb,amsfonts}
\usepackage{mathtools} % \xmapsto, \coloneqq... souvent utilisés par les modèles
\usepackage{cancel} % simplifications barrées (bilans de forces)
% \abs{z} (module, valeur absolue) et \norm{u} (norme), utilisés par les modèles
\providecommand{\abs}{}\let\abs\relax\DeclarePairedDelimiter{\abs}{\lvert}{\rvert}
\providecommand{\norm}{}\let\norm\relax\DeclarePairedDelimiter{\norm}{\lVert}{\rVert}
\usepackage[table]{xcolor} % [table] : fournit aussi \rowcolors (alternance de lignes)
\usepackage{pagecolor}
\usepackage{tikz}
\usetikzlibrary{arrows.meta,positioning,calc,shapes.geometric,shapes.misc,shapes.arrows,decorations.pathmorphing,decorations.pathreplacing,decorations.markings,patterns,shadows,mindmap,trees,fit,backgrounds,matrix,intersections,angles,quotes}
\usepackage{pgfplots}
\usepackage[version=4]{mhchem} % équations nucléaires \ce{^{235}_{92}U}
\usepackage[european,siunitx]{circuitikz} % schémas électriques
\pgfplotsset{compat=1.18}
\usepgfplotslibrary{fillbetween,groupplots} % aires entre courbes (calcul intégral)
\usepackage{enumitem}
\usepackage{fancyhdr}
% [most] charge toutes les bibliothèques tcolorbox (listings, minted, raster,
% documentation...) et gonfle inutilement la mémoire TeX sur un document long
% (~300 boîtes) : on ne charge que ce qui est réellement utilisé.
\usepackage[skins,breakable]{tcolorbox}
\usepackage{graphicx}
\usepackage{colortbl}
\usepackage{booktabs}
\usepackage{tabularx}
\usepackage{diagbox} % case d'en-tête diagonale des tableaux à double entrée
% Colonnes extensibles centrée / gauche / droite (C, L, R), souvent utilisées par les modèles
\newcolumntype{C}{>{\centering\arraybackslash}X}
\newcolumntype{L}{>{\raggedright\arraybackslash}X}
\newcolumntype{R}{>{\raggedleft\arraybackslash}X}
\usepackage{array}
\usepackage{multicol}
\usepackage{multirow}
\usepackage{siunitx}
% parse-numbers=false : accepte aussi \SI{2,5 \times 10^{-3}}{...}
\sisetup{locale=FR,output-decimal-marker={,},group-minimum-digits=4,parse-numbers=false}
\DeclareSIUnit\ohm{\ensuremath{\symup{\Omega}}} % U+2126 absent de la police
\DeclareSIUnit\division{div} % réglages d'oscilloscope (ms/div, V/div)
\DeclareSIUnit\tr{tr} % tours (vitesse de rotation, ex. \SI{1500}{\tr\per\minute})
\DeclareSIUnit\molar{M} % concentration molaire (ex. \SI{3}{\molar}, \SI{1}{\micro\molar})
\DeclareSIUnit\an{an} % année (le modèle confond parfois avec \year, un compteur TeX) — ex. \SI{5}{\kilo\meter\per\an}
\DeclareSIUnit\FCFA{FCFA} % franc CFA (ex. \SI{500}{\FCFA\per\kilo\gram})
\DeclareSIUnit\degreeBrix{\textdegree Brix} % teneur en sucre d'un sirop/jus (réfractométrie)
\DeclareSIUnit\month{mois} % le modèle utilise parfois \month, un compteur TeX, comme unité de durée
\usepackage{qrcode}
% \degree : commande fréquemment utilisée par les modèles pour un angle ou une
% température en dehors de \SI{}{} (non fournie par les paquets chargés).
\providecommand{\degree}{^{\circ}}
% \qed (fin de démonstration), utilisé par les modèles sans amsthm
\providecommand{\qed}{\hfill$\square$}
% \barre : double barre des valeurs interdites dans un tableau de variations (tkz-tab)
\providecommand{\barre}{\ensuremath{\|}}
% environnement proof (amsthm non chargé) utilisé par les modèles
\ifdefined\proof\else\newenvironment{proof}[1][Démonstration]{\par\noindent\textit{#1.}\ }{\qed\par}\fi
\usepackage{lastpage}
\usepackage{hyperref}
\hypersetup{hidelinks}

% ============================================================
% Palette « Pop » OnBuch+ — mêmes hex que la classe P (plus_ui.dart)
% ============================================================
\definecolor{poporange}{HTML}{F59321}
\definecolor{popdark}{HTML}{E07A0C}
\definecolor{popcream}{HTML}{FFF6E8}
\definecolor{popink}{HTML}{1C1714}
\definecolor{poppurple}{HTML}{7A5AE0}
\definecolor{popgreen}{HTML}{2E9E6A}
\definecolor{poppink}{HTML}{E0457B}
\definecolor{popblue}{HTML}{2F7DE1}
\definecolor{popgold}{HTML}{C98A12}
\definecolor{popmuted}{HTML}{8A7F76}
\definecolor{popline}{HTML}{EADFD2}
\definecolor{popgray}{HTML}{8A7F76}
\colorlet{popgrey}{popgray}
% Alias tolérants (le modèle nomme parfois les couleurs différemment)
\colorlet{popwhite}{white}
\colorlet{popblack}{popink}
\colorlet{popred}{poppink}
\colorlet{popyellow}{popgold}
% Teintes claires (fonds de tableaux/encadrés) — le modèle nomme parfois
% les couleurs avec un suffixe L (ex. popblueL) sans les avoir définies.
\colorlet{poporangeL}{poporange!20}
\colorlet{popdarkL}{popdark!20}
\colorlet{poppurpleL}{poppurple!20}
\colorlet{popgreenL}{popgreen!20}
\colorlet{poppinkL}{poppink!20}
\colorlet{popblueL}{popblue!20}
\colorlet{popgoldL}{popgold!20}
\colorlet{popredL}{popred!20}
\colorlet{popgrayL}{popgray!20}
% Teintes claires pour fonds de tableaux / zones
\colorlet{popblueL}{popblue!10!white}
\colorlet{poporangeL}{poporange!14!white}
\colorlet{popgreenL}{popgreen!12!white}
\colorlet{poppurpleL}{poppurple!12!white}
\colorlet{poppinkL}{poppink!12!white}
\colorlet{popgoldL}{popgold!14!white}

\pagecolor{popcream}

% ============================================================
% Typographie
% ============================================================
\defaultfontfeatures{Ligatures=TeX}
\setmainfont{PlusJakartaSans-Medium.ttf}[Path=../../../fonts/,BoldFont=PlusJakartaSans-Bold.ttf]
% Symboles, API et emojis absents de la police principale : repli sur DejaVu Sans (caractères actifs)
\newfontface\symfont{DejaVuSans.ttf}[Path=../../../fonts/]
\makeatletter
\newcommand\symactive[1]{\catcode#1=13 \begingroup\lccode`\~=#1 \lowercase{\endgroup\def~}{{\symfont\char#1}}}
\newcommand\symblank[1]{\catcode#1=13 \begingroup\lccode`\~=#1 \lowercase{\endgroup\def~}{}}
\newcommand\symhyph[1]{\catcode#1=13 \begingroup\lccode`\~=#1 \lowercase{\endgroup\def~}{\mbox{-}}}
\newcount\symc
\newcommand\symrange[2]{\symc=#1 \@whilenum\symc<#2 \do{\expandafter\symactive\expandafter{\the\symc}\advance\symc by 1 }}
\newcommand\symblankrange[2]{\symc=#1 \@whilenum\symc<#2 \do{\expandafter\symblank\expandafter{\the\symc}\advance\symc by 1 }}
\makeatother
\symrange{"0250}{"02FF}\symactive{"2713}\symactive{"2714}\symactive{"2717}\symactive{"2718}\symactive{"2610}\symactive{"2611}\symactive{"2612}
\symhyph{"2011}\symhyph{"2E1A}\symblankrange{"1F300}{"1FAFF}\symblank{"FE0F}
% Maths en sans-serif (Fira Math) avec chiffres et lettres droites de Plus
% Jakarta, pour que les formules se fondent dans le texte.
\usepackage[math-style=ISO,bold-style=ISO]{unicode-math}
\setmathfont{FiraMath-Regular.otf}[Path=../../../fonts/]
\setmathfont{PlusJakartaSans-Medium.ttf}[Path=../../../fonts/,range={up/num,up/{Latin,latin},\mathup}]
\usepackage{icomma} % virgule décimale sans espace en mode math
% Caractères mathématiques Unicode absents des polices de texte (Plus Jakarta,
% Archivo Black) : rendus par leur équivalent mathématique, en texte comme en
% formule — sinon ils s'affichent en carré vide (ex. « Arithmétique dans ℤ »).
\usepackage{newunicodechar}
\newunicodechar{ℝ}{\ensuremath{\mathbb{R}}}
\newunicodechar{ℤ}{\ensuremath{\mathbb{Z}}}
\newunicodechar{ℂ}{\ensuremath{\mathbb{C}}}
\newunicodechar{ℕ}{\ensuremath{\mathbb{N}}}
\newunicodechar{ℚ}{\ensuremath{\mathbb{Q}}}
\newunicodechar{𝔻}{\ensuremath{\mathbb{D}}}
\newunicodechar{α}{\ensuremath{\alpha}}
\newunicodechar{β}{\ensuremath{\beta}}
\newunicodechar{γ}{\ensuremath{\gamma}}
\newunicodechar{δ}{\ensuremath{\delta}}
\newunicodechar{ε}{\ensuremath{\varepsilon}}
\newunicodechar{θ}{\ensuremath{\theta}}
\newunicodechar{λ}{\ensuremath{\lambda}}
\newunicodechar{μ}{\ensuremath{\mu}}
\newunicodechar{σ}{\ensuremath{\sigma}}
\newunicodechar{τ}{\ensuremath{\tau}}
\newunicodechar{φ}{\ensuremath{\varphi}}
\newunicodechar{ψ}{\ensuremath{\psi}}
\newunicodechar{ω}{\ensuremath{\omega}}
\newunicodechar{Σ}{\ensuremath{\Sigma}}
% majuscules produites par \MakeUppercase dans les titres (α-aminés -> Α)
\newunicodechar{Α}{\ensuremath{\alpha}}
\newunicodechar{Β}{\ensuremath{\beta}}
\newunicodechar{Γ}{\ensuremath{\Gamma}}
\newunicodechar{Δ}{\ensuremath{\Delta}}
\newunicodechar{Ε}{\ensuremath{\varepsilon}}
\newunicodechar{Θ}{\ensuremath{\Theta}}
\newunicodechar{Λ}{\ensuremath{\Lambda}}
\newunicodechar{Μ}{\ensuremath{\mu}}
\newunicodechar{Τ}{\ensuremath{\tau}}
\newunicodechar{Φ}{\ensuremath{\Phi}}
\newunicodechar{Ψ}{\ensuremath{\Psi}}
\newunicodechar{Ω}{\ensuremath{\Omega}}
\newunicodechar{↦}{\ensuremath{\mapsto}}
\newunicodechar{∈}{\ensuremath{\in}}
\newunicodechar{∉}{\ensuremath{\notin}}
\newunicodechar{∧}{\ensuremath{\wedge}}
\newunicodechar{⇔}{\ensuremath{\Leftrightarrow}}
\newunicodechar{⟺}{\ensuremath{\Longleftrightarrow}}
\newunicodechar{⇒}{\ensuremath{\Rightarrow}}
\newunicodechar{⟹}{\ensuremath{\Longrightarrow}}
\newunicodechar{∘}{\ensuremath{\circ}}
\newunicodechar{∪}{\ensuremath{\cup}}
\newunicodechar{∩}{\ensuremath{\cap}}
\newunicodechar{≡}{\ensuremath{\equiv}}
\newunicodechar{⊂}{\ensuremath{\subset}}
\newunicodechar{⊥}{\ensuremath{\perp}}
\newunicodechar{∃}{\ensuremath{\exists}}
\newunicodechar{∀}{\ensuremath{\forall}}
\newunicodechar{∛}{\ensuremath{\sqrt[3]{\,}}}
\newunicodechar{∜}{\ensuremath{\sqrt[4]{\,}}}
\newunicodechar{⃗}{\ensuremath{{}^{\to}}}
\newunicodechar{ⁿ}{\ifmmode^{n}\else\textsuperscript{\ensuremath{n}}\fi}
\newunicodechar{ˣ}{\ifmmode^{x}\else\textsuperscript{\ensuremath{x}}\fi}
\newunicodechar{ᵇ}{\ifmmode^{b}\else\textsuperscript{\ensuremath{b}}\fi}
\newunicodechar{ʳ}{\ifmmode^{r}\else\textsuperscript{\ensuremath{r}}\fi}
\newunicodechar{ᵘ}{\ifmmode^{u}\else\textsuperscript{\ensuremath{u}}\fi}
\newunicodechar{ʸ}{\ifmmode^{y}\else\textsuperscript{\ensuremath{y}}\fi}
\newunicodechar{ᵃ}{\ifmmode^{a}\else\textsuperscript{\ensuremath{a}}\fi}
\newunicodechar{ᵉ}{\ifmmode^{e}\else\textsuperscript{\ensuremath{e}}\fi}
\newunicodechar{ᵏ}{\ifmmode^{k}\else\textsuperscript{\ensuremath{k}}\fi}
\newunicodechar{ᵖ}{\ifmmode^{p}\else\textsuperscript{\ensuremath{p}}\fi}
\newunicodechar{ᴺ}{\ifmmode^{N}\else\textsuperscript{\ensuremath{N}}\fi}
\newunicodechar{ᶜ}{\ifmmode^{c}\else\textsuperscript{\ensuremath{c}}\fi}
\newunicodechar{ʰ}{\ifmmode^{h}\else\textsuperscript{\ensuremath{h}}\fi}
\newunicodechar{ᵠ}{\ifmmode^{q}\else\textsuperscript{\ensuremath{q}}\fi}
\newunicodechar{ₙ}{\ifmmode_{n}\else\textsubscript{\ensuremath{n}}\fi}
\newunicodechar{ₐ}{\ifmmode_{a}\else\textsubscript{\ensuremath{a}}\fi}
\newunicodechar{ᵢ}{\ifmmode_{i}\else\textsubscript{\ensuremath{i}}\fi}
\newunicodechar{ᵣ}{\ifmmode_{r}\else\textsubscript{\ensuremath{r}}\fi}
\newunicodechar{ₖ}{\ifmmode_{k}\else\textsubscript{\ensuremath{k}}\fi}
\newunicodechar{ᵦ}{\ifmmode_{\beta}\else\textsubscript{\ensuremath{\beta}}\fi}
\newunicodechar{ₓ}{\ifmmode_{x}\else\textsubscript{\ensuremath{x}}\fi}
\newfontfamily\popdisplay{ArchivoBlack-Regular.ttf}[Path=../../../fonts/]
\newfontfamily\poplogo{SpaceGrotesk-Bold.ttf}[Path=../../../fonts/]
\newfontfamily\popsemi{PlusJakartaSans-SemiBold.ttf}[Path=../../../fonts/]
\newfontfamily\popxbold{PlusJakartaSans-ExtraBold.ttf}[Path=../../../fonts/]
\newfontfamily\popxboldls{PlusJakartaSans-ExtraBold.ttf}[Path=../../../fonts/,LetterSpace=3.0]

\setlist[enumerate]{leftmargin=1.7em,itemsep=5pt,topsep=4pt}
\setlist[itemize]{leftmargin=1.7em,itemsep=4pt,topsep=4pt,label={\color{poporange}\textbullet}}
\setlength{\parindent}{0pt}
\setlength{\parskip}{5pt}
\renewcommand{\arraystretch}{1.25}

% pgfplots : style Pop par défaut
\pgfplotsset{
  every axis/.append style={axis line style={popink,line width=1pt},tick style={popink},
    grid style={popline,line width=0.5pt},label style={font=\small},tick label style={font=\footnotesize}},
  popcurve/.style={poporange,line width=1.8pt,no markers,smooth},
}
% Même style disponible dans un \draw TikZ simple (hors pgfplots)
\tikzset{popcurve/.style={poporange,line width=1.8pt}}
\tikzset{popgrid/.style={popline,very thin}}
\colorlet{popcurve}{poporange} % le nom du style est parfois utilisé comme couleur

% ============================================================
% Pill / squiggle
% ============================================================
\newcommand{\poppill}[3]{%
  \tikz[baseline=-0.55ex]\node[draw=popink,line width=1.2pt,rounded corners=2.6pt,%
    fill=#2,inner xsep=6.5pt,inner ysep=3pt]{%
    \popxboldls\fontsize{7}{7}\selectfont\color{#3}\MakeUppercase{#1}};%
}
\newcommand{\popsquiggle}[1]{%
  \tikz[baseline=-0.4ex]{
    \draw[#1,line width=2.6pt,line cap=round]
      plot [smooth] coordinates {(0,0.09) (0.34,0.01) (0.68,0.21) (1.02,0.01) (1.36,0.10)};
  }%
}
% Mot-clé surligné (marqueur orange sous le texte)
\newcommand{\cle}[1]{{\popxbold\color{popdark}#1}}

% ============================================================
% En-tête / pied
% ============================================================
\pagestyle{fancy}
\fancyhf{}
\renewcommand{\headrulewidth}{0pt}
\renewcommand{\footrulewidth}{0pt}
\lhead{\raisebox{-0.15ex}{\poplogo\fontsize{13}{13}\selectfont\color{popink}OnBuch\color{poporange}+}}
\rhead{\poppill{\DOCMATIERE\ \textbullet\ \DOCNIVEAU\ \textbullet\ Leçon \DOCLECON}{white}{popink}}
\lfoot{\popxbold\fontsize{7.6}{7.6}\selectfont\color{popmuted}\MakeUppercase{Cours signé OnBuch+}\ \textbullet\ {\popsemi\DOCTITRE}}
\rfoot{\tikz[baseline=-0.6ex]\node[rounded corners=8.5pt,fill=popink,minimum height=17pt,minimum width=17pt,inner xsep=5pt,inner sep=0]{\popxbold\fontsize{8}{8}\selectfont\color{popcream}\thepage\,/\,\pageref*{LastPage}};}

% ============================================================
% Titres
% ============================================================
\newcommand{\chaptitre}[2]{%
  \begin{tcolorbox}[enhanced,colback=poporange,colframe=popink,boxrule=2.6pt,arc=11pt,
      drop shadow=popink,left=15pt,right=15pt,top=12pt,bottom=11pt,before skip=3pt,after skip=18pt]
    \poppill{#2}{popink}{popcream}\par\vspace{9pt}
    {\raggedright\hyphenpenalty=10000\exhyphenpenalty=10000\popdisplay\fontsize{22}{24}\selectfont\color{popcream}\MakeUppercase{#1}\par}
  \end{tcolorbox}
}
% Section numérotée : pastille orange avec le numéro + titre Archivo
\newcounter{popsec}
\newcommand{\coursec}[1]{%
  \stepcounter{popsec}\setcounter{popsub}{0}%
  \par\vspace{16pt}\noindent
  \begin{minipage}{\linewidth}
  \tikz[baseline=-0.7ex]\node[draw=popink,line width=1.6pt,fill=poporange,rounded corners=4pt,minimum size=22pt,inner sep=0]{\popdisplay\fontsize{12}{12}\selectfont\color{popcream}\thepopsec};%
  \hspace{8pt}{\popdisplay\fontsize{15}{17}\selectfont\color{popink}\MakeUppercase{#1}}\par
  \vspace{4pt}\noindent{\color{poporange}\rule{36pt}{3.6pt}}
  \end{minipage}\par\nopagebreak\vspace{8pt}
}
\newcounter{popsub}
\newcommand{\courssub}[1]{%
  \stepcounter{popsub}%
  \par\vspace{9pt}\noindent{\popxbold\fontsize{11.5}{13}\selectfont\color{popink}{\color{poporange}\thepopsec.\thepopsub}\hspace{6pt}#1}\par\nopagebreak\vspace{4pt}
}

% ============================================================
% Boîtes pédagogiques — même recette : fond blanc, bordure épaisse colorée,
% ombre dure encre, badge plein. Titre optionnel [#1].
% ============================================================
\tcbset{popbase/.style={breakable,enhanced,colback=white,boxrule=2.3pt,arc=9pt,
  shadow={3pt}{-3pt}{0pt}{popink},left=14pt,right=14pt,top=11pt,bottom=11pt,before skip=11pt,after skip=15pt,
  fontupper=\popsemi\fontsize{10.6}{14.5}\selectfont\color{popink},
  fontlower=\popsemi\fontsize{10.6}{14.5}\selectfont\color{popink}}}
% Coche dessinée (le glyphe ✓ n'existe pas dans les polices du document)
\newcommand{\popcheck}[1]{\tikz[baseline=-0.5ex]\draw[#1,line width=1.6pt,line cap=round,line join=round](-0.13,0.0)--(-0.04,-0.09)--(0.13,0.1);}
\providecommand{\coche}{\popcheck{popgreen}}
\providecommand{\resultat}[1]{\par\smallskip\noindent\fcolorbox{popgreen}{popgreenL}{\parbox{\dimexpr\linewidth-2\fboxsep-2\fboxrule\relax}{\textbf{Résultat.} #1}}\par\smallskip}
\newcommand{\popboxhead}[3]{\poppill{#1}{#2}{white}\ifx\relax#3\relax\else\hspace{6pt}{\popxbold\fontsize{10.6}{12}\selectfont\color{popink}#3}\fi\par\vspace{7pt}}

\newtcolorbox{aretenir}[1][]{popbase,colframe=popgreen,before upper={\popboxhead{À retenir}{popgreen}{#1}}}
% Texte en anglais (document de lecture, dialogue, extrait) : cadre
% d'étude. Un titre optionnel (source, auteur) s'affiche dans l'en-tête.
\newtcolorbox{textebox}[1][]{popbase,colframe=popblue,before upper={\popboxhead{Text}{popblue}{#1}\begin{otherlanguage}{english}},after upper={\end{otherlanguage}}}
% Marques ✓ / ✗ (forme correcte / fautive) souvent utilisées par les modèles
\providecommand{\cross}{\textcolor{poppink}{\ensuremath{\times}}}
\providecommand{\xmark}{\cross}\providecommand{\crossmark}{\cross}\providecommand{\cmark}{\ensuremath{\checkmark}}
% Ligne de réponse / justification à compléter (inventée par les modèles)
\providecommand{\justification}[1]{\par\noindent\textit{Justification~:} \dotfill\par}
% \textipa (package tipa non chargé) : la police ne couvre pas l'API -> simple passage
\providecommand{\textipa}[1]{#1}
% Soulignement (ulem non chargé) : \uline, \ul
\providecommand{\uline}[1]{\underline{#1}}\providecommand{\ul}[1]{\underline{#1}}
% \glossaire{\item ...} inventée par les modèles : liste de glossaire
\providecommand{\glossaire}[1]{\begin{itemize}#1\end{itemize}}
% \alert (beamer) inventée par les modèles : texte mis en évidence
\providecommand{\alert}[1]{\textbf{\textcolor{poppink}{#1}}}
% variantes ASCII d'ordinaux écrites en macro
\providecommand{\iere}{\textsuperscript{re}}\providecommand{\ieme}{\textsuperscript{e}}\providecommand{\ere}{\textsuperscript{re}}\providecommand{\eme}{\textsuperscript{e}}\providecommand{\er}{\textsuperscript{er}}\providecommand{\ie}{\textsuperscript{e}}
% Ordinaux français écrits en macro par les modèles : 1\ière, 3\ième
\newcommand{\ière}{\textsuperscript{re}}\newcommand{\ième}{\textsuperscript{e}}
% \exemple (inventée par les modèles) : étiquette « Exemple : »
\providecommand{\exemple}{\textbf{Exemple~:}\ }
% \square utilisable aussi hors mode math (cases à cocher)
\AtBeginDocument{\let\squareMath\square\renewcommand{\square}{\ensuremath{\squareMath}}}
% Mot ou phrase en anglais à l'intérieur d'un paragraphe français
\newcommand{\eng}[1]{\ifmmode\text{\textit{#1}}\else\begingroup\selectlanguage{english}\itshape #1\endgroup\fi}
\let\esp\eng % alias toléré
% flèches d'intonation inventées par les modèles
\providecommand{\textsearrow}{}\renewcommand{\textsearrow}{\ensuremath{\searrow}}\providecommand{\textnearrow}{}\renewcommand{\textnearrow}{\ensuremath{\nearrow}}
% \fr{..} : mot français (inventé par les modèles)
\providecommand{\fr}[1]{\textit{#1}}
\newtcolorbox{exemplebox}[1][]{popbase,colframe=poporange,before upper={\popboxhead{Exemple}{poporange}{#1}}}
\newtcolorbox{definition}[1][]{popbase,colframe=popblue,before upper={\popboxhead{Définition}{popblue}{#1}}}
\newtcolorbox{propriete}[1][]{popbase,colframe=poppurple,before upper={\popboxhead{Propriété}{poppurple}{#1}}}
\newtcolorbox{methode}[1][]{popbase,colframe=popgold,before upper={\popboxhead{Méthode}{popgold}{#1}}}
\newtcolorbox{attention}[1][]{popbase,colframe=poppink,before upper={\popboxhead{Attention}{poppink}{#1}}}
\newtcolorbox{experience}[1][]{popbase,colframe=popblue,colback=popblueL,before upper={\popboxhead{Expérience}{popblue}{#1}}}
% « Le savais-tu ? » : carte violette pleine (contexte, culture, Cameroun)
\newtcolorbox{savaistu}[1][]{popbase,colback=poppurple,colframe=popink,
  fontupper=\popsemi\fontsize{10.4}{14.2}\selectfont\color{white},
  before upper={\poppill{Le savais-tu\,?}{white}{poppurple}\ifx\relax#1\relax\else\hspace{6pt}{\popxbold\color{white}#1}\fi\par\vspace{7pt}}}
% Exercice résolu : énoncé en haut, \tcblower puis la solution sur fond crème
\newcounter{popexo}\newcounter{popexr}
\newtcolorbox{exoresolu}[1][]{popbase,colframe=poporange,colbacklower=popcream,
  segmentation style={popink,line width=1.2pt,dashed},
  before upper={\stepcounter{popexr}\popboxhead{Exercice résolu \thepopexr}{poporange}{#1}},
  before lower={\poppill{Solution}{popink}{popcream}\par\vspace{6pt}}}
% Exercice (non résolu en place) — la correction est dans la partie Corrigés
\newtcolorbox{exercice}[1][]{popbase,colframe=popink,
  before upper={\stepcounter{popexo}\popboxhead{Exercice \thepopexo}{popink}{#1}}}
% Corrigé
\newtcolorbox{corrige}[1][]{popbase,colframe=popgreen,colback=popgreenL,
  before upper={\popboxhead{Corrigé}{popgreen}{#1}}}
% Objectifs / prérequis / bilan
\newtcolorbox{objectifs}{popbase,colframe=popink,colback=poporangeL,
  before upper={\popboxhead{Objectifs de la leçon}{popink}{}}}
\newtcolorbox{prerequis}{popbase,colframe=popink,colback=popblueL,
  before upper={\popboxhead{Prérequis}{popblue}{}}}
\newtcolorbox{bilan}{popbase,colframe=popink,colback=white,boxrule=2.8pt,
  before upper={\popboxhead{Fiche bilan}{poporange}{L'essentiel à connaître pour l'examen}}}
% Figure encadrée avec légende
\newtcolorbox{popfigure}[1][]{enhanced,breakable=false,colback=white,colframe=popline,boxrule=1.4pt,arc=8pt,
  left=10pt,right=10pt,top=10pt,bottom=8pt,before skip=10pt,after skip=12pt,halign=center,
  lower separated=false,halign lower=center,
  fontlower=\popsemi\fontsize{9}{11.5}\selectfont\color{popmuted},
  #1}
\newcommand{\legende}[1]{\par\vspace{6pt}{\popsemi\fontsize{9}{11.5}\selectfont\color{popmuted}#1}}

% ============================================================
% Signature OnBuch+
% ============================================================
\newcommand{\poplogoinline}[1]{{\poplogo\fontsize{#1}{#1}\selectfont\color{popink}OnBuch\color{poporange}+}}
\newcommand{\signatureonbuch}{%
  \begin{tcolorbox}[enhanced,colback=popink,colframe=popink,boxrule=2.2pt,arc=10pt,
      drop shadow=poporange,left=14pt,right=14pt,top=10pt,bottom=10pt,before skip=0pt,after skip=0pt]
    \tikz[baseline=-0.7ex]\node[circle,fill=popgreen,draw=popcream,line width=1.3pt,minimum size=22pt,inner sep=0]{\popcheck{white}};%
    \hspace{9pt}\begin{minipage}[c]{0.62\linewidth}
      {\popxbold\fontsize{12}{13}\selectfont\color{popcream}Signé\ \textbullet\ {\poplogo OnBuch\color{poporange}+}}\par\vspace{2pt}
      {\raggedright\popsemi\fontsize{8}{10.5}\selectfont\color{popline}\MakeUppercase{Équipe pédagogique OnBuch+}\\\MakeUppercase{Conforme au programme officiel MINESEC}\par}
    \end{minipage}\hfill
    {\poplogo\fontsize{22}{22}\selectfont\color{popcream}OnBuch\color{poporange}+}
  \end{tcolorbox}%
}

% ============================================================
% Page de garde
% #1 titre · #2 matière · #3 niveau · #4 module · #5 n° leçon · #6 durée ·
% #7 description / objectifs (texte ou itemize)
% ============================================================
\newcommand{\pagegarde}[7]{%
  \thispagestyle{empty}
  \newgeometry{a4paper,margin=1.7cm,top=1.6cm,bottom=1.4cm}%
  \begin{tikzpicture}[remember picture,overlay]
    \fill[poporange!18] ([xshift=-3.2cm,yshift=-3.6cm]current page.north east) circle (4.6cm);
    \fill[poppurple!14] ([xshift=2.4cm,yshift=7.6cm]current page.south west) circle (3.8cm);
  \end{tikzpicture}%
  {\poplogo\fontsize{30}{30}\selectfont\color{popink}OnBuch\color{poporange}+}\hfill
  \raisebox{0.6ex}{\poppill{Cours signé \textbullet\ Premium}{popink}{popcream}}\par
  \vspace{1pt}\popsquiggle{poppurple}\par
  \vspace{8pt}
  \begin{tcolorbox}[enhanced,colback=poporange,colframe=popink,boxrule=2.8pt,arc=13pt,
      drop shadow=popink,left=18pt,right=18pt,top=12pt,bottom=13pt,after skip=10pt]
    \poppill{#2 \textbullet\ #3}{popink}{popcream}\hspace{5pt}\poppill{#4}{white}{popink}\par\vspace{8pt}
    {\popdisplay\fontsize{14}{14}\selectfont\color{popink}LEÇON #5}\par\vspace{4pt}
    {\raggedright\hyphenpenalty=10000\exhyphenpenalty=10000\popdisplay\fontsize{25}{27}\selectfont\color{popcream}\MakeUppercase{#1}\par}
    \vspace{8pt}
    {\popxbold\fontsize{9.5}{9.5}\selectfont\color{popink}\MakeUppercase{Durée conseillée : #6}}
  \end{tcolorbox}
  \begin{tcolorbox}[enhanced,colback=white,colframe=popink,boxrule=2.2pt,arc=10pt,
      drop shadow=popink,left=16pt,right=16pt,top=10pt,bottom=10pt,after skip=8pt]
    \poppill{Dans cette leçon}{poppurple}{white}\par\vspace{5pt}
    \popsemi\fontsize{9.3}{12.6}\selectfont\color{popink}#7
  \end{tcolorbox}
  \noindent\begin{minipage}[t]{0.487\linewidth}
    \begin{tcolorbox}[enhanced,colback=popgreen,colframe=popink,boxrule=2.2pt,arc=10pt,
        drop shadow=popink,left=11pt,right=11pt,top=10pt,bottom=10pt,height=3.1cm,valign=center]
      \tcbox[colback=white,colframe=popink,boxrule=1.3pt,arc=5pt,boxsep=0pt,nobeforeafter,
        left=3pt,right=3pt,top=3pt,bottom=3pt,valign=center]{\qrcode[height=1.9cm]{https://play.google.com/store/apps/details?id=com.luvvix.onbuch&hl=fr}}%
      \hspace{8pt}\begin{minipage}[c]{0.5\linewidth}
        {\popdisplay\fontsize{11}{12.5}\selectfont\color{white}\MakeUppercase{Télécharge OnBuch}}\par\vspace{4pt}
        {\raggedright\popsemi\fontsize{8.4}{11}\selectfont\color{white}Gratuit \textbullet\ Google Play\\Tuteur IA, annales, concours, résultats\par}
      \end{minipage}
    \end{tcolorbox}
  \end{minipage}\hfill
  \begin{minipage}[t]{0.487\linewidth}
    \begin{tcolorbox}[enhanced,colback=poppurple,colframe=popink,boxrule=2.2pt,arc=10pt,
        drop shadow=popink,left=11pt,right=11pt,top=10pt,bottom=10pt,height=3.1cm,valign=center]
      \tikz[baseline=-0.7ex]\node[circle,fill=white,draw=popink,line width=1.3pt,minimum size=1.6cm,inner sep=0]{\popdisplay\fontsize{24}{24}\selectfont\color{poppurple}+};%
      \hspace{8pt}\begin{minipage}[c]{0.6\linewidth}
        {\popdisplay\fontsize{11}{12.5}\selectfont\color{white}\MakeUppercase{Passe à OnBuch+}}\par\vspace{4pt}
        {\raggedright\popsemi\fontsize{8.4}{11}\selectfont\color{white}Tous les cours signés, Tuteur IA illimité, fiches PDF — onglet OnBuch+ de l'app\par}
      \end{minipage}
    \end{tcolorbox}
  \end{minipage}
  \vspace{8pt}
  \popcontenu
  \vfill
  \signatureonbuch
  \restoregeometry
  \newpage
}

% Rangée « Ce cours contient » (page de garde)
\newcommand{\poptile}[3]{% #1 couleur #2 gros texte #3 légende
  \begin{tcolorbox}[enhanced,colback=#1,colframe=popink,boxrule=2pt,arc=9pt,shadow={2.5pt}{-2.5pt}{0pt}{popink},
      left=6pt,right=6pt,top=9pt,bottom=9pt,halign=center,valign=center,height=1.75cm,width=\linewidth,nobeforeafter]
    {\popdisplay\fontsize{12.5}{13}\selectfont\color{white}\MakeUppercase{#2}}\par\vspace{3pt}
    {\centering\popsemi\fontsize{7.8}{9.5}\selectfont\color{white}#3\par}
  \end{tcolorbox}}
\newcommand{\popcontenu}{%
  \par\noindent{\popxboldls\fontsize{7.5}{7.5}\selectfont\color{popmuted}CE COURS SIGNÉ CONTIENT}\par\vspace{7pt}
  \noindent\begin{minipage}[t]{0.235\linewidth}\poptile{popblue}{Cours}{complet, illustré, pas à pas}\end{minipage}\hfill
  \begin{minipage}[t]{0.235\linewidth}\poptile{poporange}{Méthodes}{et exercices résolus}\end{minipage}\hfill
  \begin{minipage}[t]{0.235\linewidth}\poptile{poppink}{Exercices}{type examen + corrigés}\end{minipage}\hfill
  \begin{minipage}[t]{0.235\linewidth}\poptile{popgreen}{Fiche bilan}{l'essentiel à retenir}\end{minipage}\par}

% Sommaire
\newtcolorbox{sommaire}{enhanced,colback=white,colframe=popink,boxrule=2.2pt,arc=10pt,
  drop shadow=popink,left=18pt,right=18pt,top=13pt,bottom=13pt,before skip=6pt,after skip=20pt,
  before upper={\poppill{Sommaire}{poppurple}{white}\par\vspace{9pt}\popsemi\fontsize{10.6}{14.5}\selectfont\color{popink}}}

% Signature de fin de cours — poussée tout en bas de la dernière page
\newcommand{\findecours}{%
  \par\vfill\nopagebreak
  \begin{center}\popsquiggle{poporange}\end{center}\vspace{4pt}
  \signatureonbuch
}
\AtBeginDocument{\def\llbracket{\mathopen{[\mkern-3.5mu[}}\def\rrbracket{\mathclose{]\mkern-3.5mu]}}} % intervalles d'entiers (glyphe ⟦ absent de la police)
% Glyphes absents de Fira Math : redéfinis à partir de glyphes disponibles
\AtBeginDocument{%
  \def\setminus{\mathbin{\backslash}}\let\smallsetminus\setminus
  \def\vdots{\mathord{\vbox{\baselineskip3.2pt\lineskiplimit0pt\kern4pt\hbox{.}\hbox{.}\hbox{.}}}}%
  \def\ddots{\mathinner{\mkern1mu\raise6.5pt\hbox{.}\mkern2mu\raise3.6pt\hbox{.}\mkern2mu\raise0.7pt\hbox{.}\mkern1mu}}%
  \def\triangle{\mathord{\scalebox{0.9}{$\symup{\Delta}$}}}%
  \def\nabla{\mathord{\raisebox{\depth}{\scalebox{1}[-1]{$\symup{\Delta}$}}}}%
  \def\checkmark{\popcheck{popdark}}%
  \def\star{\mathbin{\ast}}\def\bigstar{\mathord{\ast}}%
  \def\diamond{\mathbin{\scalebox{0.6}{$\Diamond$}}}%
  \def\bigcup{\mathop{\mathchoice{\raisebox{-0.3em}{\scalebox{1.7}{$\cup$}}}{\scalebox{1.25}{$\cup$}}{\cup}{\cup}}\displaylimits}%
  \def\bigcap{\mathop{\mathchoice{\raisebox{-0.3em}{\scalebox{1.7}{$\cap$}}}{\scalebox{1.25}{$\cap$}}{\cap}{\cap}}\displaylimits}%
  \def\lesseqqgtr{\mathrel{\lessgtr}}\def\bigoplus{\mathop{\oplus}\displaylimits}%
}
\providecommand{\ssi}{si et seulement si} % abréviation fréquente
\AtBeginDocument{\providecommand{\wideparen}{\overparen}} % arc de cercle

%% FIGFIX-BEGIN v3 — correctifs globaux des figures (docs/onbuch-generation/tools/figfix)
\usepackage{adjustbox}\usepackage{etoolbox}
% toute figure TikZ/pgfplots plus large que la ligne est réduite à la largeur disponible
\BeforeBeginEnvironment{tikzpicture}{\begin{adjustbox}{max width=\linewidth}}
\AfterEndEnvironment{tikzpicture}{\end{adjustbox}}
% légendes pgfplots lisibles par-dessus les courbes
\pgfplotsset{every axis/.append style={legend style={fill=white,fill opacity=0.92,text opacity=1,draw=black!15,inner sep=3pt}}}
% étiquettes d'arêtes (syntaxe edge["..."]) avec fond blanc
\tikzset{every edge quotes/.append style={fill=white,inner sep=1.2pt}}
% bulles de cartes mentales : texte à la ligne dans la bulle, bulle un peu plus grande
\tikzset{every concept/.append style={text width=2.0cm,align=center,minimum size=2.3cm,inner sep=2pt}}
%% FIGFIX-END
\begin{document}

\pagegarde{Lesson 18 — Vocabulary: Words and expressions related to banking services}{Anglais}{Tle A}{Chapter 2 : Economic Life and Occupations}{EN18}{2 h}{Cette leçon de vocabulaire présente les mots et expressions essentiels liés aux services bancaires (comptes, opérations, crédits, moyens de paiement). Elle permet aux élèves de Tle A de comprendre et d'employer ce lexique en contexte, à l'oral comme à l'écrit, en évitant les faux amis fréquents.
\vspace{4pt}
\begin{multicols}{2}\small
\begin{itemize}
\item À la fin de cette leçon, je sais identifier et définir le vocabulaire de base des services bancaires (account, deposit, withdrawal, loan...)
\item classer ce lexique par champs thématiques (personnes, opérations, moyens de paiement, crédit)
\item employer les collocations courantes du domaine bancaire (open an account, pay interest, take out a loan...)
\item reconnaître et éviter les faux amis et les calques du français (to save ≠ salver, a cheque ≠ un chèque au sens de 'check')
\item utiliser les mots de la même famille (bank, banker, banking, bankable ; lend, lender, lending)
\item prononcer correctement les mots clés du champ lexical (debt, mortgage, receipt...)
\end{itemize}\end{multicols}}

\begin{sommaire}
\begin{multicols}{2}
\begin{enumerate}
\item Introduction : the world of banking
\item Champ lexical 1 : people and places
\item Champ lexical 2 : accounts and operations
\item Champ lexical 3 : payment methods, credit and savings
\item Faux amis, familles de mots et expressions
\item Lexique en contexte : dialogues et production guidée
\item Activité d'intégration
\item Méthodes et exercices résolus
\item Exercices
\item Corrigés des exercices
\item Fiche bilan
\end{enumerate}
\end{multicols}
\end{sommaire}

\begin{objectifs}
\begin{itemize}
\item À la fin de cette leçon, je sais identifier et définir le vocabulaire de base des services bancaires (account, deposit, withdrawal, loan...)
\item classer ce lexique par champs thématiques (personnes, opérations, moyens de paiement, crédit)
\item employer les collocations courantes du domaine bancaire (open an account, pay interest, take out a loan...)
\item reconnaître et éviter les faux amis et les calques du français (to save ≠ salver, a cheque ≠ un chèque au sens de 'check')
\item utiliser les mots de la même famille (bank, banker, banking, bankable ; lend, lender, lending)
\item prononcer correctement les mots clés du champ lexical (debt, mortgage, receipt...)
\item réemployer le lexique dans des phrases, dialogues et courts textes en contexte camerounais et anglophone
\item comprendre un court texte ou dialogue bancaire au Bac et répondre aux questions de vocabulaire
\end{itemize}
\end{objectifs}

\begin{prerequis}
\begin{itemize}
\item Vocabulaire général de l'argent et des achats (money, price, pay, buy, sell) vu en classe de Première
\item Les noms dénombrables et indénombrables (money, cash, information)
\item Les verbes irréguliers courants (lend–lent–lent, pay–paid–paid, withdraw–withdrew–withdrawn)
\item Les structures de base du dialogue (questions avec wh- words, présent simple)
\item Notion de prépositions après verbes (pay for, depend on)
\end{itemize}
\end{prerequis}

\newpage

\coursec{Introduction : the world of banking}

\courssub{Situation d'accroche : Amina's problem}

Imaginez la scène : Amina, a student in Buea, has just received her first scholarship. Her uncle tells her: \eng{“Open a bank account, don't keep the money under your mattress!”} At the bank, she hears words like \eng{deposit}, \eng{withdrawal}, \eng{interest rate}, \eng{overdraft}... and she is lost. She smiles politely, signs the papers, but understands almost nothing of what the \eng{bank clerk} is explaining.

Cette situation est celle de millions de jeunes Camerounais. Aujourd'hui, avec \eng{mobile money} (MTN MoMo, Orange Money), les \eng{microfinance institutions} et les \eng{commercial banks} présentes partout — de Yaoundé à Bamenda —, comprendre le vocabulaire bancaire en anglais est une compétence de la vie réelle (\eng{a real-life skill}). C'est aussi un sujet fréquent dans les textes du Bac sur la vie économique (\eng{economic life}).

\textbf{Problématique de la leçon :} quels mots et quelles expressions faut-il maîtriser pour parler des services bancaires en anglais, les comprendre dans un texte du Bac, et les réemployer sans calquer le français ?

\courssub{Brainstorming : where do people keep their money?}

Avant tout vocabulaire, posons-nous une question simple : \eng{Where do people keep their money?} (Où les gens gardent-ils leur argent ?) En classe, les réponses fusent :

\begin{itemize}
\item \eng{in a commercial bank} — dans une banque commerciale (Afriland, BICEC, SGBC...) ;
\item \eng{in a microfinance institution} — dans une institution de microfinance (les « caisses » comme MC2 ou les coopératives) ;
\item \eng{on a mobile money account} — sur un compte mobile money (MTN MoMo, Orange Money) ;
\item \eng{in a tontine / njangui} — dans une tontine (mot camerounais admis en anglais dans le contexte local !) ;
\item \eng{under the mattress} — sous le matelas (expression idiomatique anglaise aussi : \eng{to keep money under the mattress}).
\end{itemize}

\begin{attention}[Tontine : un cas particulier]
Le mot \eng{tontine} existe en anglais (emprunté au français !), mais beaucoup d'Anglo-Saxons ne le connaissent pas. Dans une rédaction du Bac, on peut écrire \eng{a tontine (a rotating savings group)} pour expliquer. Évitez de traduire \eng{njangui} par autre chose : gardez le mot local et ajoutez une courte explication entre parenthèses.
\end{attention}

\courssub{The main types of financial institutions}

Observons d'abord ce petit tableau comparatif, puis nous dégagerons les définitions.

\begin{center}
\begin{tabularx}{\linewidth}{|X|X|X|}\hline
\rowcolor{popblueL} Institution & English definition & French equivalent \\ \hline
commercial bank & a large company that keeps, lends and transfers money for the public & banque commerciale \\ \hline
microfinance institution (MFI) & a small financial organisation that gives small loans to people who cannot use big banks & institution de microfinance \\ \hline
mobile money service & a telephone-based service to send, receive and keep money (MTN MoMo, Orange Money) & service d'argent mobile \\ \hline
central bank & the national bank that controls the other banks (BEAC in Central Africa) & banque centrale \\ \hline
tontine / njangui & an informal rotating savings group & tontine \\ \hline
\end{tabularx}
\end{center}

\begin{definition}[Banking services]
\eng{Banking services} are all the operations a bank offers to its customers: \cle{keeping money} (garder l'argent), \cle{lending money} (prêter de l'argent), \cle{transferring money} (transférer de l'argent) and \cle{providing means of payment} (fournir des moyens de paiement : cartes, chèques, virements).
\end{definition}

\begin{exemplebox}[Institutions in sentences]
\eng{Amina opened an account at a commercial bank in Buea.} « Amina a ouvert un compte dans une banque commerciale à Buéa. »

\eng{Market women in Douala often borrow from microfinance institutions.} « Les vendeuses du marché à Douala empruntent souvent auprès d'institutions de microfinance. »

\eng{My father sends school fees by mobile money.} « Mon père envoie les frais de scolarité par mobile money. »

\eng{Every month, our njangui collects 10,000 francs from each member.} « Chaque mois, notre tontine collecte 10 000 francs auprès de chaque membre. »
\end{exemplebox}

\courssub{Reading : Amina's First Bank Account}

\begin{methode}[How to read a vocabulary text — trois étapes]
\begin{enumerate}
\item \textbf{First reading for global meaning} : lisez le texte une première fois, rapidement, pour comprendre le sens général. Qui ? Où ? Que se passe-t-il ? Ne vous arrêtez pas sur les mots inconnus.
\item \textbf{Underline the theme words} : relisez et soulignez tous les mots qui appartiennent au champ lexical du thème (ici : la banque).
\item \textbf{Guess before you check} : devinez le sens de chaque mot souligné grâce au contexte \emph{avant} de consulter le glossaire. Le contexte donne presque toujours la réponse.
\end{enumerate}
\end{methode}

Appliquons cette méthode au texte support. Première lecture : de quoi parle le texte ?

\begin{textebox}[Amina's First Bank Account]
Last Monday, Amina went to a commercial bank in Buea to open a savings account. She had just received her scholarship, and her uncle had told her: “Don't keep your money under the mattress!”

At the entrance, a security guard directed her to the customer service desk. A friendly bank clerk welcomed her and asked for her ID card and an initial deposit of 5,000 francs CFA. Amina filled in a form with her name, address and telephone number, and signed at the bottom.

The clerk opened the account on the computer and gave Amina her account number and a brand-new bank card. “With this card,” he explained, “you can deposit and withdraw money at any branch of the bank, or at any ATM in town. You can also check your balance on your phone.”

He added that her savings would earn interest every year — about three per cent — and that the bank could give her a loan one day, if she wanted to start a small business after her studies. “But be careful,” he smiled, “never spend more than you have, or you will have an overdraft, and the bank will charge you fees.”

Amina thanked the clerk and walked out proudly. Her money was now safe, and it was even working for her. That evening, she sent a message to her uncle: “Account opened! No more money under the mattress!”
\end{textebox}

\textbf{Glossaire (mots difficiles) :}

\begin{center}
\begin{tabularx}{\linewidth}{|X|X|X|}\hline
\rowcolor{popblueL} Word & Nature & Traduction \\ \hline
savings account & nom & compte d'épargne \\ \hline
bank clerk & nom & employé(e) de banque, guichetier \\ \hline
ID card & nom & carte d'identité \\ \hline
initial deposit & nom & dépôt initial, premier versement \\ \hline
to fill in a form & verbe & remplir un formulaire \\ \hline
account number & nom & numéro de compte \\ \hline
bank card & nom & carte bancaire \\ \hline
to deposit & verbe & déposer (de l'argent) \\ \hline
to withdraw & verbe & retirer (de l'argent) \\ \hline
branch & nom & agence, succursale \\ \hline
ATM & nom & guichet automatique (DAB) \\ \hline
balance & nom & solde (du compte) \\ \hline
to earn interest & expression & rapporter des intérêts \\ \hline
loan & nom & prêt, emprunt \\ \hline
overdraft & nom & découvert (bancaire) \\ \hline
to charge fees & expression & prélever des frais \\ \hline
\end{tabularx}
\end{center}

\textbf{Comprehension questions :}

\begin{enumerate}
\item Why did Amina go to the bank? (\eng{To open a savings account.})
\item What documents did the clerk ask for? (\eng{Her ID card and an initial deposit of 5,000 francs CFA.})
\item What can Amina do with her bank card? (\eng{Deposit and withdraw money at any branch or ATM, and check her balance.})
\item How much interest will her savings earn? (\eng{About three per cent every year.})
\item What happens if she spends more than she has? (\eng{She will have an overdraft and the bank will charge her fees.})
\end{enumerate}

\courssub{Spotting the lexical field in the text}

Deuxième étape de la méthode : soulignons les mots du champ lexical. Dans le texte, on relève, dans l'ordre d'apparition : \eng{commercial bank}, \eng{savings account}, \eng{bank clerk}, \eng{ID card}, \eng{initial deposit}, \eng{account number}, \eng{bank card}, \eng{deposit}, \eng{withdraw}, \eng{branch}, \eng{ATM}, \eng{balance}, \eng{earn interest}, \eng{loan}, \eng{overdraft}, \eng{charge fees}.

Troisième étape : deviner par le contexte. Par exemple, \eng{withdraw} apparaît dans la phrase \eng{you can deposit and withdraw money} : le mot est opposé à \eng{deposit} (déposer), donc il signifie « retirer ». De même, \eng{overdraft} est expliqué par \eng{spend more than you have} : dépenser plus que ce qu'on possède, c'est être « à découvert ».

\begin{exoresolu}[Guess the meaning from the context]
Énoncé : devinez le sens des mots en gras grâce au contexte, sans dictionnaire.

1. \eng{The bank will \textbf{charge} you fees.} 2. \eng{You can check your \textbf{balance} on your phone.} 3. \eng{The bank could give her a \textbf{loan} to start a small business.}
\tcblower
Solution détaillée :

1. La phrase parle d'une punition en cas de découvert (\eng{never spend more than you have... the bank will charge you fees}) : la banque va \textbf{prélever/facturer} des frais. \eng{Charge} = prélever (et non « charger » au sens de mettre une charge !).

2. On peut vérifier son \eng{balance} sur son téléphone : c'est ce qui reste sur le compte, donc le \textbf{solde}.

3. \eng{To start a small business}, la banque peut donner quelque chose : de l'argent prêté, donc un \textbf{prêt}. Le contexte (après les études, créer une entreprise) confirme.
\end{exoresolu}

\courssub{First mind map : BANKING SERVICES}

Organisons maintenant ce premier vocabulaire en carte mentale. Au centre, le thème ; autour, quatre grandes familles que nous étudierons dans les sections suivantes : les personnes et lieux (\eng{People}), les opérations (\eng{Operations}), les moyens de paiement (\eng{Payment methods}), le crédit et l'épargne (\eng{Credit \& Savings}).

\begin{popfigure}
\begin{center}
\begin{tikzpicture}[mindmap, grow cyclic,
every node/.style={concept, align=center, font=\small},
root concept/.append style={concept color=popblue, font=\bfseries\small, minimum size=3.2cm},
level 1/.append style={level distance=4.2cm, sibling angle=90},
level 2/.append style={level distance=2.8cm, sibling angle=45, font=\scriptsize}]
\node[root concept, align=center]{BANKING\\SERVICES}
child[concept color=poporange]{ node[align=center]{People\\\& Places}
  child[concept color=poporangeL]{ node[align=center]{bank clerk\\customer\\manager} }
  child[concept color=poporangeL]{ node[align=center]{branch\\ATM\\counter} }
}
child[concept color=popgreen]{ node {Operations}
  child[concept color=popgreenL]{ node[align=center]{deposit\\withdraw\\transfer} }
  child[concept color=popgreenL]{ node[align=center]{open / close\\an account} }
}
child[concept color=popgold]{ node[align=center]{Payment\\methods}
  child[concept color=popgoldL]{ node[align=center]{cash\\cheque\\bank card} }
  child[concept color=popgoldL]{ node[align=center]{mobile\\money} }
}
child[concept color=poppurple]{ node[align=center]{Credit \&\\Savings}
  child[concept color=poppurpleL]{ node[align=center]{loan\\interest\\overdraft} }
  child[concept color=poppurpleL]{ node[align=center]{savings\\account\\balance} }
};
\end{tikzpicture}
\end{center}
\legende{Carte mentale du champ lexical : les quatre branches de BANKING SERVICES, avec les premiers mots rencontrés dans le texte.}
\end{popfigure}

\begin{aretenir}[Ce qu'il faut retenir de cette introduction]
\begin{itemize}
\item \eng{Banking services} = garder, prêter, transférer l'argent et fournir des moyens de paiement.
\item Les institutions : \eng{commercial bank}, \eng{microfinance institution}, \eng{mobile money service}, \eng{central bank}, \eng{tontine}.
\item Méthode de lecture : sens global $\rightarrow$ mots du thème soulignés $\rightarrow$ sens deviné par le contexte.
\item Premiers mots-clés : \eng{account, deposit, withdraw, balance, interest, loan, overdraft, fees}.
\end{itemize}
\end{aretenir}

\begin{attention}[Premiers pièges pour les francophones]
\begin{itemize}
\item \eng{Bank} se prononce « bank » (a comme dans « sac »), pas « benk ».
\item Ne dites pas \eng{*to retire money} pour « retirer de l'argent » : \eng{retire} signifie « prendre sa retraite » ! Dites \eng{to withdraw money}.
\item \eng{Balance} ne signifie pas « balance » (l'instrument) dans ce contexte, mais « solde ».
\item Attention à l'orthographe : \eng{account} (deux c), \eng{clerk} (prononcé « clârk » en anglais britannique), \eng{interest} (un seul r, deux e).
\end{itemize}
\end{attention}

\begin{savaistu}[Two banking worlds]
In Cameroon, millions of people use mobile money every day — to pay for bendskins, market goods or school fees. But the Bac often uses texts about traditional banking in Britain or the USA, with words like \eng{current account} (compte courant), \eng{chequebook} (chéquier) or \eng{high street bank} (grande banque de réseau). You must know both worlds! In Britain, people say \eng{current account}; in the USA, they say \eng{checking account} — same thing, different word.
\end{savaistu}

\begin{exercice}[To go further]
1. Recopiez la carte mentale et ajoutez deux mots du texte dans chaque branche. 2. Écrivez trois phrases vraies sur votre vie (ou celle de votre famille) avec \eng{deposit}, \eng{mobile money} et \eng{loan}. 3. Question orale : \eng{Where do you keep your money, and why?} — préparez une réponse de trois phrases.
\end{exercice}

\coursec{Champ lexical 1 : people and places}

Dans la section précédente, nous avons découvert le monde de la banque (\eng{the world of banking}) : pourquoi les gens confient leur argent à une banque, et quels services elle propose. Pour aller plus loin, il faut maintenant pouvoir nommer \cle{les personnes} qui travaillent dans une banque ou qui la fréquentent, et \cle{les lieux} qui la composent. C'est l'objectif de cette section : construire un vocabulaire précis, bien prononcé et sans calque du français.

\courssub{Observation : a first visit to the bank}

Commençons par observer le vocabulaire en situation. Lisez ce court texte original et repérez les mots qui désignent des personnes et des lieux.

\begin{textebox}[A morning at the bank — original text]
Last Monday, Mrs Amina Bello went to her local bank in Buea. It is not the head office, just a small branch near the market. Outside the building, a customer was using the cash machine to withdraw some money. Inside, Amina joined the queue in front of the counter. The bank clerk behind the counter smiled and asked how he could help her. She wanted to pay some money into her savings account, because she is a careful saver. While she was waiting, she saw the bank manager come out of his office at the back to speak to a borrower who wanted a loan for his shop. The manager explained that the bank, as a lender, must check every request carefully. The branch accountant was busy checking the figures on a computer. Before leaving, Amina asked the clerk where the safe was kept. He laughed and said, “In the vault, of course — customers are not allowed to see it!”
\end{textebox}

\textbf{Glossaire :}
\begin{itemize}
\item \eng{branch} (n.) : agence, succursale
\item \eng{cash machine} (n.) : distributeur automatique de billets
\item \eng{to withdraw} (v.) : retirer (de l'argent)
\item \eng{counter} (n.) : guichet, comptoir
\item \eng{saver} (n.) : épargnant(e)
\item \eng{borrower} (n.) : emprunteur(-euse)
\item \eng{lender} (n.) : prêteur(-euse)
\item \eng{loan} (n.) : prêt
\item \eng{safe / vault} (n.) : coffre-fort / chambre forte
\end{itemize}

\textbf{Comprehension questions :}
\begin{enumerate}
\item Where is Amina's branch situated?
\item Who was using the cash machine outside?
\item What did Amina want to do at the counter?
\item Why did the bank manager speak to the borrower?
\item Where is the money kept safely?
\end{enumerate}

\courssub{The people : les personnes de la banque}

\begin{definition}[People in the banking world]
\begin{itemize}
\item \cle{customer} : a person who uses the services of a bank (client d'une banque).
\item \cle{bank clerk} (GB) / \cle{teller} (US) : the employee who serves customers at the counter (employé de banque au guichet).
\item \cle{bank manager} : the person who directs a branch (directeur d'agence).
\item \cle{accountant} : a person who keeps and checks financial records (comptable).
\item \cle{borrower} : a person who borrows money (emprunteur).
\item \cle{lender} : a person or institution that lends money (prêteur).
\item \cle{saver} : a person who saves money regularly (épargnant).
\end{itemize}
\end{definition}

Observez la construction des mots \eng{borrower}, \eng{lender}, \eng{saver} : on ajoute le suffixe \cle{-er} au verbe (\eng{to borrow}, \eng{to lend}, \eng{to save}) pour obtenir la personne qui fait l'action. C'est un procédé très productif en anglais : \eng{to teach → teacher}, \eng{to work → worker}.

\begin{exemplebox}[Exemples traduits]
\eng{The bank clerk counted the money carefully.} = « L'employé de banque a compté l'argent attentivement. »

\eng{The manager refused the loan because the borrower had no regular income.} = « Le directeur a refusé le prêt parce que l'emprunteur n'avait pas de revenu régulier. »

\eng{My grandmother is a great saver; she puts money aside every month.} = « Ma grand-mère est une grande épargnante ; elle met de l'argent de côté chaque mois. »

\eng{The accountant discovered a mistake in the branch accounts.} = « Le comptable a découvert une erreur dans les comptes de l'agence. »
\end{exemplebox}

\begin{attention}[Customer ou client ?]
En français, « client » convient partout. En anglais, on distingue : \eng{customer} pour les services courants (banque, magasin, supermarché) et \eng{client} pour les professions libérales (avocat, architecte, consultant). On dit donc \eng{a bank customer}, jamais \eng{a bank client} dans le langage courant. De même, \eng{guest} désigne le client d'un hôtel.
\end{attention}

\begin{attention}[Clerk ou teller ?]
\eng{Bank clerk} est le terme britannique ; \eng{teller} est le terme américain. Les deux désignent l'employé au guichet. Attention à la prononciation : \eng{clerk} se prononce « klâk » en anglais britannique (comme \eng{Clark} Kent !) et « kleurk » en anglais américain. Ne prononcez jamais le « e » comme en français !
\end{attention}

\courssub{The places : les lieux de la banque}

\begin{definition}[Places in the banking world]
\begin{itemize}
\item \cle{bank} : the financial institution (la banque).
\item \cle{branch} : a local office of a bank (une agence, une succursale).
\item \cle{head office} : the main administrative centre (le siège social).
\item \cle{counter} : the desk where customers are served (le guichet, le comptoir).
\item \cle{cash machine} (GB) / \cle{ATM} (US) : the machine that gives cash (le distributeur).
\item \cle{safe} / \cle{vault} : the strong room where money and valuables are kept (le coffre-fort / la chambre forte).
\end{itemize}
\end{definition}

\begin{definition}[ATM]
\cle{ATM} signifie \eng{Automated Teller Machine} : littéralement « machine-guichet automatique ». En Grande-Bretagne, on dit plutôt \eng{cash machine} ou \eng{cashpoint} ; aux États-Unis, \eng{ATM}. Traduction française : distributeur automatique de billets (DAB).
\end{definition}

\begin{exemplebox}[Exemples traduits]
\eng{You can withdraw cash from the ATM outside.} = « Vous pouvez retirer de l'argent au distributeur à l'extérieur. »

\eng{The head office of the bank is in Douala, but there are branches in Bamenda and Yaoundé.} = « Le siège de la banque est à Douala, mais il y a des agences à Bamenda et à Yaoundé. »

\eng{Please wait at the counter; the clerk will serve you in a minute.} = « Veuillez patienter au guichet ; l'employé vous servira dans une minute. »

\eng{All the documents are kept in the vault overnight.} = « Tous les documents sont conservés dans la chambre forte pendant la nuit. »
\end{exemplebox}

\begin{propriete}[Prononciation des mots clés]
\begin{itemize}
\item \eng{customer} : « kas-tè-meur » — l'accent tombe sur la première syllabe : \textbf{CUStomer}.
\item \eng{clerk} : « klâk » (GB) / « kleurk » (US).
\item \eng{manager} : « ma-na-djeur » — accent sur \textbf{MAnager} ; le « g » se prononce « dj ».
\item \eng{counter} : « kaoun-teur ».
\item \eng{vault} : « volt » — le « au » se prononce « o ».
\end{itemize}
\end{propriete}

\courssub{Tableau récapitulatif}

\begin{center}
\begin{tabularx}{\linewidth}{|X|X|X|}
\hline
\rowcolor{popblueL} English & Français & Remarque \\
\hline
customer & client (de banque) & pas \eng{client} pour une banque \\
\hline
bank clerk (GB) / teller (US) & employé de banque, guichetier & « klâk » en GB \\
\hline
bank manager & directeur d'agence & accent sur la 1\`{e}re syllabe \\
\hline
accountant & comptable & verbe : \eng{to account for} \\
\hline
borrower / lender & emprunteur / prêteur & suffixe -er \\
\hline
saver & épargnant & de \eng{to save} \\
\hline
branch & agence, succursale & prononcé « brantch » \\
\hline
head office & siège social & aussi \eng{headquarters} \\
\hline
counter & guichet, comptoir & \eng{at the counter} \\
\hline
cash machine (GB) / ATM (US) & distributeur de billets & \eng{Automated Teller Machine} \\
\hline
safe / vault & coffre-fort / chambre forte & \eng{vault} = « volt » \\
\hline
\end{tabularx}
\end{center}

\courssub{Visualisation : le plan d'une agence}

\begin{popfigure}
\begin{center}
\begin{tikzpicture}[font=\small]
\draw[thick, popdark, fill=popcream] (0,0) rectangle (8,5);
\node[popdark] at (4,5.35) {\textbf{A bank branch}};
\draw[thick, popblue, fill=popblueL] (0.3,0.3) rectangle (2.2,1.5);
\node[align=center, popdark] at (1.25,0.9) {entrance\\(entrée)};
\draw[thick, poporange, fill=poporangeL] (3,2.2) rectangle (5,3.2);
\node[align=center, popdark] at (4,2.7) {counter\\(guichet)};
\node[circle, draw=popgreen, fill=popgreenL, minimum size=1.1cm, align=center] (clerk) at (4,4.1) {clerk};
\node[circle, draw=poppurple, fill=poppinkL, minimum size=1.1cm, align=center] (cust) at (4,1.3) {customer};
\draw[thick, poppurple, fill=poppurpleL] (5.8,3.6) rectangle (7.7,4.7);
\node[align=center, popdark] at (6.75,4.15) {manager's\\office};
\draw[thick, popgold, fill=popgoldL] (5.8,0.3) rectangle (7.7,1.4);
\node[align=center, popdark] at (6.75,0.85) {safe / vault\\(at the back)};
\draw[thick, popgreen, fill=popgreenL] (-2.3,0.3) rectangle (-0.4,1.5);
\node[align=center, popdark] at (-1.35,0.9) {ATM\\(outside)};
\draw[-{Stealth}, thick, poppink] (1.25,1.5) -- (3.4,2.1);
\draw[-{Stealth}, thick, popblue] (4,2.1) -- (cust);
\draw[-{Stealth}, thick, popgreen] (cust) -- (4,2.0);
\end{tikzpicture}
\end{center}
\legende{Plan simplifié d'une agence bancaire : le client entre, se rend au guichet où l'attend l'employé ; le bureau du directeur et le coffre sont au fond ; le distributeur est à l'extérieur.}
\end{popfigure}

\courssub{Mini-dialogue de repérage : at the counter}

\begin{textebox}[At the counter — dialogue]
Clerk: Good morning, madam. How can I help you?

Customer: Good morning. I would like to pay this cheque into my account, please.

Clerk: Certainly. Can I see your bank card?

Customer: Here you are. By the way, is the cash machine outside working today?

Clerk: Yes, it is. You can withdraw money there at any time.

Customer: Thank you. And could I also speak to the bank manager about a loan?

Clerk: Of course. His office is at the back, next to the vault. Please take a seat; he will see you in a few minutes.

Customer: Thank you very much. You are very kind.

Clerk: You are welcome, madam. Have a nice day!
\end{textebox}

\textbf{Glossaire :} \eng{cheque} (n.) : chèque ; \eng{by the way} : à propos ; \eng{to take a seat} : s'asseoir ; \eng{you are welcome} : je vous en prie.

\begin{exoresolu}[Repérage et réemploi du lexique]
Énoncé — Complétez les phrases avec le mot qui convient : \eng{branch}, \eng{clerk}, \eng{borrower}, \eng{ATM}, \eng{manager}, \eng{vault}, \eng{customer}.
\begin{enumerate}
\item The \ldots\ldots\ldots{} served me at the counter very politely.
\item There was no money left in the \ldots\ldots\ldots{}, so I could not withdraw any cash.
\item The \ldots\ldots\ldots{} of the bank must repay the loan with interest.
\item This bank has a new \ldots\ldots\ldots{} in Bamenda, near the commercial avenue.
\item Every \ldots\ldots\ldots{} must show an identity card at the counter.
\item The bank \ldots\ldots\ldots{} approved my application for a loan.
\item The gold is kept in the \ldots\ldots\ldots{} at the back of the building.
\end{enumerate}
\tcblower
Solution détaillée :
\begin{enumerate}
\item \eng{clerk} — c'est l'employé qui sert au guichet (\eng{at the counter}).
\item \eng{ATM} — on retire de l'argent au distributeur ; s'il est vide, impossible de retirer.
\item \eng{borrower} — celui qui emprunte doit rembourser le prêt avec des intérêts.
\item \eng{branch} — une banque ouvre une agence locale ; le siège serait \eng{head office}.
\item \eng{customer} — c'est le client qui présente sa carte d'identité au guichet.
\item \eng{manager} — c'est le directeur d'agence qui approuve (\eng{to approve}) une demande de prêt.
\item \eng{vault} — la chambre forte, au fond du bâtiment, protège les objets de valeur.
\end{enumerate}
\end{exoresolu}

\begin{attention}[Pièges fréquents des francophones]
\begin{itemize}
\item Ne dites pas \eng{the bank's office} pour « l'agence » : dites \eng{the branch}.
\item \eng{Counter} ne signifie pas « compteur » (d'électricité) : c'est le guichet. Le compteur électrique se dit \eng{meter}.
\item Attention à l'ordre des mots : \eng{bank manager}, \eng{bank clerk}, \eng{head office} — en anglais, le nom qui précise se place \emph{avant} le nom principal, contrairement au français (« directeur de banque »).
\item N'oubliez pas l'article : on dit \eng{at the counter}, \eng{to the bank}, \eng{from the ATM}.
\end{itemize}
\end{attention}

\begin{savaistu}[The bank on the high street]
En Grande-Bretagne, on entend souvent l'expression \eng{the bank on the high street} pour parler de son agence de quartier. La \eng{high street} est la rue commerçante principale d'une ville britannique — l'équivalent de notre « avenue commerciale » à Douala ou à Yaoundé. Les grandes banques britanniques sont même appelées \eng{high street banks} !
\end{savaistu}

\begin{aretenir}[L'essentiel de la section]
\begin{itemize}
\item Personnes : \eng{customer}, \eng{clerk} (GB) / \eng{teller} (US), \eng{manager}, \eng{accountant}, \eng{borrower}, \eng{lender}, \eng{saver}.
\item Lieux : \eng{bank}, \eng{branch}, \eng{head office}, \eng{counter}, \eng{cash machine} / \eng{ATM}, \eng{safe} / \eng{vault}.
\item Prononciation : \eng{clerk} = « klâk » (GB) ; accent sur la première syllabe pour \eng{CUStomer} et \eng{MAnager}.
\item Le suffixe \cle{-er} transforme un verbe en nom de personne : \eng{borrow → borrower}, \eng{lend → lender}, \eng{save → saver}.
\end{itemize}
\end{aretenir}

Maintenant que vous savez nommer les personnes et les lieux de la banque, la section suivante vous permettra de parler des \cle{comptes} et des \cle{opérations} bancaires : ouvrir un compte, déposer, retirer, virer de l'argent.

\coursec{Champ lexical 2 : accounts and operations}

Dans la section précédente, nous avons découvert les \cle{people} et les \cle{places} du monde bancaire : le \eng{bank manager}, le \eng{cashier}, le \eng{customer}, la \eng{bank}, la \eng{branch} et le \eng{cash machine}. Mais une banque ne sert à rien si l'on ne sait pas ce que l'on y \emph{fait}. Cette section est donc consacrée au cœur du sujet : les \cle{accounts} (les comptes) et les \cle{operations} (les opérations) — ouvrir un compte, déposer de l'argent, retirer de l'argent, faire un virement, vérifier son solde. C'est le vocabulaire indispensable pour parler de la vie économique quotidienne, au Cameroun comme dans les pays anglophones.

\courssub{Observation : a day at the bank}

Lisons d'abord un court texte original. Souligne mentalement tous les mots liés aux comptes et aux opérations bancaires.

\begin{textebox}[A busy morning at the bank]
Last Monday, Mrs Etame went to her bank in Yaoundé very early. She had three things to do. First, she wanted to open a savings account for her daughter, who would start university the following year. The clerk asked her to fill in a form and to make a first deposit of 10,000 francs. Mrs Etame deposited the money into the new account and received a receipt.

Then she went to the counter to withdraw some money from her current account. She made a withdrawal of 50,000 francs because she had to pay the school fees of her two sons. Before leaving, she checked her balance at the cash machine: she still had 120,000 francs left. “Good,” she thought, “I can also transfer 30,000 francs to my brother in Douala. He needs it to pay his rent.”

At the end of the month, Mrs Etame always reads her bank statement carefully. It shows all her deposits, withdrawals and transfers. “You must always know how much money goes in and how much goes out,” she tells her children. “That is the first rule of good management.”
\end{textebox}

\textbf{Glossaire}\par
\begin{center}
\begin{tabularx}{\linewidth}{|l|l|X|}\hline
\rowcolor{popblueL} Mot anglais & Nature & Traduction \\ \hline
clerk & nom & employé(e) de banque, guichetier(-ière) \\ \hline
to fill in a form & expression & remplir un formulaire \\ \hline
receipt & nom & reçu \\ \hline
counter & nom & guichet, comptoir \\ \hline
school fees & nom & frais de scolarité \\ \hline
rent & nom & loyer \\ \hline
statement & nom & relevé (de compte) \\ \hline
management & nom & gestion \\ \hline
\end{tabularx}
\end{center}

\textbf{Questions de compréhension}\par
\begin{enumerate}
\item Why did Mrs Etame open a savings account?
\item How much money did she withdraw, and why?
\item What was her balance after the withdrawal?
\item What does her bank statement show?
\end{enumerate}

\courssub{Les comptes : types de comptes bancaires}

\begin{definition}[Bank account et les types de comptes]
Un \cle{bank account} (compte bancaire) est un compte qu'un client ouvre dans une banque pour y placer son argent. Les principaux types de comptes sont :
\begin{itemize}
\item \cle{current account} (GB) / \cle{checking account} (US) : le compte courant, utilisé pour les opérations quotidiennes (salaire, paiements, retraits).
\item \cle{savings account} : le compte d'épargne, pour mettre de l'argent de côté ; il rapporte des \eng{interest} (intérêts).
\item \cle{joint account} : le compte joint, partagé par deux personnes (par exemple un couple).
\end{itemize}
\end{definition}

\begin{attention}[Current account : attention au faux ami]
\eng{Current account} ne veut pas dire « compte actuel » au sens de « celui qu'on utilise maintenant » : c'est le terme technique britannique pour \emph{compte courant}. De même, \eng{actually} signifie « en fait, en réalité », et non « actuellement » (qui se dit \eng{currently} ou \eng{at present}).
\end{attention}

\begin{exemplebox}[Les comptes en contexte]
\eng{My father has a current account and a savings account.} « Mon père a un compte courant et un compte d'épargne. »\par
\eng{My parents opened a joint account when they got married.} « Mes parents ont ouvert un compte joint quand ils se sont mariés. »\par
\eng{In the United States, people say “checking account” instead of “current account”.} « Aux États-Unis, on dit “checking account” au lieu de “current account”. »
\end{exemplebox}

\courssub{Les opérations : les verbes essentiels}

\begin{definition}[Les cinq opérations de base]
\begin{itemize}
\item \cle{to open an account} : ouvrir un compte ; \cle{to close an account} : fermer un compte.
\item \cle{to deposit money} : déposer de l'argent (le mettre sur le compte).
\item \cle{to withdraw money} : retirer de l'argent (le sortir du compte).
\item \cle{to transfer money} : virer de l'argent, faire un virement.
\item \cle{to check one's balance} : vérifier son solde.
\end{itemize}
\end{definition}

\begin{propriete}[La règle des collocations bancaires]
En anglais, chaque opération a son verbe attitré ; on ne peut pas les intervertir :
\begin{itemize}
\item on dit \eng{to OPEN an account} — jamais \eng{to create an account} dans ce contexte (\eng{create} s'emploie pour un site web ou un profil, pas pour un compte bancaire) ;
\item on dit \eng{to MAKE a deposit} et \eng{to MAKE a withdrawal} — c'est le verbe \eng{make} qui s'emploie avec les noms, alors que \eng{deposit} et \eng{withdraw} s'emploient comme verbes avec l'argent : \eng{to deposit money}, \eng{to withdraw money} ;
\item on dit \eng{to CHECK your balance} — jamais \eng{to verify your balance} (calque du français « vérifier »).
\end{itemize}
\end{propriete}

\begin{methode}[Mémoriser : verbe + préposition + nom]
Pour ne plus jamais hésiter, apprends chaque verbe \emph{avec sa préposition et son nom} comme un bloc unique :
\begin{enumerate}
\item \eng{to deposit money INTO an account} — déposer de l'argent \emph{sur} un compte.
\item \eng{to withdraw money FROM an account} — retirer de l'argent \emph{d'}un compte.
\item \eng{to transfer money TO another account} — virer de l'argent \emph{vers} un autre compte.
\end{enumerate}
Répète ces trois blocs à voix haute, puis fabrique ta propre phrase pour chacun. La préposition fait partie du mot : sans elle, la phrase est fausse ou incompréhensible.
\end{methode}

\courssub{Les noms associés aux opérations}

Chaque verbe a un nom correspondant. Attention : la formation n'est pas toujours régulière.

\begin{center}
\begin{tabularx}{\linewidth}{|X|X|X|}\hline
\rowcolor{popblueL} Verbe & Nom & Traduction du nom \\ \hline
to deposit & a deposit & un dépôt, un versement \\ \hline
to withdraw & a withdrawal & un retrait \\ \hline
to transfer & a transfer & un virement, un transfert \\ \hline
to pay & a payment & un paiement \\ \hline
to save & savings & l'épargne, des économies \\ \hline
--- & a balance & un solde \\ \hline
--- & a statement & un relevé de compte \\ \hline
--- & a receipt & un reçu \\ \hline
\end{tabularx}
\end{center}

\begin{attention}[Withdrawal : orthographe et prononciation]
Le nom de \eng{withdraw} est \cle{withdrawal}, avec un \emph{-al} final : w-i-t-h-d-r-a-w-a-l. Les francophones écrivent souvent \eng{withdraw} pour le nom — erreur ! Prononciation : « ouiz-dro-ol » (verbe : « ouiz-dro »). De même, \eng{receipt} (reçu) se prononce « ri-site » : le \emph{p} est muet, comme dans \eng{psychology}.
\end{attention}

\courssub{Money IN, money OUT : la logique du compte}

Un compte bancaire fonctionne comme un réservoir : de l'argent entre (\cle{money in}) et de l'argent sort (\cle{money out}). Le \cle{balance} (solde) est ce qui reste.

\begin{popfigure}
\begin{tikzpicture}
\node[draw=popgreen, fill=popgreenL, rounded corners, thick, minimum width=5.6cm, minimum height=0.9cm, align=center] (in) at (-3.4,0) {\textbf{Money IN}\\ \emph{l'argent qui entre}};
\node[draw=poporange, fill=poporangeL, rounded corners, thick, minimum width=5.6cm, minimum height=0.9cm, align=center] (out) at (3.4,0) {\textbf{Money OUT}\\ \emph{l'argent qui sort}};
\node[align=left, text=popdark] at (-3.4,-1.4) {a deposit (un dépôt)};
\node[align=left, text=popdark] at (-3.4,-2.1) {a transfer in (virement reçu)};
\node[align=left, text=popdark] at (-3.4,-2.8) {interest (les intérêts)};
\node[align=left, text=popdark] at (3.4,-1.4) {a withdrawal (un retrait)};
\node[align=left, text=popdark] at (3.4,-2.1) {a transfer out (virement émis)};
\node[align=left, text=popdark] at (3.4,-2.8) {charges, fees (les frais)};
\draw[-{Stealth}, very thick, popgreen] (-3.4,-0.6) -- (-3.4,-1.0);
\draw[-{Stealth}, very thick, poporange] (3.4,-0.6) -- (3.4,-1.0);
\node[draw=popblue, fill=popblueL, rounded corners, thick, align=center, minimum width=4.2cm] at (0,-3.9) {\textbf{BALANCE} = money in -- money out\\ \emph{le solde}};
\end{tikzpicture}
\legende{Les deux mouvements d'un compte bancaire : ce qui entre (en vert) et ce qui sort (en orange).}
\end{popfigure}

\begin{exemplebox}[Exemples traduits]
\eng{She deposited 20,000 francs into her savings account.} « Elle a déposé 20 000 francs sur son compte d'épargne. »\par
\eng{I need to check my balance before I pay.} « Je dois vérifier mon solde avant de payer. »\par
\eng{He made a withdrawal of 15,000 francs at the cash machine.} « Il a fait un retrait de 15 000 francs au distributeur. »\par
\eng{The bank transferred the money to her brother's account.} « La banque a viré l'argent sur le compte de son frère. »\par
\eng{Always keep the receipt after a deposit.} « Garde toujours le reçu après un dépôt. »
\end{exemplebox}

\begin{attention}[Balance : solde, pas équilibre !]
Dans le contexte bancaire, \cle{balance} signifie \emph{solde}, c'est-à-dire le montant restant sur le compte : \eng{My balance is low.} « Mon solde est bas. » Le sens « équilibre » existe aussi, mais dans un autre contexte : \eng{She lost her balance and fell.} « Elle a perdu l'équilibre et elle est tombée. »\par
Piège inverse : le mot français \emph{solde} au sens de « réduction de prix » ne se dit \textbf{pas} \eng{balance} mais \cle{sale} ! \eng{The shoes are on sale.} « Les chaussures sont en solde. »
\end{attention}

\courssub{Verbes irréguliers du monde bancaire}

Trois verbes irréguliers sont indispensables pour parler d'argent. Apprends-les par cœur avec leurs trois formes.

\begin{center}
\begin{tabular}{|l|l|l|l|}\hline
\rowcolor{popblueL} Base verbale & Prétérit & Participe passé & Traduction \\ \hline
withdraw & withdrew & withdrawn & retirer \\ \hline
lend & lent & lent & prêter \\ \hline
pay & paid & paid & payer \\ \hline
\end{tabular}
\end{center}

\begin{exemplebox}[Les verbes irréguliers en phrases]
\eng{Yesterday, I withdrew 10,000 francs from my account.} « Hier, j'ai retiré 10 000 francs de mon compte. »\par
\eng{She has withdrawn all her savings.} « Elle a retiré toutes ses économies. »\par
\eng{The bank lent him 500,000 francs last year.} « La banque lui a prêté 500 000 francs l'année dernière. »\par
\eng{Have you paid the school fees yet?} « As-tu déjà payé les frais de scolarité ? »
\end{exemplebox}

\begin{attention}[Lend ou borrow ? Le grand classique]
Les francophones confondent souvent \cle{lend} et \cle{borrow} parce que le français « prêter » couvre parfois les deux idées dans la tête des élèves :
\begin{itemize}
\item \eng{to lend something TO somebody} : prêter quelque chose \emph{à} quelqu'un (on donne temporairement). \eng{Can you lend me your pen?} « Peux-tu me prêter ton stylo ? »
\item \eng{to borrow something FROM somebody} : emprunter quelque chose \emph{à} quelqu'un (on reçoit temporairement). \eng{Can I borrow your pen?} « Puis-je t'emprunter ton stylo ? »
\end{itemize}
La banque \emph{lends} ; le client \emph{borrows}. Et attention à l'orthographe du participe : \eng{paid}, jamais \eng{payed} !
\end{attention}

\courssub{Exercice résolu et entraînement}

\begin{exoresolu}[Phrases à trous : la vie scolaire et la banque]
Complète avec : \emph{deposit, withdraw, balance, transfer, statement, opened, paid, lent}.\par
1. My mother \ldots\ a savings account for me last month.\par
2. The school received my \ldots\ of 25,000 francs this morning.\par
3. I want to \ldots\ 5,000 francs from my account to buy my books.\par
4. Before you pay, always check your \ldots.\par
5. My uncle will \ldots\ the money to the school's account directly.\par
6. The bank \ldots\ my father the money for my school fees.\par
7. I have already \ldots\ the examination fees.\par
8. My bank \ldots\ shows all my operations for the month.
\tcblower
\textbf{Corrigé détaillé :}\par
1. \textbf{opened} — on dit \eng{to open an account} (collocation obligatoire ; prétérit car \eng{last month}).\par
2. \textbf{deposit} — l'école a reçu un dépôt (nom, après le possessif \eng{my}).\par
3. \textbf{withdraw} — \eng{to withdraw money from an account} : retirer de l'argent pour acheter des livres.\par
4. \textbf{balance} — \eng{to check your balance} : vérifier son solde avant de payer.\par
5. \textbf{transfer} — \eng{to transfer money to an account} : faire un virement direct.\par
6. \textbf{lent} — prétérit de \eng{lend} : la banque a prêté l'argent des frais de scolarité.\par
7. \textbf{paid} — participe passé de \eng{pay} avec \eng{have} (present perfect ; attention : \eng{paid}, pas \eng{payed}).\par
8. \textbf{statement} — le relevé de compte montre toutes les opérations du mois.
\end{exoresolu}

\begin{exercice}[À toi de jouer : production guidée]
\textbf{A. Traduis en anglais :}\par
1. Elle a déposé 30 000 francs sur son compte d'épargne.\par
2. Je dois retirer de l'argent pour payer mes frais de scolarité.\par
3. Mes parents ont ouvert un compte joint.\par
4. Vérifie ton solde avant de faire un virement.\par
\textbf{B. Production personnelle :} écris trois phrases vraies ou imaginées sur ta vie scolaire en utilisant : \eng{deposit}, \eng{withdraw}, \eng{transfer} (par exemple : payer les frais de scolarité, recevoir une bourse — \eng{a scholarship} —, économiser pour les fournitures).
\end{exercice}

\begin{corrige}[Exercice « À toi de jouer »]
\textbf{A. Traductions attendues :}\par
1. \eng{She deposited 30,000 francs into her savings account.}\par
2. \eng{I need to withdraw some money to pay my school fees.}\par
3. \eng{My parents opened a joint account.}\par
4. \eng{Check your balance before you make a transfer.}\par
\textbf{B. Exemples de production :}\par
\eng{Every term, my father deposits money into my account to pay the school fees.}\par
\eng{When I received my scholarship, I withdrew 10,000 francs to buy my school supplies.}\par
\eng{The ministry transferred the scholarship directly to the students' accounts.}
\end{corrige}

\begin{aretenir}[L'essentiel de la section]
\begin{itemize}
\item Types de comptes : \eng{current/checking account} (compte courant), \eng{savings account} (compte d'épargne), \eng{joint account} (compte joint).
\item Collocations obligatoires : \eng{to OPEN/CLOSE an account}, \eng{to MAKE a deposit/a withdrawal}, \eng{to CHECK your balance}.
\item Prépositions : \eng{deposit INTO}, \eng{withdraw FROM}, \eng{transfer TO}.
\item \eng{Balance} = solde (pas équilibre) ; \eng{sale} = solde (réduction).
\item Verbes irréguliers : \eng{withdraw–withdrew–withdrawn}, \eng{lend–lent–lent}, \eng{pay–paid–paid}.
\end{itemize}
\end{aretenir}

\coursec{Champ lexical 3 : payment methods, credit and savings}

Dans la section précédente, nous avons découvert les \cle{comptes} (\eng{current account}, \eng{savings account}) et les \cle{opérations bancaires} (\eng{to deposit}, \eng{to withdraw}, \eng{balance}). Mais une fois l'argent déposé sur un compte, comment le dépense-t-on ? Comment emprunte-t-on pour acheter une moto ou un terrain ? Et comment fait-on fructifier ses économies ? C'est l'objet de cette troisième section : les \cle{moyens de paiement}, le \cle{crédit} et l'\cle{épargne}. Ce lexique est absolument central au Bac, tant à l'écrit (compréhension de textes économiques) qu'à l'oral.

\courssub{Observation : a text to start with}

Lisons d'abord ce court texte original. Soulignez mentalement tous les mots qui appartiennent au champ lexical de la banque.

\begin{textebox}[Banking in Douala: a young trader's story]
Aminatou runs a small clothing shop near Marché Mboppi in Douala. For years, she kept all her money in \textbf{cash} at home, which was risky. Last year, she finally opened a \textbf{savings account} and received a \textbf{debit card}. Now her customers can pay by \textbf{bank transfer} or even by \textbf{mobile money}, which is extremely popular in Cameroon. She rarely accepts \textbf{cheques} anymore, because some of them bounce.

Business was good, so Aminatou decided to expand. She \textbf{applied for a loan} of two million francs CFA. The bank agreed to \textbf{lend} her the money at an \textbf{interest rate} of eight per cent. She now \textbf{repays} the loan every month, and she is careful never to have an \textbf{overdraft}, because the \textbf{charges} are very high. “I \textbf{borrowed} money from the bank, not from my family,” she laughs, “so nobody can quarrel with me at Christmas!”

Aminatou also tries to \textbf{save} a little every month. “My \textbf{savings} are my security,” she explains. “One day, I would like to \textbf{invest} in a second shop in Bafoussam. A good \textbf{investment} today means a better life for my children tomorrow.” She only regrets one thing: the bank \textbf{fees} and the \textbf{commission} on international transfers, which she finds too expensive when she orders clothes from Nigeria.
\end{textebox}

\textbf{Glossaire :}

\begin{center}
\begin{tabularx}{\linewidth}{|l|X|}\hline
\rowcolor{popblueL} \textbf{English} & \textbf{Français} \\ \hline
cash (n.) & l'argent liquide, les espèces \\ \hline
debit card (n.) & la carte bancaire (à débit immédiat) \\ \hline
bank transfer (n.) & le virement bancaire \\ \hline
mobile money (n.) & l'argent mobile (Orange Money, MTN MoMo) \\ \hline
cheque (n.) & le chèque ; to bounce = être rejeté (chèque sans provision) \\ \hline
to apply for a loan & demander un prêt \\ \hline
interest rate (n.) & le taux d'intérêt \\ \hline
to repay (n. repayment) & rembourser (le remboursement) \\ \hline
overdraft (n.) & le découvert bancaire \\ \hline
charges / fees (n.) & les frais (bancaires) \\ \hline
to save / savings (n.) & économiser / les économies, l'épargne \\ \hline
to invest / investment (n.) & investir / l'investissement \\ \hline
commission (n.) & la commission \\ \hline
\end{tabularx}
\end{center}

\textbf{Questions de compréhension (oral) :}
\begin{enumerate}
\item Why was it risky for Aminatou to keep her money at home?
\item Which payment methods do her customers use?
\item How much did she borrow, and at what interest rate?
\item Why does she avoid an overdraft?
\item What is her project for the future?
\end{enumerate}

\courssub{Payment methods : les moyens de paiement}

\begin{definition}[Payment methods — les moyens de paiement]
\eng{cash} = l'argent liquide (espèces, billets et pièces) ; \eng{bank card / debit card} = la carte bancaire (le montant est débité immédiatement du compte) ; \eng{credit card} = la carte de crédit (on paie plus tard, la banque avance l'argent) ; \eng{cheque} (\emph{américain :} \eng{check}) = le chèque ; \eng{cheque book} (\emph{am.} \eng{checkbook}) = le chéquier / carnet de chèques ; \eng{bank transfer} = le virement ; \eng{mobile money} = le paiement par téléphone mobile, très répandu au Cameroun.
\end{definition}

\begin{exemplebox}[Collocations utiles avec les moyens de paiement]
\begin{itemize}
\item \eng{to pay in cash} = payer en liquide. \eng{The bendskin driver only accepts cash.} « Le moto-taxi n'accepte que le liquide. »
\item \eng{to pay by card / by cheque / by transfer} = payer par carte / par chèque / par virement. \eng{Can I pay by card?} « Puis-je payer par carte ? »
\item \eng{to write a cheque} = faire / rédiger un chèque ; \eng{to cash a cheque} = encaisser un chèque.
\item \eng{to withdraw cash from an ATM} = retirer du liquide à un distributeur.
\item \eng{to send money by mobile money} = envoyer de l'argent par mobile money.
\end{itemize}
\end{exemplebox}

\begin{attention}[Prépositions : in cash, by card, by cheque]
On dit \eng{to pay \textbf{in} cash} mais \eng{to pay \textbf{by} card / cheque / transfer}. Ne dites jamais \eng{pay by cash} ni \eng{pay in card}. De même : \eng{to pay \textbf{for} something} (payer quelque chose) : \eng{She paid for the books.} « Elle a payé les livres. » — sans \eng{for} si l'objet est une personne : \eng{She paid the taxi driver.}
\end{attention}

\courssub{Credit : le crédit — borrow, lend, interest, debt}

\begin{definition}[Le vocabulaire du crédit]
\eng{loan} (n.) = le prêt, l'emprunt ; \eng{to take out a loan} = contracter un prêt ; \eng{interest} (n.) = les intérêts ; \eng{interest rate} = le taux d'intérêt ; \eng{debt} (n.) = la dette ; \eng{to be in debt} = être endetté ; \eng{overdraft} (n.) = le découvert ; \eng{borrower} (n.) = l'emprunteur ; \eng{lender} (n.) = le prêteur.
\end{definition}

\begin{definition}[Overdraft]
An \cle{overdraft} is a situation where you spend more money than you have in your account — the bank lets your account go below zero, but charges you for it. = un découvert bancaire. \eng{He has an overdraft of 50,000 francs.} « Il a un découvert de 50 000 francs. »
\end{definition}

\begin{propriete}[La règle d'or : borrow FROM / lend TO]
Ces deux verbes décrivent la \emph{même opération} vue des deux côtés :
\begin{itemize}
\item \cle{to borrow something FROM somebody} = emprunter quelque chose \emph{à} quelqu'un (celui qui reçoit). \eng{I borrowed 10,000 francs from my brother.} « J'ai emprunté 10 000 francs à mon frère. »
\item \cle{to lend something TO somebody} (ou \eng{to lend somebody something}) = prêter quelque chose \emph{à} quelqu'un (celui qui donne). \eng{My brother lent me 10,000 francs.} / \eng{My brother lent 10,000 francs to me.} « Mon frère m'a prêté 10 000 francs. »
\end{itemize}
Formes irrégulières : \eng{lend – lent – lent} ; \eng{borrow} est régulier (\eng{borrowed, borrowed}).
\end{propriete}

\begin{attention}[Piège majeur des francophones : ne confondez pas borrow et lend]
En français, « prêter » et « emprunter » sont proches ; en anglais, l'erreur est très grave de sens. \eng{Can you borrow me your pen?} est FAUX : il faut \eng{Can you \textbf{lend} me your pen?} « Peux-tu me prêter ton stylo ? ». Astuce : \emph{borrow} = je reçois ; \emph{lend} = je donne. De même, \eng{interest} signifie « intérêts » bancaires \emph{et} « intérêt » (curiosité), mais attention : « être intéressé par » se dit \eng{to be interested \textbf{in}}, jamais \eng{interested by}.
\end{attention}

\begin{exemplebox}[Le crédit en contexte — exemples traduits]
\begin{itemize}
\item \eng{He took out a loan to buy a taxi.} « Il a contracté un prêt pour acheter un taxi. »
\item \eng{The bank charges high fees for international transfers.} « La banque prélève des frais élevés pour les virements internationaux. »
\item \eng{She repays her loan every month.} « Elle rembourse son prêt chaque mois. »
\item \eng{The interest rate is too high for small traders.} « Le taux d'intérêt est trop élevé pour les petits commerçants. »
\item \eng{After the funeral expenses, the family was heavily in debt.} « Après les frais de funérailles, la famille était très endettée. »
\item \eng{The bank lent the cooperative 5 million francs.} « La banque a prêté 5 millions de francs à la coopérative. »
\end{itemize}
\end{exemplebox}

\begin{popfigure}
\begin{center}
\begin{tikzpicture}[font=\small]
\draw[-{Stealth}, very thick, popblue] (0,0) -- (12.6,0);
\foreach \x/\txt/\col in {0.4/{apply for\\a loan}/popblue, 4.4/{the bank\\lends money}/popgreen, 8.4/{the borrower\\repays monthly}/poporange, 12.2/{with\\interest}/poppurple}{
  \fill[\col] (\x,0) circle (0.09);
  \node[draw=\col, fill=\col!15, rounded corners=2pt, align=center, anchor=south] at (\x, 0.4) {\txt};
}
\node[below=0.25cm of {6.3,0}, text=popmuted] {\emph{The credit cycle — le cycle du crédit}};
\end{tikzpicture}
\end{center}
\legende{Le cycle du crédit : demander un prêt, la banque prête, l'emprunteur rembourse chaque mois, avec des intérêts.}
\end{popfigure}

\courssub{Savings and investment : l'épargne et l'investissement}

\begin{definition}[La famille de save et invest]
\eng{to save} (v.) = économiser, épargner ; \eng{savings} (n., toujours au pluriel) = les économies, l'épargne ; \eng{saver} (n.) = l'épargnant ; \eng{savings account} = le compte d'épargne ; \eng{to invest} (v.) = investir ; \eng{investment} (n.) = l'investissement ; \eng{investor} (n.) = l'investisseur.
\end{definition}

\begin{exemplebox}[L'épargne en contexte]
\begin{itemize}
\item \eng{Young people should save at least a small amount every month.} « Les jeunes devraient épargner au moins une petite somme chaque mois. »
\item \eng{She keeps her savings in a savings account.} « Elle garde ses économies sur un compte d'épargne. »
\item \eng{They invested their savings in a cocoa farm near Ebolowa.} « Ils ont investi leurs économies dans une plantation de cacao près d'Ebolowa. »
\item \eng{Education is the best investment.} « L'éducation est le meilleur investissement. »
\end{itemize}
\end{exemplebox}

\begin{attention}[Faux amis : economy, savings, to invest]
\begin{itemize}
\item \eng{savings} (pluriel) = les économies d'argent ; \eng{economy} = l'économie (d'un pays) ; \eng{economics} = les sciences économiques ; \eng{economical} = économe (qui évite le gaspillage) ; \eng{economic} = économique (relatif à l'économie).
\item \eng{to save money} = économiser de l'argent, mais \eng{to save a life} = sauver une vie. « Économiser du temps » = \eng{to save time}.
\item \eng{to invest} = investir (argent) ; ne pas confondre avec \eng{to invent} = inventer.
\end{itemize}
\end{attention}

\courssub{Bank charges : les frais bancaires}

\begin{definition}[Charges, fees, commission]
\eng{charges} et \eng{fees} = les frais (bancaires) ; \eng{commission} = la commission (pourcentage prélevé sur une opération) ; \eng{free of charge} = gratuit. \eng{The bank charges a fee for each withdrawal.} « La banque prélève des frais à chaque retrait. » \eng{The agent takes a two per cent commission.} « L'agent prend une commission de deux pour cent. »
\end{definition}

\begin{attention}[Prononciations pièges à l'oral du Bac]
\begin{itemize}
\item \eng{debt} se prononce « \textbf{dett} » : le \emph{b} est muet (comme \eng{doubt}, \eng{subtle}).
\item \eng{receipt} (le reçu) se prononce « \textbf{rissit} » : le \emph{p} est muet.
\item \eng{mortgage} se prononce « \textbf{morguidge} » : le \emph{t} est muet.
\item \eng{cheque} se prononce exactement comme \eng{check} : « \textbf{tchèk} ».
\end{itemize}
\end{attention}

\begin{savaistu}[(Complément Bac) What is a mortgage?]
A \cle{mortgage} is a special long-term loan — often for 20 or 25 years — that people take out to buy a house or a flat. The house itself is the bank's guarantee: if the borrower cannot repay, the bank can take the house. Mortgages are extremely common in British and American texts, so you will often meet expressions like \eng{to take out a mortgage} (contracter un prêt immobilier), \eng{to pay off a mortgage} (finir de rembourser) and \eng{monthly mortgage payments} (mensualités). In Cameroon, many people build their houses progressively instead, but the word is essential for reading English newspapers.
\end{savaistu}

\courssub{Tableau récapitulatif du champ lexical}

\begin{center}
\begin{tabularx}{\linewidth}{|l|X|l|}\hline
\rowcolor{popblueL} \textbf{English} & \textbf{Français} & \textbf{Famille / collocation} \\ \hline
cash & l'argent liquide & to pay in cash \\ \hline
debit / credit card & la carte bancaire / de crédit & to pay by card \\ \hline
cheque / cheque book & le chèque / le chéquier & to write / cash a cheque \\ \hline
bank transfer & le virement & to make a transfer \\ \hline
loan & le prêt & to take out / repay a loan \\ \hline
to borrow (from) & emprunter (à) & borrower (n.) \\ \hline
to lend (to) & prêter (à) & lender (n.) ; lent, lent \\ \hline
interest / interest rate & les intérêts / le taux d'intérêt & high / low interest \\ \hline
overdraft & le découvert & to have an overdraft \\ \hline
debt & la dette & to be in debt (« dett ») \\ \hline
mortgage & le prêt immobilier & (« morguidge ») \\ \hline
to save / savings / saver & épargner / l'épargne / l'épargnant & savings account \\ \hline
to invest / investment & investir / l'investissement & investor (n.) \\ \hline
charges / fees / commission & les frais / la commission & free of charge \\ \hline
receipt & le reçu & (« rissit ») \\ \hline
\end{tabularx}
\end{center}

\courssub{Exercice résolu et entraînement}

\begin{exoresolu}[Borrow or lend ? Complétez et traduisez]
Énoncé — Complete with the correct form of \emph{borrow} or \emph{lend}, then translate into French.

1. Can you \dots\ me your bicycle this afternoon?

2. I \dots\ 20,000 francs from the bank to pay my school fees.

3. Banks \dots\ money to businesses and charge interest.

4. She never \dots\ her books to anyone.
\tcblower
Solution détaillée :

1. \eng{Can you \textbf{lend} me your bicycle this afternoon?} — « Peux-tu me prêter ton vélo cet après-midi ? » (celui qui parle \emph{reçoit} l'objet, donc l'autre personne \emph{donne} : \emph{lend}).

2. \eng{I \textbf{borrowed} 20,000 francs from the bank to pay my school fees.} — « J'ai emprunté 20 000 francs à la banque pour payer mes frais de scolarité. » (le sujet \emph{reçoit} : \emph{borrow from} ; passé : \emph{borrowed}).

3. \eng{Banks \textbf{lend} money to businesses and charge interest.} — « Les banques prêtent de l'argent aux entreprises et prélèvent des intérêts. » (les banques \emph{donnent} : \emph{lend to}).

4. \eng{She never \textbf{lends} her books to anyone.} — « Elle ne prête jamais ses livres à personne. » (elle \emph{donne} : \emph{lend} ; présent simple, 3\textsuperscript{e} personne : \emph{lends}).
\end{exoresolu}

\begin{exercice}[À vous de jouer]
1. Traduisez en anglais : (a) « J'ai contracté un prêt pour acheter une moto. » (b) « Peux-tu me prêter 5 000 francs ? Je te rembourserai lundi. » (c) « La banque prélève des frais élevés sur les virements. » (d) « Il épargne un peu d'argent chaque mois sur son compte d'épargne. »

2. Prononciation : lisez à voix haute \eng{debt}, \eng{receipt}, \eng{mortgage}, \eng{cheque} en respectant les lettres muettes.

3. Production guidée : en cinq phrases, racontez comment un jeune commerçant de Bamenda utilise sa carte bancaire, le mobile money, un prêt et son compte d'épargne. Utilisez au moins six mots du tableau récapitulatif.
\end{exercice}

\begin{corrige}[Exercice : éléments de correction]
1. (a) \eng{I took out a loan to buy a motorbike.} (b) \eng{Can you lend me 5,000 francs? I will repay you on Monday.} (c) \eng{The bank charges high fees on transfers.} (d) \eng{He saves a little money every month in his savings account.}

2. « dett », « rissit », « morguidge », « tchèk ».

3. Modèle possible : \eng{Manka is a young trader in Bamenda. His customers pay him by mobile money or by bank transfer. He pays his suppliers with his debit card. Last year, he borrowed money from a microfinance bank and took out a small loan. He repays it every month with interest. He also saves money in a savings account because he wants to invest in a bigger shop.}
\end{corrige}

\begin{aretenir}[L'essentiel de la section]
\begin{itemize}
\item Moyens de paiement : \eng{cash} (\eng{pay \textbf{in} cash}), \eng{card}, \eng{cheque}, \eng{transfer}, \eng{mobile money} (\eng{pay \textbf{by} card}).
\item Crédit : \eng{to borrow FROM} (recevoir) / \eng{to lend TO} (donner) ; \eng{loan}, \eng{interest rate}, \eng{overdraft}, \eng{debt}, \eng{mortgage}.
\item Épargne : \eng{to save}, \eng{savings}, \eng{saver}, \eng{savings account}, \eng{to invest}, \eng{investment}.
\item Frais : \eng{charges}, \eng{fees}, \eng{commission}.
\item Prononciation : \eng{debt} « dett », \eng{receipt} « rissit », \eng{mortgage} « morguidge ».
\end{itemize}
\end{aretenir}

\coursec{Faux amis, familles de mots et expressions}

Dans la section précédente, nous avons étudié le vocabulaire des moyens de paiement, du crédit et de l'épargne (\eng{payment methods, credit and savings}). Nous avons rencontré des mots comme \eng{to save}, \eng{a loan}, \eng{a bank note}. Mais attention : certains de ces mots ressemblent à des mots français sans avoir le même sens ! Ce sont les fameux \cle{faux amis} (\eng{false friends} ou \eng{false cognates}). Dans cette section, nous allons d'abord débusquer les faux amis du monde bancaire, puis corriger les calques du français les plus fréquents, et enfin apprendre à construire des \cle{familles de mots} (\eng{word families}) grâce aux suffixes — une compétence décisive pour l'épreuve d'anglais du Baccalauréat.

\courssub{Les faux amis du monde bancaire}

\begin{definition}[Faux ami (\eng{false friend})]
Un \cle{faux ami} est un mot anglais qui ressemble par sa forme à un mot français, mais dont le sens est différent. Exemple : \eng{a bank note} ressemble à « une note de banque », mais signifie « un billet de banque ». À l'inverse, certains mots français trompeurs n'ont pas l'équivalent anglais qu'on imagine : « un banc » n'est pas \eng{a bank} mais \eng{a bench}.
\end{definition}

Observons d'abord ces deux phrases authentiques :

\eng{The lifeguard saved the child from drowning.} « Le maître-nageur a sauvé l'enfant de la noyade. »

\eng{She saved 5,000 francs every month to pay her school fees.} « Elle économisait 5 000 francs chaque mois pour payer ses frais de scolarité. »

Le même verbe \eng{to save} a ici deux sens différents. C'est le piège numéro un des francophones.

\begin{attention}[Faux ami majeur : \eng{to save}]
\eng{To save money} signifie « économiser / épargner de l'argent » ; \eng{to save somebody} signifie « sauver quelqu'un » (d'un danger). Ne traduisez jamais « sauver » par \eng{to save} quand il s'agit d'un secours dans un contexte où l'on veut insister sur l'opération de sauvetage : on préfère souvent \eng{to rescue}.
\begin{itemize}
\item \eng{She saved 5,000 francs every month.} « Elle économisait 5 000 francs chaque mois. »
\item \eng{The lifeguard saved the child.} « Le maître-nageur a sauvé l'enfant. »
\item \eng{Firefighters rescued two people from the burning building.} « Les pompiers ont sauvé deux personnes de l'immeuble en feu. »
\end{itemize}
Autre piège : « sauvegarder un fichier » (informatique) se dit bien \eng{to save a file}, mais « faire des économies » en vue d'un achat se dit \eng{to save up} : \eng{I am saving up for a new phone.} « J'économise pour m'acheter un nouveau téléphone. »
\end{attention}

\begin{center}
\begin{tabularx}{\linewidth}{|X|X|X|}
\hline
\rowcolor{popblueL} Français & English & Remarque \\
\hline
économiser / épargner (de l'argent) & to save (money), to save up & \eng{He saved 20,000 francs.} \\
\hline
sauver (quelqu'un d'un danger) & to save, to rescue & \eng{They rescued the driver.} \\
\hline
un billet de banque & a bank note (GB), a bill (US) & pas « a note of bank » \\
\hline
une banque & a bank & \eng{I went to the bank.} \\
\hline
un banc (siège) & a bench & \eng{We sat on a bench.} \\
\hline
un banc de sable / de poissons & a sandbank / a school of fish & faux sens de \eng{bank} \\
\hline
une économie (somme épargnée) & savings (au pluriel) & \eng{my savings} = « mes économies » \\
\hline
\end{tabularx}
\end{center}

\begin{exemplebox}[Deux faux amis en contexte]
\eng{The cashier counted the bank notes carefully.} « Le caissier a compté les billets attentivement. »

\eng{The old man was sitting on a bench in front of the bank.} « Le vieil homme était assis sur un banc devant la banque. »

Remarquez la différence de prononciation : \eng{bank} se prononce « banke » (a bref), \eng{bench} se prononce « bentch ». Confondre les deux à l'oral peut changer tout le sens de la phrase !
\end{exemplebox}

\courssub{Les calques du français à éviter}

Un \cle{calque} est une traduction mot à mot du français qui produit une phrase anglaise incorrecte ou non naturelle. En voici quatre très fréquents dans le domaine bancaire.

\begin{attention}[Calques à bannir]
\begin{itemize}
\item « Faire un crédit » ne se dit pas \eng{to make a credit}. On dit \eng{to take out a loan} ou \eng{to get a loan} : \eng{They took out a loan to build their house.} « Ils ont fait un crédit pour construire leur maison. »
\item « Toucher de l'argent » (au sens de recevoir un paiement/salaire) ne se dit pas \eng{to touch money}. On dit \eng{to receive money} ou \eng{to be paid} : \eng{Workers are paid at the end of the month.} « Les travailleurs touchent leur salaire à la fin du mois. »
\item « Retirer de l'argent » (au distributeur) se dit \eng{to withdraw money} (et non \eng{to retire money} ! \eng{to retire} = partir à la retraite) : \eng{I withdrew 10,000 francs from the ATM.} « J'ai retiré 10 000 francs au distributeur. »
\item « Verser de l'argent sur un compte » se dit \eng{to deposit money into an account} (ou \eng{to pay money into an account}) : \eng{She deposited her salary into her account.} « Elle a versé son salaire sur son compte. »
\end{itemize}
\end{attention}

\begin{propriete}[Prononciation des verbes d'opération bancaire]
\begin{itemize}
\item \eng{to withdraw} : « wi-dro » (le \emph{w} ne se prononce pas, accent sur \emph{draw}).
\item \eng{to deposit} : « di-po-zite » (accent sur \emph{po}).
\item \eng{to lend} : « lende » (comme \eng{send}).
\item \eng{to borrow} : « bo-ro » (le \emph{ow} sonne \emph{o}).
\item \eng{to invest} : « in-veste » (accent sur \emph{vest}).
\end{itemize}
\end{propriete}

\courssub{Les familles de mots : les suffixes bancaires}

Observons ces phrases :

\eng{My father is a banker. He works in banking. He advised us to open a bank account.} « Mon père est banquier. Il travaille dans la banque. Il nous a conseillé d'ouvrir un compte bancaire. »

Trois mots de la même famille : \eng{bank}, \eng{banker}, \eng{banking}. L'anglais construit ses familles de mots avec des \cle{suffixes} réguliers.

\begin{propriete}[Formation des familles de mots : les suffixes fréquents]
\begin{itemize}
\item \cle{-er} (ou \cle{-or}) : la personne qui fait l'action → \eng{banker} (banquier), \eng{lender} (prêteur), \eng{saver} (épargnant), \eng{borrower} (emprunteur), \eng{employer} (employeur), \eng{investor} (investisseur).
\item \cle{-ing} : l'activité, le domaine → \eng{banking} (le secteur bancaire), \eng{borrowing} (l'emprunt, l'action d'emprunter), \eng{savings} (les économies — toujours au pluriel dans ce sens).
\item \cle{-ment} : l'action ou son résultat → \eng{payment} (paiement), \eng{investment} (investissement), \eng{agreement} (accord).
\item \cle{-able} : ce qui peut être fait → \eng{payable} (payable), \eng{bankable} (bancable, rentable pour une banque).
\end{itemize}
Attention à l'orthographe : \eng{pay} → \eng{payment} (on garde le \emph{y} devant \emph{-ment}), mais \eng{payable} garde aussi le \emph{y}. De même \eng{lay} → \eng{layout} (différent), mais pour les verbes en \emph{-y} précédés d'une consonne, le \emph{y} ne change pas devant \emph{-ment} ni \emph{-able} (sauf \eng{try} → \eng{trial}).
\end{propriete}

\begin{center}
\begin{tabular}{|l|l|l|l|}
\hline
\rowcolor{popblueL} Verbe & Personne (-er/-or) & Action / Résultat & Traduction du verbe \\
\hline
to bank & banker & banking & déposer en banque / avoir un compte \\
\hline
to lend & lender & loan (nom irrégulier) & prêter (à qqn) \\
\hline
to borrow & borrower & borrowing & emprunter (à qqn) \\
\hline
to save & saver & savings (pluriel) & économiser \\
\hline
to pay & payer (payeur rare) & payment & payer \\
\hline
to invest & investor & investment & investir \\
\hline
\end{tabular}
\end{center}

\begin{attention}[\eng{To lend} ou \eng{to borrow} ?]
Piège classique ! \eng{To lend} = prêter (on donne) ; \eng{to borrow} = emprunter (on reçoit). Le sens est opposé :
\eng{Can you lend me 2,000 francs?} « Peux-tu me prêter 2 000 francs ? »
\eng{Can I borrow 2,000 francs from you?} « Puis-je t'emprunter 2 000 francs ? »
Ne dites jamais \eng{Can you borrow me money?} — c'est une erreur très fréquente chez les francophones.
Rappel : \eng{lend} (irrégulier) → \eng{lent} → \eng{lent} ; \eng{borrow} (régulier) → \eng{borrowed} → \eng{borrowed}.
\end{attention}

\begin{popfigure}
\begin{center}
\begin{tikzpicture}[
  root/.style={rectangle, rounded corners, draw=popblue, fill=popblue, text=white, font=\bfseries, align=center, minimum width=2.2cm},
  child1/.style={rectangle, rounded corners, draw=poporange, fill=poporangeL, align=center, font=\small},
  child2/.style={rectangle, rounded corners, draw=popgreen, fill=popgreenL, align=center, font=\small},
  grow=right, level distance=3.4cm,
  level 1/.style={sibling distance=1.6cm},
  edge from parent/.style={draw=popline, thick}
]
\node[root, align=center]{BANK\\ « banke »}
  child { node[child2, align=center]{bank statement\\ (relevé de compte)} }
  child { node[child2, align=center]{bank account\\ (compte bancaire)} }
  child { node[child2, align=center]{bank card\\ (carte bancaire)} }
  child { node[child1, align=center]{bankable\\ (bancable)} }
  child { node[child1, align=center]{banking\\ (secteur bancaire)} }
  child { node[child1, align=center]{banker\\ (banquier)} };
\end{tikzpicture}
\end{center}
\legende{Arbre lexical : la famille du mot \eng{bank} (en orange, les dérivés par suffixe ; en vert, les mots composés).}
\end{popfigure}

\courssub{Expressions idiomatiques du monde de l'argent}

Les Anglais utilisent beaucoup d'expressions imagées liées à la banque. En voici quatre à connaître absolument.

\begin{definition}[Expressions idiomatiques bancaires]
\begin{itemize}
\item \eng{To be in the red} : être à découvert (compte débiteur). Les banques écrivaient autrefois les dettes à l'encre rouge.
\item \eng{To be in the black} : avoir un compte créditeur, être en positif (bénéficiaire).
\item \eng{To break the bank} : coûter très cher, coûter une fortune (littéralement « faire sauter la banque »).
\item \eng{(As) safe as the Bank of England} : sûr et certain, sans aucun risque.
\end{itemize}
\end{definition}

\begin{exemplebox}[Les expressions en contexte]
\eng{My account is in the red.} « Mon compte est à découvert. »

\eng{After the cocoa season, the cooperative is finally in the black.} « Après la saison du cacao, la coopérative est enfin bénéficiaire. »

\eng{Don't worry, this car won't break the bank.} « Ne t'inquiète pas, cette voiture ne coûte pas une fortune. »

\eng{Investing in government bonds is as safe as the Bank of England.} « Investir dans les bons du Trésor est on ne peut plus sûr. »
\end{exemplebox}

\begin{savaistu}[Pourquoi « in the red » ?]
Dans les registres comptables traditionnels, les comptables écrivaient les montants négatifs à l'encre rouge et les montants positifs à l'encre noire. Cette pratique a donné naissance aux expressions \eng{in the red} et \eng{in the black}, encore utilisées chaque jour dans la presse économique anglophone. Le vendredi des soldes aux États-Unis s'appelle \eng{Black Friday} parce que, grâce aux ventes massives, les commerçants repassent « dans le noir », c'est-à-dire deviennent bénéficiaires.
\end{savaistu}

\courssub{Lexique en contexte : dialogue et production guidée}

Pour fixer le vocabulaire, étudions un dialogue typique au guichet d'une banque camerounaise (ex. : UBA, Afriland First Bank, BICEC).

\begin{textebox}[Dialogue : Opening a savings account]
\textbf{Bank Clerk} : Good morning, sir. How can I help you?
\textbf{Customer} : Good morning. I would like to open a savings account.
\textbf{Bank Clerk} : Certainly. Do you have your identity card and a proof of address?
\textbf{Customer} : Yes, here is my national ID card and an electricity bill.
\textbf{Bank Clerk} : Thank you. You will need to make an initial deposit of 5,000 francs. The interest rate is 3.5\% per year.
\textbf{Customer} : That sounds fine. I want to deposit 20,000 francs right now.
\textbf{Bank Clerk} : Perfect. Please sign this form. Here is your passbook and your bank card. Your PIN will be sent by SMS.
\textbf{Customer} : Thank you very much. Goodbye.
\textbf{Bank Clerk} : You're welcome. Have a nice day.
\end{textebox}

\begin{experience}[Jeu de rôle : À la banque]
\textbf{Consigne} : Par binômes. L'un joue le rôle du \textbf{guichetier} (bank clerk), l'autre celui du \textbf{client} (customer). Utilisez le vocabulaire de la leçon.
\begin{enumerate}
\item \textbf{Scénario A} : Le client veut retirer de l'argent (\eng{withdraw money}) mais son compte est \eng{in the red}. Le guichetier lui propose un découvert autorisé (\eng{overdraft}) ou un prêt (\eng{loan}).
\item \textbf{Scénario B} : Le client veut \eng{take out a loan} pour acheter une moto. Le guichetier demande des garanties (\eng{guarantees / collateral}) et explique le \eng{interest rate} et les \eng{repayment terms} (modalités de remboursement).
\item \textbf{Scénario C} : Le client a perdu sa \eng{bank card}. Il doit la faire opposition (\eng{report the card lost/stolen}) et commander une nouvelle.
\end{enumerate}
\textbf{Expressions utiles} : \eng{I would like to...}, \eng{Could you please...?}, \eng{Here is my...}, \eng{How much would you like to...?}, \eng{You need to provide...}, \eng{The rate is...}, \eng{Please sign here.}
\end{experience}

\begin{methode}[Pour le Bac : les questions de vocabulaire]
Dans les épreuves de compréhension de texte, une question fréquente est : \eng{Find in the text a word from the same family as...} ou \eng{Give the noun/adjective derived from...}. Voici la démarche :
\begin{enumerate}
\item Repérez le mot de base dans la question (ex. \eng{to pay}).
\item Identifiez la catégorie grammaticale demandée : nom de personne ? nom d'action ? adjectif ?
\item Cherchez dans le texte un mot qui commence par la même racine (\eng{pay-}) : \eng{payment}, \eng{payer}, \eng{payable}.
\item Vérifiez que le mot trouvé convient grammaticalement (un nom après \eng{the}, un adjectif avant un nom).
\end{enumerate}
Le mot de la même famille que celui du texte est très souvent la réponse attendue : si le texte contient \eng{investment} et que la question demande le verbe, la réponse est \eng{to invest}.
\end{methode}

\begin{exoresolu}[Dérivez le bon mot de la famille]
Énoncé — Complétez chaque phrase avec le mot de la famille indiquée entre parenthèses :
\begin{enumerate}
\item My uncle is a \ldots\ldots\ldots ; he manages a branch in Douala. (bank)
\item The \ldots\ldots\ldots of the loan must be made before Friday. (pay)
\item Small \ldots\ldots\ldots can borrow money from the microfinance institution. (borrow)
\item She put all her \ldots\ldots\ldots in a savings account. (save)
\item Online \ldots\ldots\ldots is very convenient for young customers. (bank)
\item The bank refused because the project was not \ldots\ldots\ldots . (bank)
\end{enumerate}
\tcblower
Solution détaillée :
\begin{enumerate}
\item \eng{banker} — on demande la personne (suffixe \emph{-er}) : « Mon oncle est banquier. »
\item \eng{payment} — on demande le nom de l'action (suffixe \emph{-ment}) après l'article \eng{the} : « Le paiement (remboursement) du prêt. »
\item \eng{borrowers} — la personne qui emprunte (suffixe \emph{-er}), au pluriel à cause de \eng{Small... can} sans article : « Les petits emprunteurs. »
\item \eng{savings} — le nom « économies », toujours au pluriel dans ce sens : « Elle a placé toutes ses économies. »
\item \eng{banking} — le nom de l'activité (suffixe \emph{-ing}) : \eng{online banking} = « la banque en ligne ».
\item \eng{bankable} — l'adjectif (suffixe \emph{-able}) après \eng{was not} : « le projet n'était pas bancable ».
\end{enumerate}
\end{exoresolu}

\begin{exercice}[À vous de jouer]
\begin{enumerate}
\item Traduisez en anglais : « Mon compte est à découvert, je dois économiser davantage. »
\item Corrigez l'erreur : \eng{The bank borrowed me 50,000 francs.}
\item Trouvez le mot de la famille de \eng{to lend} qui complète : \eng{Money \ldots\ldots\ldots charge high interest rates.}
\item Vrai ou faux : \eng{to save a child} signifie « économiser un enfant ».
\item Complétez avec la famille de \eng{to invest} : \eng{Foreign \ldots\ldots\ldots made a huge \ldots\ldots\ldots in the port of Kribi.}
\end{enumerate}
\end{exercice}

\begin{corrige}[Exercice « À vous de jouer »]
\begin{enumerate}
\item \eng{My account is in the red, so I must save more (money).} — \eng{in the red} pour « à découvert », \eng{to save} pour « économiser ».
\item \eng{The bank lent me 50,000 francs.} — la banque prête, c'est moi qui emprunte ; attention au prétérit irrégulier \eng{lend → lent → lent}.
\item \eng{lenders} — \eng{money lenders} = « les prêteurs d'argent » (personne, pluriel, sujet de \eng{charge}).
\item Faux : \eng{to save a child} signifie « sauver un enfant » ; « économiser » ne s'applique qu'à l'argent, au temps ou à l'énergie (\eng{to save time/energy}).
\item \eng{investors} (personnes, pluriel) et \eng{investment} (nom de l'action) : « Les investisseurs étrangers ont fait un énorme investissement dans le port de Kribi. »
\end{enumerate}
\end{corrige}

\begin{aretenir}[L'essentiel de la section]
\begin{itemize}
\item \eng{To save money} = économiser ; \eng{to save/rescue somebody} = sauver quelqu'un.
\item \eng{A bank note} = un billet ; \eng{a bank} = une banque ; \eng{a bench} = un banc.
\item Jamais \eng{to make a credit} → \eng{to take out a loan} ; jamais \eng{to touch money} → \eng{to be paid / to receive money} ; jamais \eng{to retire money} → \eng{to withdraw money}.
\item Suffixes : \emph{-er/-or} (personne), \emph{-ing} (activité), \emph{-ment} (action/résultat), \emph{-able} (possibilité).
\item \eng{To lend} (prêter) ≠ \eng{to borrow} (emprunter) ; \eng{lend → lent → lent}.
\item Expressions : \eng{in the red} (à découvert), \eng{in the black} (créditeur/bénéficiaire), \eng{to break the bank} (coûter une fortune), \eng{safe as the Bank of England} (sans risque).
\end{itemize}
\end{aretenir}

\coursec{Lexique en contexte : dialogues et production guidée}

Dans la section précédente, nous avons appris à repérer les \cle{faux amis}, à construire des \cle{familles de mots} et à manier les expressions idiomatiques du monde bancaire. Mais un vocabulaire n'est vraiment acquis que lorsqu'il est \emph{employé} : à l'oral, dans un dialogue authentique, et à l'écrit, dans un texte cohérent. C'est précisément l'objectif de cette dernière section : passer du \emph{savoir} (connaître les mots) au \emph{savoir-faire} (les utiliser correctement, avec politesse et aisance). Nous allons d'abord observer un dialogue modèle à la banque, puis lire un texte de compréhension type Bac, et enfin produire nos propres dialogues et paragraphes.

\courssub{Dialogue modèle : “At the bank”}

\courssub{Observation : écoutons et lisons}

Imaginez la scène : un jeune client entre dans une agence bancaire à Yaoundé. Il veut ouvrir un compte d'épargne. Lisez le dialogue ci-dessous, puis observez comment le vocabulaire étudié dans cette leçon est réemployé naturellement.

\begin{textebox}[At the bank — a model dialogue]
\textbf{Clerk:} Good morning, how can I help you?

\textbf{Customer:} Good morning. I'd like to open a savings account.

\textbf{Clerk:} Certainly. Do you have your ID card?

\textbf{Customer:} Yes, here it is.

\textbf{Clerk:} How much would you like to deposit today?

\textbf{Customer:} 10,000 CFA francs.

\textbf{Clerk:} Fine. Your savings will earn 3\% interest per year. Here is your account number and your bank card.

\textbf{Customer:} Thank you very much!

\textbf{Clerk:} You're welcome. You can withdraw money at any branch, and you can also check your balance online.

\textbf{Customer:} That's very convenient. Have a nice day!
\end{textebox}

\textbf{Glossaire :}

\begin{center}
\begin{tabularx}{\linewidth}{|l|l|X|}
\hline
\rowcolor{popblueL} \textbf{English} & \textbf{Nature} & \textbf{Français} \\
\hline
clerk & nom & employé(e) de banque, guichetier(-ière) \\
\hline
savings account & nom & compte d'épargne \\
\hline
ID card & nom & carte d'identité \\
\hline
to deposit & verbe & déposer (de l'argent) \\
\hline
interest & nom & intérêt (taux d'intérêt : interest rate) \\
\hline
account number & nom & numéro de compte \\
\hline
bank card & nom & carte bancaire \\
\hline
to withdraw & verbe & retirer (de l'argent) \\
\hline
branch & nom & agence, succursale \\
\hline
balance & nom & solde (du compte) \\
\hline
convenient & adjectif & pratique, commode \\
\hline
\end{tabularx}
\end{center}

\begin{attention}[La politesse à l'oral : un piège culturel]
À l'oral, ne dites jamais \eng{I want to open an account.} « Je veux ouvrir un compte. » : en anglais, \eng{I want} sonne trop direct, presque impoli, comme un ordre. Le français « je veux » est déjà familier, mais « je voudrais » est la norme polie ; en anglais, c'est encore plus marqué. Préférez toujours la forme conditionnelle atténuée :

\eng{I'd like to open an account.} « Je voudrais ouvrir un compte. »

De même, évitez \eng{Give me my money!} « Donne-moi mon argent ! » et dites \eng{I'd like to withdraw some money, please.} La petite formule \eng{please} et le conditionnel \eng{I'd like...} (contraction de \emph{I would like}) sont les deux clés de la politesse bancaire en anglais.
\end{attention}

\courssub{Les formules fonctionnelles à mémoriser}

Un dialogue à la banque repose sur un petit nombre de \cle{formules fonctionnelles} (ou \emph{functional language}) : des phrases toutes faites qui remplissent une fonction précise (demander, proposer, s'informer, remercier). Les voici, classées par fonction :

\begin{center}
\begin{tabularx}{\linewidth}{|l|X|X|}
\hline
\rowcolor{popblueL} \textbf{Fonction} & \textbf{English} & \textbf{Français} \\
\hline
Saluer / offrir de l'aide & Good morning, how can I help you? & Bonjour, comment puis-je vous aider ? \\
\hline
Formuler une demande & I'd like to open an account. & Je voudrais ouvrir un compte. \\
\hline
Demander un document & Do you have your ID card? & Avez-vous votre carte d'identité ? \\
\hline
Poser une question (montant) & How much would you like to deposit? & Combien souhaitez-vous déposer ? \\
\hline
S'informer (taux) & What is the interest rate? & Quel est le taux d'intérêt ? \\
\hline
S'informer (possibilité) & Can I withdraw money at any branch? & Puis-je retirer de l'argent dans n'importe quelle agence ? \\
\hline
Remercier & Thank you very much! / You're welcome. & Merci beaucoup ! / Je vous en prie. \\
\hline
\end{tabularx}
\end{center}

\begin{exemplebox}[Exemples traduits à réemployer]
\eng{I'd like to withdraw some money, please.} « Je voudrais retirer de l'argent, s'il vous plaît. »

\eng{What is the interest rate on a savings account?} « Quel est le taux d'intérêt sur un compte d'épargne ? »

\eng{How much would you like to deposit today?} « Combien souhaitez-vous déposer aujourd'hui ? »

\eng{Can I check my balance at any branch?} « Puis-je consulter mon solde dans n'importe quelle agence ? »

\eng{I'd like to apply for a loan.} « Je voudrais demander un prêt. »
\end{exemplebox}

\begin{propriete}[Structure de la demande polie]
Forme : \textbf{I'd like to + base verbale} (I'd like = I would like).

\eng{I'd like to open a current account.} « Je voudrais ouvrir un compte courant. »

Pour la question, on inverse l'auxiliaire : \eng{Would you like to deposit some money today?} « Souhaitez-vous déposer de l'argent aujourd'hui ? »

Attention à la question avec \emph{how much} : \eng{How much would you like to deposit?} — \emph{how much} (et non \emph{how many}) car l'argent (\emph{money}) est indénombrable en anglais.
\end{propriete}

\courssub{La structure du dialogue : un organigramme}

Tout dialogue bancaire suit le même parcours en cinq étapes. Le schéma ci-dessous le résume ; apprenez-le, il vous servira de plan pour vos propres productions orales.

\begin{popfigure}
\begin{center}
\begin{tikzpicture}[
  node distance=0.55cm,
  etape/.style={rectangle, rounded corners, draw=popblue, fill=popblueL, thick, align=center, text width=4.6cm, inner sep=5pt},
  fleche/.style={-{Stealth[length=3mm]}, thick, popdark}
]
\node[etape, align=center] (e1){\textbf{1. Greeting}\\ \eng{Good morning!}};
\node[etape, below=of e1, align=center] (e2){\textbf{2. Request}\\ \eng{I'd like to open an account.}};
\node[etape, below=of e2, align=center] (e3){\textbf{3. Questions}\\ \eng{ID card? How much deposit?}};
\node[etape, below=of e3, align=center] (e4){\textbf{4. Explanation of services}\\ \eng{interest rate, bank card, branches}};
\node[etape, below=of e4, align=center] (e5){\textbf{5. Thanks}\\ \eng{Thank you very much!}};
\draw[fleche] (e1) -- (e2);
\draw[fleche] (e2) -- (e3);
\draw[fleche] (e3) -- (e4);
\draw[fleche] (e4) -- (e5);
\end{tikzpicture}
\end{center}
\legende{Organigramme du dialogue type à la banque : greeting $\rightarrow$ request $\rightarrow$ questions $\rightarrow$ explanation of services $\rightarrow$ thanks.}
\end{popfigure}

\begin{methode}[Réussir le dialogue à la banque (jeu de rôle)]
Pour jouer ou rédiger un dialogue bancaire réussi, suivez ces quatre étapes :
\begin{enumerate}
\item \textbf{Saluez} : \eng{Good morning / Good afternoon.} Le clerk ajoute \eng{How can I help you?}
\item \textbf{Formulez votre demande} avec \eng{I'd like to...} : \eng{I'd like to open a savings account.}
\item \textbf{Répondez aux questions du clerk} (identité, montant du dépôt) et posez vos propres questions : \eng{What is the interest rate? Can I withdraw money at any branch?}
\item \textbf{Remerciez} : \eng{Thank you very much!}
\end{enumerate}
Consigne de qualité : utilisez \textbf{au moins 5 mots du lexique} de la leçon (account, deposit, withdraw, interest rate, balance, bank card, branch, loan...) et au moins une \cle{collocation} (\eng{open an account}, \eng{earn interest}, \eng{check your balance}).
\end{methode}

\begin{experience}[Jeu de rôle : At the bank]
Travaillez par deux. L'élève A est le \eng{clerk}, l'élève B est le \eng{customer}. Le customer choisit une mission : ouvrir un compte d'épargne, retirer de l'argent, ou demander un prêt (\eng{loan}). Le dialogue doit comporter au moins 8 répliques, suivre l'organigramme en 5 étapes et contenir au moins 5 mots du lexique bancaire. Puis échangez les rôles. La classe écoute et coche les mots du lexique entendus.
\end{experience}

\courssub{Compréhension de l'écrit : “Mobile Money vs Banks in Cameroon”}

Voici un court texte de compréhension, dans l'esprit des épreuves du Baccalauréat. Lisez-le attentivement, puis répondez aux questions.

\begin{textebox}[Mobile Money vs Banks in Cameroon]
In Cameroon today, there are two main ways to keep and move money: traditional banks and mobile money services. Banks offer savings accounts, current accounts, loans and bank cards. To open an account, a customer needs an ID card and a first deposit. The money in a savings account earns interest, usually between 2\% and 4\% per year. Customers can withdraw cash at any branch or at an ATM, and they can check their balance online.

Mobile money, on the other hand, works with a simple telephone. There is no need to visit a branch: people can deposit, transfer and withdraw money at small agents' shops in every neighbourhood, even in villages where there is no bank. This is why mobile money has become extremely popular, especially among young people, market traders and motorcycle taxi drivers, the famous \eng{bendskin} riders of Douala and Yaoundé.

However, mobile money has limits. Transfers can be expensive, and the money kept on a telephone does not earn any interest. Moreover, if you lose your telephone, you may have problems accessing your money. Banks, by contrast, are safer for large amounts and offer services that mobile money does not, such as loans for students or business projects.

Many Cameroonians now use both systems. They keep their savings in a bank account, where the money earns interest, and they use mobile money for small daily payments. For young people, financial experts often give one piece of advice: open a savings account early, deposit a little money every month, and watch your savings grow. As the saying goes, “a penny saved is a penny earned.”
\end{textebox}

\textbf{Glossaire :}

\begin{center}
\begin{tabularx}{\linewidth}{|l|l|X|}
\hline
\rowcolor{popblueL} \textbf{English} & \textbf{Nature} & \textbf{Français} \\
\hline
mobile money & nom & argent mobile (transfert par téléphone) \\
\hline
current account & nom & compte courant \\
\hline
ATM & nom & distributeur automatique de billets \\
\hline
transfer & nom & transfert (d'argent) \\
\hline
agent & nom & agent, point de vente agréé \\
\hline
trader & nom & commerçant(e) \\
\hline
safer & adjectif comparatif & plus sûr (safe $\rightarrow$ safer) \\
\hline
a penny saved is a penny earned & proverbe & un sou économisé est un sou gagné \\
\hline
\end{tabularx}
\end{center}

\begin{exoresolu}[Questions de compréhension type Bac]
\textbf{A. Vrai ou faux ? Justifiez avec une phrase du texte.}

1. You need an ID card to open a bank account.

2. Money kept on a telephone earns interest.

3. Some Cameroonians use both banks and mobile money.

\textbf{B. QCM.} Mobile money is very popular because :

a) it pays high interest \quad b) it works everywhere, even in villages \quad c) it gives loans to students.

\textbf{C. Réponse courte.} What advice do financial experts give to young people?

\tcblower

\textbf{A.} 1. \textbf{True} — \eng{To open an account, a customer needs an ID card and a first deposit.}

2. \textbf{False} — \eng{The money kept on a telephone does not earn any interest.}

3. \textbf{True} — \eng{Many Cameroonians now use both systems.}

\textbf{B.} Réponse \textbf{b} : \eng{People can deposit, transfer and withdraw money at small agents' shops in every neighbourhood, even in villages where there is no bank.}

\textbf{C.} \eng{They advise young people to open a savings account early, deposit a little money every month, and watch their savings grow.} (« Ils conseillent aux jeunes d'ouvrir tôt un compte d'épargne, d'y déposer un peu d'argent chaque mois et de regarder leur épargne grandir. »)
\end{exoresolu}

\courssub{Production guidée : à vous d'écrire}

\begin{exercice}[Tâche 1 — Un dialogue à la banque]
Rédigez un dialogue de \textbf{8 à 10 répliques} entre un clerk et un customer. Situation au choix : ouvrir un compte d'épargne, retirer de l'argent à un guichet, ou demander des informations sur un prêt (\eng{loan}). Votre dialogue doit suivre l'organigramme en 5 étapes et contenir au moins 8 mots du lexique bancaire.
\end{exercice}

\begin{exercice}[Tâche 2 — Un paragraphe argumentatif]
Rédigez un paragraphe de 8 à 10 lignes intitulé \eng{Why young people should open a savings account}. Utilisez le plan suivant : 1) introduire le sujet ; 2) donner deux ou trois arguments (safety, interest, future projects) ; 3) conclure avec un conseil. Employez au moins 8 mots du lexique et une expression idiomatique.
\end{exercice}

\begin{corrige}[Tâche 2 — modèle de paragraphe]
\eng{Young people should open a savings account as early as possible. First, a bank is a safe place to keep money: unlike cash kept at home, savings in a bank cannot easily be lost or stolen. Second, the money in a savings account earns interest, so it grows little by little every year. Third, having an account teaches young people to manage their money and to prepare future projects, such as university studies or a small business. They can deposit a small amount every month and check their balance regularly. As the saying goes, “a penny saved is a penny earned.” My advice is simple: go to the nearest branch, ask the clerk how to open an account, and start saving today.}

Ce modèle contient plus de 10 mots du lexique (\emph{savings account, bank, savings, interest, deposit, balance, branch, clerk, open an account, cash}), deux collocations (\emph{earn interest}, \emph{check their balance}) et une expression idiomatique (le proverbe final).
\end{corrige}

\begin{methode}[Grille d'auto-évaluation de ma production]
Relisez votre dialogue ou votre paragraphe et cochez chaque critère atteint :
\begin{enumerate}
\item Ai-je utilisé \textbf{au moins 8 mots du lexique} bancaire (account, deposit, withdraw, balance, interest rate, branch, loan, bank card, clerk, savings...) ?
\item Ai-je employé \textbf{au moins 2 collocations} correctes (\eng{open an account}, \eng{earn interest}, \eng{withdraw money}, \eng{check my balance}, \eng{apply for a loan}) ?
\item Ai-je intégré \textbf{au moins 1 expression idiomatique} ou un proverbe (\eng{a penny saved is a penny earned}, \eng{save for a rainy day}) ?
\item Ai-je évité les \textbf{faux amis} (\emph{money} indénombrable, \eng{to withdraw} et non « to retire money ») et les calques du français ?
\item Dans un dialogue : ai-je utilisé la forme polie \eng{I'd like to...} et non \eng{I want...} ?
\end{enumerate}
Si vous cochez les cinq cases, votre production est réussie. Sinon, corrigez le point faible et relisez.
\end{methode}

\begin{attention}[Dernières erreurs à traquer avant de rendre votre copie]
\begin{itemize}
\item \eng{How much money do you have?} et non \eng{How many money...} : \emph{money} est indénombrable.
\item \eng{I went to the bank to withdraw money.} et non « to retire money » : \emph{retirer} de l'argent se dit \eng{to withdraw}.
\item \eng{The interest rate is 3\% per year.} : \emph{per year}, pas « by year ».
\item Ordre des mots dans la question indirecte : \eng{Can you tell me what the interest rate is?} (et non « what is the interest rate » après \emph{can you tell me}).
\end{itemize}
\end{attention}

\begin{aretenir}[L'essentiel de la section]
Un dialogue bancaire suit cinq étapes : salutation, demande polie avec \eng{I'd like to...}, questions et réponses, explication des services, remerciements. À l'écrit comme à l'oral, la qualité d'une production se mesure à la réutilisation du lexique : visez au moins 8 mots du champ lexical, 2 collocations et 1 expression idiomatique. Enfin, relisez toujours votre texte à la chasse des faux amis et des calques du français : c'est la signature d'un excellent élève d'anglais.
\end{aretenir}

\newpage

\coursec{Activité d'intégration}

\coursec{Lesson 18 — Vocabulary: Words and expressions related to banking services}

\courssub{Objectifs de l'activité}

À la fin de cette activité d'intégration, tu seras capable de :

\begin{itemize}
\item \cle{réemployer} à l'oral, dans une situation de communication authentique, le vocabulaire des services bancaires étudié dans la leçon (\eng{open an account, withdraw money, deposit, interest rate, bank charges, loan, overdraft}\ldots) ;
\item \cle{conduire} un échange formel client--employé de banque en utilisant les formules de politesse appropriées (\eng{Could you\ldots? I would like to\ldots; May I help you?}) ;
\item \cle{comprendre} et \cle{produire} des questions et des réponses sur les documents, les frais et les conditions d'un compte bancaire ;
\item \cle{évaluer} la production de tes camarades à l'aide d'une grille d'observation (au moins 8 mots du lexique, 2 collocations, 1 formule polie) ;
\item \cle{transposer} l'oral à l'écrit en rédigeant un court paragraphe de 10 lignes intitulé \eng{My visit to the bank}.
\end{itemize}

\courssub{Situation}

\begin{textebox}[Notice at the entrance of Buea Commercial Bank (BCB)]
\textbf{Buea Commercial Bank — Molyko Branch}

\textbf{Welcome to BCB! Your money, our commitment.}

\textbf{Our services:}
\begin{itemize}
\item Open a savings account or a current account in only 15 minutes. Minimum deposit: 5,000 FCFA.
\item Withdraw and deposit money at the counter or at our cash machines (ATMs), 24 hours a day.
\item Apply for a small business loan at a low interest rate. Bring your ID card, a passport photograph and proof of income.
\item Get your BCB bank card and pay in shops without cash.
\item Send and receive money transfers anywhere in Cameroon and abroad.
\end{itemize}

\textbf{Please note:} account holders pay no monthly fees during the first year. Overdrafts are allowed only on current accounts. Our customer service desk is open from 8 a.m. to 4 p.m., Monday to Friday.

\emph{BCB — Banking made simple for the people of the South West Region.}
\end{textebox}

\begin{definition}[Glossaire du document]
\begin{itemize}
\item \eng{current account} (n.) : compte courant ; \eng{savings account} (n.) : compte d'épargne ;
\item \eng{minimum deposit} (n.) : dépôt minimum ; \eng{cash machine / ATM} (n.) : distributeur automatique de billets ;
\item \eng{to apply for a loan} (v.) : demander un prêt ; \eng{interest rate} (n.) : taux d'intérêt ;
\item \eng{proof of income} (n.) : justificatif de revenus ; \eng{fee} (n.) : frais (bancaires) ;
\item \eng{overdraft} (n.) : découvert (autorisé) ; \eng{account holder} (n.) : titulaire d'un compte.
\end{itemize}
\end{definition}

\textbf{Mise en situation.} Tu es élève en Terminale A au lycée de Buea. Avec deux camarades, tu visites l'agence \eng{Buea Commercial Bank} de Molyko. L'un de vous est un jeune commerçant (ou une jeune commerçante) qui veut ouvrir un compte, retirer de l'argent ou demander un petit prêt pour développer son commerce (vente de beignets, boutique de \emph{call-box}, élevage de poulets). Le deuxième joue l'employé(e) de banque qui accueille le client, pose des questions et explique les services. Le troisième est l'observateur : il écoute la scène et coche sur une grille les mots du lexique bancaire effectivement utilisés.

\courssub{Consignes}

\begin{enumerate}
\item \textbf{Préparation du lexique (10 min).} En silence, relis l'affiche de la banque et souligne tous les mots et expressions liés aux services bancaires. Classe-les en trois colonnes : \emph{people and places}, \emph{accounts and operations}, \emph{payment, credit and savings}. Compare ta liste avec celle de ton voisin.

\item \textbf{Compréhension du document (5 min).} Réponds oralement aux questions suivantes en phrases complètes :
\begin{enumerate}
\item How much money do you need to open an account at BCB?
\item What documents must you bring to apply for a loan?
\item Which type of account allows an overdraft?
\item When does a new account holder stop paying monthly fees?
\end{enumerate}

\item \textbf{Préparation du jeu de rôle (10 min).} Par groupes de trois, choisissez les rôles : \emph{customer}, \emph{bank clerk}, \emph{observer}. Choisissez une mission pour le client : (a) open a savings account ; (b) withdraw money and ask about bank charges ; (c) apply for a small business loan. Préparez ensemble un dialogue d'environ 12 répliques. Le client doit poser au moins trois questions ; l'employé doit expliquer au moins deux services de l'affiche.

\item \textbf{Performance et observation (10 min par scène).} Jouez la scène. L'observateur remplit la grille ci-dessous et vérifie que le dialogue contient : au moins \cle{8 mots} du champ lexical bancaire, \cle{2 collocations} (ex. \eng{open an account}, \eng{pay interest}) et \cle{1 formule de politesse} (ex. \eng{Could you tell me\ldots?}).

\begin{center}
\begin{tabularx}{\linewidth}{|X|X|X|}
\hline
\rowcolor{popblueL} \textbf{Critère} & \textbf{Exemples attendus} & \textbf{Observé ? (oui/non)} \\
\hline
8 mots du lexique bancaire & account, deposit, withdraw, loan, interest rate, fees, bank card, transfer & \\
\hline
2 collocations correctes & open an account, apply for a loan, pay fees & \\
\hline
1 formule de politesse & May I help you? / I would like to\ldots & \\
\hline
Prononciation claire & « oui-zdro », « inté-rest », « loune » & \\
\hline
\end{tabularx}
\end{center}

\item \textbf{Rotation des rôles (10 min).} Échangez les rôles : l'observateur devient client, le client devient employé, l'employé devient observateur. Choisissez une mission différente et rejouez la scène.

\item \textbf{Représentation devant la classe (10 min).} Les deux meilleurs groupes jouent leur scène devant toute la classe. Les autres élèves notent sur leur cahier les mots du lexique qu'ils reconnaissent.

\item \textbf{Production écrite individuelle (15 min).} Rédige sur ton cahier un paragraphe de 10 lignes intitulé \eng{My visit to the bank}. Raconte ta visite à la banque au \cle{prétérit} : ce que tu voulais faire, les questions posées, les explications de l'employé, le résultat. Utilise au moins 6 mots du lexique bancaire et 2 collocations.
\end{enumerate}

\courssub{Ressources à mobiliser}

\begin{aretenir}[Boîte à outils linguistique]
\textbf{Lexique utile :}
\begin{itemize}
\item Personnes et lieux : \eng{bank clerk, customer, account holder, branch, counter, cash machine (ATM)}.
\item Comptes et opérations : \eng{savings account, current account, to open/close an account, to deposit money, to withdraw money, balance, statement, transfer}.
\item Paiement, crédit et épargne : \eng{bank card, cash, cheque, loan, interest rate, fees/charges, overdraft, to save money}.
\end{itemize}

\textbf{Collocations à réemployer :} \eng{open an account} (ouvrir un compte), \eng{apply for a loan} (demander un prêt), \eng{pay interest} (payer des intérêts), \eng{withdraw cash} (retirer de l'argent liquide), \eng{pay monthly fees} (payer des frais mensuels).

\textbf{Formules de politesse (registre formel) :}
\begin{itemize}
\item L'employé : \eng{Good morning, may I help you? / What can I do for you today? / Let me explain our services.}
\item Le client : \eng{Good morning, I would like to open an account. / Could you tell me what the interest rate is? / How much are the monthly fees? / Thank you very much for your help.}
\end{itemize}

\textbf{Grammaire :} questions avec inversion (\eng{How much do I need?}), \eng{would like to} + base verbale pour exprimer un souhait poli, prétérit simple pour le récit écrit (\eng{I went, I asked, the clerk explained, I opened}).
\end{aretenir}

\begin{attention}[Pièges à éviter]
\begin{itemize}
\item Ne dis pas \eng{*I want to open an account} (trop direct) : préfère \eng{I would like to open an account}.
\item \eng{money} est indénombrable : \eng{I withdrew some money}, jamais \eng{*moneys}.
\item Attention au faux ami : \eng{to demand} signifie « exiger » ; pour « demander (un prêt) », dis \eng{to apply for a loan} ou \eng{to ask for a loan}.
\item Prononciation : \eng{withdraw} se dit « oui-zdro », \eng{loan} se dit « loune », \eng{cheque} se prononce exactement comme \eng{check} (« tchèk »).
\item Orthographe : \eng{cheque} (anglais britannique) et non \eng{check} (américain) ; \eng{programme} (britannique) mais \eng{program} en informatique.
\end{itemize}
\end{attention}

\begin{savaistu}[Banque et culture : Cameroun et monde anglophone]
Au Cameroun, le secteur bancaire est bilingue : les guichets affichent souvent \eng{Depot/Retrait} et \eng{Deposit/Withdrawal}. Les termes anglais dominent dans les banques à capital international (Standard Chartered, UBA, Ecobank) et dans la zone anglophone (Buea, Bamenda, Limbe). Au Royaume-Uni, on dit \eng{current account} (compte courant) et \eng{cheque} ; aux États-Unis, \eng{checking account} et \eng{check}. Le \eng{sort code} (UK) équivaut au \eng{routing number} (US) pour identifier la banque. Le \eng{IBAN} est utilisé en Europe, le \eng{SWIFT/BIC} pour les virements internationaux. Au Cameroun, le \eng{mobile money} (MTN MoMo, Orange Money) coexiste avec la banque classique ; on dit \eng{cash in / cash out} pour déposer/retirer chez un agent.
\end{savaistu}

\courssub{Proposition de production et corrigé commenté}

\begin{exoresolu}[Modèle de dialogue et de paragraphe]
\textbf{Énoncé.} Produis (a) un dialogue client--employé d'environ 12 répliques pour la mission « apply for a small business loan », puis (b) le paragraphe écrit \eng{My visit to the bank} (10 lignes).

\tcblower

\textbf{(a) Dialogue modèle.}

\textbf{Clerk:} “Good morning, madam. May I help you?”

\textbf{Customer:} “Good morning. I would like to apply for a small business loan, please.”

\textbf{Clerk:} “Certainly. What kind of business do you run?”

\textbf{Customer:} “I sell beignets and puff-puff near the Molyko market, and I want to buy a bigger frying pot.”

\textbf{Clerk:} “I see. Are you an account holder at our bank?”

\textbf{Customer:} “Yes, I opened a savings account here last year.”

\textbf{Clerk:} “Very good. Could you show me your ID card and a proof of income?”

\textbf{Customer:} “Here you are. Could you tell me what the interest rate is?”

\textbf{Clerk:} “For small business loans, the interest rate is very low, and you can repay the loan over twelve months.”

\textbf{Customer:} “And how much are the fees?”

\textbf{Clerk:} “There are no monthly fees during your first year. With your BCB bank card, you can also withdraw cash at our cash machine at any time.”

\textbf{Customer:} “That is wonderful. Thank you very much for your help.”

\textbf{Clerk:} “You are welcome. Please fill in this form, and come back tomorrow with a passport photograph.”

\textbf{(b) Paragraphe modèle — My visit to the bank.}

\eng{Yesterday morning, I went to Buea Commercial Bank in Molyko to apply for a small business loan. At the counter, a friendly bank clerk greeted me and asked how she could help me. I explained that I sell beignets near the market and that I needed money to buy a bigger frying pot. She checked that I was an account holder and asked for my ID card and a proof of income. Then she explained that the interest rate was low and that I could repay the loan over twelve months. She also told me that there were no monthly fees during the first year. Finally, I filled in a form, and she gave me an appointment for the next day. I was very satisfied because the clerk was polite and everything was clear. I hope the bank will accept my application, so that my small business can grow.}

\textbf{Commentaire des choix de langue.}
\begin{itemize}
\item \cle{Politesse} : le client utilise \eng{I would like to\ldots} et \eng{Could you tell me\ldots?} plutôt que l'impératif ou \eng{I want}, ce qui correspond au registre formel d'une banque.
\item \cle{Collocations} réussies : \eng{apply for a loan}, \eng{open a savings account}, \eng{repay the loan}, \eng{withdraw cash}, \eng{fill in a form} — verbe + nom correctement associés, sans calque du français.
\item \cle{Lexique} : plus de 12 mots du champ bancaire sont réemployés (\eng{loan, interest rate, fees, account holder, counter, cash machine, ID card, proof of income}\ldots).
\item \cle{Grammaire} : questions avec inversion (\eng{How much are the fees?}), prétérit simple régulier dans le récit (\eng{greeted, explained, checked, filled}), présent de vérité générale pour les conditions de la banque (\eng{the interest rate is low}).
\item \cle{Faux amis évités} : « demander un prêt » est rendu par \eng{apply for a loan}, et non par \eng{*demand a loan}.
\end{itemize}
\end{exoresolu}

\courssub{Bilan de l'activité}

Grâce à ce jeu de rôle, tu as réinvesti le vocabulaire des services bancaires dans une communication orale authentique et structurée, proche d'une situation réelle de la vie économique camerounaise. Tu as appris à :

\begin{itemize}
\item identifier dans un document authentique (l'affiche de la banque) les mots-clés du champ lexical et les classer par sous-thèmes ;
\item employer à l'oral les collocations essentielles (\eng{open an account, apply for a loan, withdraw cash}) et les formules de politesse du registre formel ;
\item écouter et évaluer la production d'autrui à l'aide d'une grille de critères précis — une compétence utile pour l'épreuve orale du baccalauréat ;
\item transposer un échange oral en récit écrit au prétérit, en veillant à l'orthographe (\eng{cheque, withdraw, receipt}) et à la prononciation des mots difficiles.
\end{itemize}

\begin{methode}[Pour aller plus loin]
\begin{enumerate}
\item Lors de ta prochaine visite réelle dans une banque ou un établissement de microfinance (MC\textsuperscript{2}, coopérative de crédit), note trois mots anglais que tu entends ou lis sur les affiches.
\item Rédige un deuxième paragraphe au présent, intitulé \eng{The services of my bank}, en décrivant les services d'une banque que tu connais.
\item En binôme, invente une nouvelle scène : un client mécontent se plaint de frais bancaires trop élevés (\eng{to complain about high bank charges}).
\end{enumerate}
\end{methode}

Cette activité clôt la leçon : le lexique bancaire n'est plus une liste à mémoriser, mais un outil vivant que tu sais mobiliser pour comprendre, dialoguer et écrire dans le monde économique anglophone.

\newpage

\coursec{Méthodes et exercices résolus}

\coursec{Lesson 18 — Vocabulary: Words and expressions related to banking services}

\courssub{Méthodes et exercices résolus}

\begin{methode}[Construire et structurer son lexique bancaire par champs sémantiques]
    \textbf{Objectif :} Identifier, classer et mémoriser le vocabulaire technique de la banque en trois grands domaines : \emph{People \& Places}, \emph{Accounts \& Operations}, \emph{Payment \& Credit}.
    \begin{enumerate}
        \item \textbf{Collecter} les mots-clés à partir d'un texte support (article, dialogue, publicité bancaire) ou d'une écoute.
        \item \textbf{Classer} chaque mot dans un tableau à trois colonnes selon sa fonction : \textit{Acteurs/Lieux}, \textit{Produits/Opérations}, \textit{Moyens de paiement/Crédit/Épargne}.
        \item \textbf{Associer} les \cle{collocations} (associations figées) : \eng{open/close an account}, \eng{apply for a loan}, \eng{make a deposit/withdrawal}, \eng{charge a fee}, \eng{earn interest}.
        \item \textbf{Visualiser} avec une carte mentale (mind map) centrée sur \eng{Banking Services} pour relier les sous-thèmes.
        \item \textbf{Réviser} par lots de 7 à 10 mots (technique de la répétition espacée) en produisant une phrase personnelle par mot.
    \end{enumerate}
    \textbf{Structures à retenir :} \eng{The \textbf{teller} handles \textbf{cash} deposits.} (Le guichetier gère les dépôts d'espèces.) \quad \eng{You need a \textbf{valid ID} to \textbf{open a current account}.} (Il vous faut une pièce d'identité valide pour ouvrir un compte courant.)
\end{methode}

\begin{exoresolu}[Classification lexicale et collocations — Type Bac : Vocabulary in Context]
    \textbf{Énoncé :} Voici une liste de mots et expressions liés aux services bancaires. 
    \begin{center}
        \texttt{overdraft \quad teller \quad standing order \quad branch \quad mortgage \quad withdrawal \quad PIN \quad interest rate \quad loan officer \quad cheque book \quad ATM \quad savings account \quad direct debit \quad credit limit \quad deposit slip}
    \end{center}
    \textbf{A.} Complétez le tableau ci-dessous en rangeant chaque terme dans la bonne catégorie. (Un terme par case, certains termes peuvent convenir à deux catégories mais placez-les au plus pertinent).
    \begin{center}
        \begin{tabularx}{\linewidth}{|X|X|X|}\hline 
        \rowcolor{popblueL} \textbf{People \& Places} & \textbf{Accounts \& Operations} & \textbf{Payment, Credit \& Savings} \\ \hline
        1. & 1. & 1. \\ \hline
        2. & 2. & 2. \\ \hline
        3. & 3. & 3. \\ \hline
        4. & 4. & 4. \\ \hline
        5. & 5. & 5. \\ \hline
    \end{tabularx}
    \end{center}
    \textbf{B.} Choisissez le verbe correct pour former la collocation naturelle (barrez l'intrus).
    \begin{enumerate}
        \item To \textbf{make / do / take} a deposit.
        \item To \textbf{apply for / ask / demand} a loan.
        \item To \textbf{pay / withdraw / take out} money from an ATM.
        \item To \textbf{charge / load / put} a fee / interest.
        \item To \textbf{set up / establish / create} a standing order.
    \end{enumerate}
    \textbf{C.} Traduisez en anglais : « Le directeur d'agence a refusé ma demande de découvert autorisé car mon taux d'endettement est trop élevé. »
    
    \tcblower
    \textbf{Solution détaillée :}
    \textbf{A. Classification :}
    \begin{center}
        \begin{tabularx}{\linewidth}{|X|X|X|}\hline 
        \rowcolor{popblueL} \textbf{People \& Places} & \textbf{Accounts \& Operations} & \textbf{Payment, Credit \& Savings} \\ \hline
        \textbf{teller} (guichetier) & \textbf{withdrawal} (retrait) & \textbf{overdraft} (découvert) \\ \hline
        \textbf{branch} (agence/succursale) & \textbf{deposit slip} (bordereau de dépôt) & \textbf{mortgage} (prêt hypothécaire) \\ \hline
        \textbf{loan officer} (chargé de clientèle/prêts) & \textbf{standing order} (virement permanent) & \textbf{interest rate} (taux d'intérêt) \\ \hline
        \textbf{ATM} (Distributeur automatique / DAB) & \textbf{savings account} (compte d'épargne) & \textbf{credit limit} (plafond de crédit) \\ \hline
        \textbf{PIN} (Code secret) & \textbf{cheque book} (chéquier) & \textbf{direct debit} (prélèvement automatique) \\ \hline
    \end{tabularx}
    \end{center}
    \textit{Justification :} \eng{Teller} (US) = \eng{Cashier} (UK). \eng{Branch} = lieu physique. \eng{Loan officer} = personne. \eng{ATM} = lieu/machine (\eng{Automated Teller Machine}). \eng{PIN} = outil d'identification (\eng{Personal Identification Number}). \eng{Standing order} = ordre de virement récurrent émis par le client (différent de \eng{Direct debit} émis par le bénéficiaire).
    
    \textbf{B. Collocations (Verbes corrects en gras) :}
    \begin{enumerate}
        \item To \textbf{make} a deposit. (\eng{Do/Make} : \eng{make} pour les opérations financières : \eng{make a payment, make a profit, make a withdrawal}).
        \item To \textbf{apply for} a loan. (\eng{Apply for} = postuler/demander formellement. \eng{Ask for} trop familier, \eng{Demand} trop agressif).
        \item To \textbf{withdraw} / \textbf{take out} money. (\eng{Withdraw} = terme technique formel ; \eng{Take out} = phrasal verb courant. \eng{Pay money} = déposer/payer).
        \item To \textbf{charge} a fee / interest. (\eng{Charge} = facturer/débiter des frais. La banque \eng{charges} le client).
        \item To \textbf{set up} a standing order. (\eng{Set up} = mettre en place/configurer. \eng{Establish} possible mais plus formel/institutionnel).
    \end{enumerate}
    
    \textbf{C. Traduction :}
    \eng{The branch manager turned down my overdraft request because my debt-to-income ratio is too high.}
    \textit{ou} \eng{The bank manager refused my application for an overdraft as my debt ratio is too high.}
    \begin{itemize}
        \item \textbf{Branch manager} (Directeur d'agence) — \eng{Bank manager} aussi correct.
        \item \textbf{Turn down} (refuser/rejeter une demande) — registre soutenu/administratif. \eng{Refuse} + \eng{application/request} correct aussi.
        \item \textbf{Overdraft} (découvert autorisé) — \textbf{Faux ami :} \eng{Overdraft} \neq \eng{Overdraft fee} (frais de découvert). \eng{Unauthorized overdraft} = découvert non autorisé.
        \item \textbf{Debt-to-income ratio} / \textbf{Debt ratio} (taux d'endettement).
    \end{itemize}
    \textbf{Conclusion :} La maîtrise des collocations (\eng{make a deposit}, \eng{apply for}, \eng{charge a fee}) et des termes techniques précis (\eng{overdraft}, \eng{standing order} vs \eng{direct debit}) est cruciale pour le Bac.
\end{exoresolu}

\begin{methode}[Démasquer les faux amis et maîtriser les familles de mots (Word Formation)]
    \textbf{Objectif :} Éviter les calques lexicaux français/anglais et construire le mot grammaticalement adapté (nom, verbe, adjectif, adverbe) grâce aux suffixes productifs.
    \begin{enumerate}
        \item \textbf{Lister} les \cle{faux amis} bancaires fréquents : \eng{Benefit} (avantage/bénéfice \neq \emph{bénéfice} comptable = \eng{Profit}), \eng{Agenda} (ordre du jour \neq \emph{agenda} = \eng{Diary/Planner}), \eng{Commission} (commission \neq \emph{commission} parlementaire = \eng{Committee}), \eng{Receipt} (reçu \neq \emph{recette} = \eng{Revenue/Recipe}), \eng{Signature} (signature \neq \emph{signature} de traité = \eng{Signing/Ratification}), \eng{Credit} (crédit/avoir \neq \emph{crédit} seulement = aussi \eng{reputation/merit}).
        \item \textbf{Construire} le tableau de dérivation (Word Formation) pour les racines clés : \eng{Bank, Finance, Invest, Borrow, Lend, Save, Pay, Transact}.
        \item \textbf{Repérer} les suffixes nominaux : \cle{-tion/-sion} (\eng{transaction}), \cle{-ment} (\eng{payment}), \cle{-ance/-ence} (\eng{finance $\to$ finance} mais \eng{insure $\to$ insurance}), \cle{-al} (\eng{refuse $\to$ refusal} — attention irrégulier), \cle{-ity} (\eng{secure $\to$ security}).
        \item \textbf{Noter} l'accent tonique (stress shift) : \eng{In\textbf{vest}} (v) $\to$ \eng{\textbf{In}vestment} (n) ; \eng{Fi\textbf{nance}} (v/n) $\to$ \eng{Fi\textbf{nan}cial} (adj) ; \eng{Trans\textbf{fer}} (v/n) $\to$ \eng{\textbf{Trans}ferable} (adj).
        \item \textbf{S'exercer} à transformer des phrases : \eng{The bank \textbf{approved} the loan.} $\to$ \eng{The bank gave its \textbf{approval} for the loan.}
    \end{enumerate}
    \textbf{Règle d'or :} En anglais, la catégorie grammaticale (part of speech) change souvent la forme du mot (suffixe) et parfois l'accent. En français, le radical reste souvent stable (\emph{banque, banquier, bancaire}).
\end{methode}

\begin{exoresolu}[Faux amis, dérivation et correction d'erreurs — Type Bac : Grammar/Vocabulary]
    \textbf{Énoncé :}
    \textbf{A.} Les phrases suivantes contiennent chacune une erreur lexicale (faux ami ou mauvaise forme dérivée). Soulignez l'erreur, expliquez-la en français et réécrivez la phrase correcte.
    \begin{enumerate}
        \item The customer asked for a \textbf{benefit} of 500,000 FCFA to start his business.
        \item The \textbf{borrower} signed the \textbf{agreement} without reading the \textbf{clauses}.
        \item The bank \textbf{refused} to \textbf{agree} the loan because the project was risky.
        \item You must \textbf{depose} your money in the \textbf{safe} before 3 PM.
        \item The \textbf{interests} rates on \textbf{economies} accounts are low this year.
    \end{enumerate}
    \textbf{B.} Complétez le tableau de dérivation suivant. Attention aux changements d'orthographe et d'accent.
    \begin{center}
        \begin{tabular}{|l|l|l|l|l|}\hline 
        \rowcolor{popblueL} \textbf{Verbe} & \textbf{Nom (action/résultat)} & \textbf{Nom (agent/personne)} & \textbf{Adjectif} & \textbf{Adverbe} \\ \hline
        \textbf{Invest} & 1. \_\_\_\_\_\_\_\_\_\_\_ & 2. \_\_\_\_\_\_\_\_\_\_\_ & 3. \_\_\_\_\_\_\_\_\_\_\_ & 4. \_\_\_\_\_\_\_\_\_\_\_ \\ \hline
        \textbf{Borrow} & 5. \_\_\_\_\_\_\_\_\_\_\_ & 6. \_\_\_\_\_\_\_\_\_\_\_ & 7. \_\_\_\_\_\_\_\_\_\_\_ (rare) & -- \\ \hline
        \textbf{Save} & 8. \_\_\_\_\_\_\_\_\_\_\_ & 9. \_\_\_\_\_\_\_\_\_\_\_ & 10. \_\_\_\_\_\_\_\_\_\_\_ & -- \\ \hline
        \textbf{Transact} & 11. \_\_\_\_\_\_\_\_\_\_\_ & -- & 12. \_\_\_\_\_\_\_\_\_\_\_ & 13. \_\_\_\_\_\_\_\_\_\_\_ \\ \hline
        \textbf{Secure} & 14. \_\_\_\_\_\_\_\_\_\_\_ & -- & 15. \_\_\_\_\_\_\_\_\_\_\_ & 16. \_\_\_\_\_\_\_\_\_\_\_ \\ \hline
    \end{tabular}
    \end{center}
    \textbf{C.} Traduisez en évitant les pièges : « Le conseiller financier m'a conseillé de diversifier mes placements pour réduire les risques. »
    
    \tcblower
    \textbf{Solution détaillée :}
    \textbf{A. Correction des erreurs :}
    \begin{enumerate}
        \item \textbf{Erreur :} \eng{Benefit}. \textbf{Explication :} \eng{Benefit} = avantage, prestation (ex: \eng{fringe benefits} = avantages sociaux). Pour « somme prêtée / crédit », on utilise \eng{Loan} ou \eng{Credit facility}. \textbf{Correction :} The customer asked for a \textbf{loan} of 500,000 FCFA...
        \item \textbf{Erreur :} \textbf{Aucune}. \eng{Borrower} (emprunteur), \eng{Agreement} (accord/contrat), \eng{Clauses} (clauses) sont corrects. \textit{(Piège : vérifier si l'élève ne confond pas \eng{Borrower} et \eng{Lender} / \eng{Creditor}).}
        \item \textbf{Erreur :} \eng{Agree} (verbe). \textbf{Explication :} \eng{Agree a loan} n'existe pas. La banque \eng{approves} (approuve), \eng{grants} (accorde) ou \eng{authorizes} un prêt. \textbf{Correction :} The bank refused to \textbf{approve / grant} the loan...
        \item \textbf{Erreur :} \eng{Depose} / \eng{Safe}. \textbf{Explication :} \eng{Depose} = déposer (un corps, une arme, un témoignage), pas de l'argent. Verbe bancaire : \eng{Deposit}. \eng{Safe} (adj = sûr / n = coffre-fort). Dans une banque, on dépose à la caisse (\eng{counter/teller}) ou sur son \eng{account}. \textbf{Correction :} You must \textbf{deposit} your money \textbf{into your account} / \textbf{at the counter} before 3 PM.
        \item \textbf{Erreur :} \eng{Interests} (pluriel) / \eng{Economies}. \textbf{Explication :} \eng{Interest} (taux d'intérêt) est \textbf{incomptable} (uncountable) $\to$ jamais de \textit{s}. \eng{Economies} = \emph{économies} (savings) ou \emph{économie} (science). Compte d'épargne = \eng{Savings account}. \textbf{Correction :} The \textbf{interest} rates on \textbf{savings} accounts are low this year.
    \end{enumerate}
    
    \textbf{B. Tableau de dérivation (Word Formation) :}
    \begin{center}
        \begin{tabular}{|l|l|l|l|l|}\hline 
        \rowcolor{popblueL} \textbf{Verbe} & \textbf{Nom (action/résultat)} & \textbf{Nom (agent)} & \textbf{Adjectif} & \textbf{Adverbe} \\ \hline
        \textbf{Invest} & \textbf{Investment} & \textbf{Investor} & \textbf{Investment} (adj attrib.) / \textbf{Investable} & -- \\ \hline
        \textbf{Borrow} & \textbf{Borrowing} & \textbf{Borrower} & \textbf{Borrowed} (participe) & -- \\ \hline
        \textbf{Save} & \textbf{Savings} (pluriel usuel) / \textbf{Saving} & \textbf{Saver} & \textbf{Saving} (ex: \eng{saving grace}) & -- \\ \hline
        \textbf{Transact} & \textbf{Transaction} & -- & \textbf{Transactional} & \textbf{Transactionally} \\ \hline
        \textbf{Secure} & \textbf{Security} & -- & \textbf{Secure} / \textbf{Secured} & \textbf{Securely} \\ \hline
    \end{tabular}
    \end{center}
    \textit{Notes phonologiques :} \eng{In\textbf{vest}} $\to$ \eng{\textbf{In}vestment} (accent sur 1ère syllabe). \eng{Fi\textbf{nance}} $\to$ \eng{Fi\textbf{nan}cial}. \eng{Se\textbf{cure}} $\to$ \eng{Se\textbf{cu}rity} (son /ʒ/), \eng{Se\textbf{cure}ly}.
    
    \textbf{C. Traduction :}
    \eng{The financial advisor advised me to diversify my investments to reduce the risks.}
    \textit{ou} \eng{The financial adviser recommended (that) I diversify my portfolio to minimise risk.}
    \begin{itemize}
        \item \textbf{Financial advisor / adviser} (conseiller financier — \eng{Adviser} plus courant UK, \eng{Advisor} US).
        \item \textbf{Advise} (verbe) vs \textbf{Advice} (nom, incomptable) — \textbf{Faux ami orthographique}.
        \item \textbf{Investments} (placements — sens financier). \eng{Placements} n'existe pas en anglais financier.
        \item \textbf{Reduce / Minimise / Mitigate} risks.
    \end{itemize}
\end{exoresolu}

\begin{methode}[Comprendre un dialogue bancaire authentique (Listening / Reading Comprehension)]
    \textbf{Objectif :} Extraire l'information explicite et implicite d'un échange oral ou écrit entre un client et un employé de banque (guichet, conseiller, service clientèle).
    \begin{enumerate}
        \item \textbf{Anticiper} : Lire les questions \textbf{avant} d'écouter/lire. Identifier les mots-clés attendus (montants, noms, dates, numéros, problèmes).
        \item \textbf{Écouter/Lire} une 1ère fois pour la \cle{compréhension globale} : Qui ? Où ? Quel service ? Quel problème ? Issue résolue ?
        \item \textbf{Écouter/Lire} une 2ème fois pour les \cle{détails} : Chiffres (account number, amount, interest rate), noms propres (orthographe), termes techniques (\eng{overdraft limit, standing order, direct debit mandate}).
        \item \textbf{Noter} les \cle{marqueurs de relation} : \eng{However, unfortunately, alternatively, let me check, I'll put you through} (indiquent un changement, un refus, une solution).
        \item \textbf{Répondre} en phrases complètes si exigé, ou en mots-clés pour le \textit{note-taking}. Respecter l'orthographe des termes techniques entendus.
    \end{enumerate}
    \textbf{Réflexes :} \eng{Could you spell that, please?} $\to$ Épeler. \eng{Let me repeat the details.} $\to$ Vérification. \eng{To put someone on hold} = mettre en attente. \eng{To sort out} = régler un problème.
\end{methode}

\begin{exoresolu}[Compréhension d'un dialogue : Ouverture de compte et services numériques — Type Bac : Listening/Reading]
    \textbf{Énoncé :} Lisez le dialogue suivant entre Mme \textbf{Alice Mvondo}, cliente à la \textbf{BCEC Bank} (Douala), et M. \textbf{Kouam}, \textbf{Customer Service Officer}. Répondez aux questions en anglais.

    \begin{textebox}[Dialogue at BCEC Bank, Douala]
    \textbf{Mr Kouam:} Good morning, Madam. Welcome to BCEC Bank. How can I help you today?\\
    \textbf{Mrs Mvondo:} Good morning. I would like to open a current account for my small business — a cocoa beans export venture. I also need a debit card and access to online banking.\\
    \textbf{Mr Kouam:} Certainly. We have the \textbf{“Business Starter”} package. It includes a cheque book, a Visa debit card with a daily withdrawal limit of 500,000 FCFA, and free access to our mobile app and internet banking. The monthly maintenance fee is 2,500 FCFA, waived if the average monthly balance stays above 1,000,000 FCFA.\\
    \textbf{Mrs Mvondo:} That sounds suitable. What documents are required?\\
    \textbf{Mr Kouam:} Your National Identity Card, a proof of address (electricity bill or tenancy agreement less than three months old), your business registration certificate (\textit{Extrait K-bis}) and a minimum opening deposit of 50,000 FCFA.\\
    \textbf{Mrs Mvondo:} I have everything here. One question: can I set up \textbf{standing orders} to pay my suppliers in Yaoundé and Bafoussam automatically?\\
    \textbf{Mr Kouam:} Absolutely. You can schedule recurring transfers via the app. You can also set up \textbf{direct debits} for utility bills like ENEO or Camwater. For international payments to your buyers in Europe, you will need to request a \textbf{SWIFT/BIC code} and provide the beneficiary's \textbf{IBAN}.\\
    \textbf{Mrs Mvondo:} Perfect. I am a bit worried about security. What if I lose my phone?\\
    \textbf{Mr Kouam:} The app uses \textbf{two-factor authentication (2FA)}. If you lose your device, call our 24/7 hotline immediately to block the app access. Your card can be blocked instantly via the app or by calling the number on the back of the card. We also send \textbf{SMS alerts} for every transaction.\\
    \textbf{Mrs Mvondo:} Reassuring. Let's proceed.
    \end{textebox}

    \textbf{Questions :}
    \begin{enumerate}
        \item What type of account does Mrs Mvondo want to open and for what purpose?
        \item List \textbf{three} features included in the “Business Starter” package.
        \item What is the condition to avoid paying the monthly maintenance fee?
        \item Name \textbf{four} documents required to open the account.
        \item Explain the difference between a \eng{standing order} and a \eng{direct debit} based on the dialogue.
        \item What security measures does the bank offer for mobile banking? (Give at least three).
        \item \textbf{Vocabulary in context :} Find in the text the words/expressions matching these definitions:
        \begin{itemize}
            \item[a)] An official document proving where you live: \_\_\_\_\_\_\_\_\_\_\_
            \item[b)] The maximum amount you can take out in 24 hours: \_\_\_\_\_\_\_\_\_\_\_
            \item[c)] A code used for international bank transfers: \_\_\_\_\_\_\_\_\_\_\_
            \item[d)] Money put into the account to activate it: \_\_\_\_\_\_\_\_\_\_\_
        \end{itemize}
    \end{enumerate}
    
    \tcblower
    \textbf{Solution détaillée :}
    \begin{enumerate}
        \item She wants to open a \textbf{current account} for her \textbf{small business} (cocoa beans export venture).
        \item (Any three) : \textbf{Cheque book} ; \textbf{Visa debit card} (with daily limit 500,000 FCFA) ; \textbf{Free mobile app / internet banking access} ; \textbf{Monthly fee waiver condition}.
        \item The \textbf{average monthly balance} must stay \textbf{above 1,000,000 FCFA}.
        \item 1. \textbf{National Identity Card}. 2. \textbf{Proof of address} (bill/tenancy agreement < 3 months). 3. \textbf{Business registration certificate (Extrait K-bis)}. 4. \textbf{Minimum opening deposit (50,000 FCFA)}.
        \item \textbf{Standing order} : Instruction given by the \textbf{account holder (payer)} to pay a \textbf{fixed amount} regularly to another account (suppliers). \textbf{Direct debit} : Authorization given by the \textbf{account holder} to a \textbf{third party (payee/beneficiary)} (ENEO, Camwater) to collect \textbf{variable amounts} from the account.
        \item 1. \textbf{Two-factor authentication (2FA)}. 2. \textbf{24/7 hotline} to block app/card. 3. \textbf{Instant card blocking} via app/phone. 4. \textbf{SMS alerts} for every transaction.
        \item \textbf{Vocabulary :}
        \begin{itemize}
            \item[a)] \textbf{Proof of address} (ou \textbf{tenancy agreement} / \textbf{electricity bill}).
            \item[b)] \textbf{Daily withdrawal limit}.
            \item[c)] \textbf{SWIFT/BIC code}.
            \item[d)] \textbf{Minimum opening deposit}.
        \end{itemize}
    \end{enumerate}
    \textbf{Conclusion :} Ce dialogue couvre le vocabulaire cible (package, maintenance fee, waiver, standing order vs direct debit, SWIFT/IBAN, 2FA, alerts). Notez l'usage de \eng{waived} (exempté/supprimé) et \eng{venture} (entreprise/activité risquée).
\end{exoresolu}

\begin{methode}[Rédiger une correspondance bancaire formelle (Formal Letter / Email)]
    \textbf{Objectif :} Produire un écrit formel (réclamation, demande de renseignements, contestation de frais, clôture de compte) respectant les conventions anglo-saxonnes.
    \begin{enumerate}
        \item \textbf{Structure rigide :}
            \begin{itemize}
                \item \textbf{Sender's Address} (en haut à droite) — \textbf{Date} (sous l'adresse, format : \eng{15th October 2025} ou \eng{October 15, 2025}).
                \item \textbf{Recipient's Address} (à gauche, en dessous de la date) : \eng{The Branch Manager, [Bank Name], [Branch Address]}.
                \item \textbf{Subject Line} (objet, gras, centré ou aligné gauche) : \eng{Re: Complaint regarding unauthorized transaction / Account No: 123456}.
                \item \textbf{Salutation} : \eng{Dear Sir/Madam,} (si nom inconnu) ou \eng{Dear Mr Kouam,}.
                \item \textbf{Body} : 3 paragraphes : \textbf{1. Reason for writing} (référence compte/date). \textbf{2. Details/Evidence} (faits, chiffres, pièces jointes). \textbf{3. Call to action / Expected outcome} (remboursement, explication, rendez-vous).
                \item \textbf{Closing} : \eng{Yours faithfully,} (si \eng{Dear Sir/Madam}) / \eng{Yours sincerely,} (si nom connu).
                \item \textbf{Signature} + \textbf{Printed Name} + \textbf{Account Number} / \textbf{Phone Number}.
            \end{itemize}
        \item \textbf{Registre (Tone) :} Formel, poli, factuel, sans émotion excessive. Utiliser le \textbf{passif} pour l'objectivité : \eng{A fee was charged...}, \eng{The transaction was processed...}. Utiliser les \textbf{modaux de politesse} : \eng{I would be grateful if you could...}, \eng{I look forward to hearing from you...}.
        \item \textbf{Vocabulaire technique imposé :} \eng{Transaction reference}, \eng{Statement}, \eng{Dispute}, \eng{Refund}, \eng{Rectify}, \eng{Overdraft charges}, \eng{Standing order mandate}.
    \end{enumerate}
    \textbf{Modèle de phrase « passe-partout » :} \eng{I am writing to formally dispute a charge of [Amount] on my account [Number] dated [Date], which appears on my statement issued on [Date].}
\end{methode}

\begin{exoresolu}[Rédaction d'un email formel de réclamation — Type Bac : Writing]
    \textbf{Énoncé :} Vous êtes \textbf{Paul Atangana}, titulaire du compte n° \textbf{4455667788} à la \textbf{Afriland First Bank, Agence Mvan (Yaoundé)}. Le 10 octobre 2025, vous avez constaté sur votre relevé bancaire en ligne un prélèvement de \textbf{15,000 FCFA} libellé \eng{“COMMISSION INTERBANCAIRE”} pour un virement reçu de votre cousin en France (montant : 500 EUR). Vous n'avez jamais été informé de cette commission. Rédigez un email formel au \textbf{Directeur d'Agence} pour contester ce prélèvement et demander son remboursement. (Environ 120-150 mots).
    
    \tcblower
    \textbf{Solution détaillée (Modèle de réponse) :}
    
    \begin{textebox}[Formal Email — Complaint]
    \textbf{From:} paul.atangana@email.com \\
    \textbf{To:} branch.manager.mvan@afrilandfirstbank.cm \\
    \textbf{Date:} October 15, 2025 \\
    \textbf{Subject:} \textbf{Re: Dispute of Interbank Commission — Account No. 4455667788}
    
    \textbf{Dear Sir/Madam,}
    
    I am writing to formally dispute a charge of \textbf{15,000 FCFA} debited from my current account (No. 4455667788) on \textbf{October 10, 2025}, described on my online statement as \eng{“Interbank Commission”}.
    
    This fee relates to an incoming international transfer of \textbf{500 EUR} received from a relative abroad. I was not informed of any such commission at the time of opening the account, nor does my account agreement (\eng{“Convention de compte”}) mention a flat fee for receiving foreign currency credits. According to the tariff guide available on your website, incoming SEPA transfers in EUR should be free of charge for the beneficiary.
    
    I would be grateful if you could \textbf{investigate this matter}, \textbf{waive this charge}, and \textbf{reimburse the 15,000 FCFA} to my account within the next five business days. I have attached a screenshot of the transaction details for your reference.
    
    I look forward to your prompt resolution of this issue.
    
    \textbf{Yours faithfully,}
    
    \textbf{Paul Atangana} \\
    Account Holder — 4455667788 \\
    Tel: +237 6XX XX XX XX
    \end{textebox}
    
    \textbf{Analyse des points clés (Grille d'évaluation) :}
    \begin{itemize}
        \item \textbf{Format :} En-têtes email complets, Objet clair (\eng{Re: ... Account No.}), Salutation/Closing formels (\eng{Yours faithfully} car destinataire inconnu).
        \item \textbf{Contenu :} Paragraphe 1 = Objet précis (Montant, Date, Libellé, N° compte). Paragraphe 2 = Argumentation factuelle (Montant virement, origine, référence tarifaire, absence d'information). Paragraphe 3 = Demande d'action concrète (Enquête, Annulation, Remboursement, Délai) + Pièce jointe mentionnée.
        \item \textbf{Langue :} Passif (\eng{was debited, was not informed, should be free}). Modaux polis (\eng{would be grateful, could investigate}). Vocabulaire précis (\eng{incoming transfer, beneficiary, tariff guide, waive, reimburse, screenshot, transaction details}).
        \item \textbf{Orthographe/Grammaire :} Pas de fautes. \eng{Advice} (nom) vs \eng{Advise} (verbe) non piégé ici. \eng{Commission} (frais) correct.
    \end{itemize}
\end{exoresolu}

\begin{methode}[Réussir une simulation orale (Role-play / Guided Conversation)]
    \textbf{Objectif :} Interagir oralement dans une situation bancaire (guichet, téléphone, rendez-vous conseiller) avec fluidité, politesse et vocabulaire précis.
    \begin{enumerate}
        \item \textbf{Analyser la fiche de rôle :} Qui suis-je ? (Client / Employé). Quel est mon but ? (Ouvrir compte, contester, demander prêt, bloquer carte). Quel ton ? (Formel / Courant).
        \item \textbf{Préparer le "script mental" :} Salutation $\to$ Exposé du motif $\to$ Échange d'infos (questions/réponses : pièces justificatives, montants, délais) $\to$ Négociation/Validation $\to$ Conclusion/Remerciements.
        \item \textbf{Maîtriser les "Gambits" (routines interactionnelles) :}
            \begin{itemize}
                \item \textbf{Demander :} \eng{Could you please...?} / \eng{I would like to...} / \eng{May I have...?}
                \item \textbf{Préciser :} \eng{What exactly do you mean by...?} / \eng{Let me check...} / \eng{Just to confirm...}
                \item \textbf{Problème :} \eng{I'm afraid there's an issue...} / \eng{There seems to be a discrepancy...}
                \item \textbf{Clôturer :} \eng{Is there anything else I can help with?} / \eng{Thank you for your time.}
            \end{itemize}
        \item \textbf{Prononciation clé :} Chiffres (\eng{15 / fifty} vs \eng{50 / fifteen}), Montants (\eng{FCFA} = \eng{F-C-F-A} ou \eng{FRANCS}), Acronymes (\eng{IBAN, SWIFT, PIN, ATM, KYC} — \eng{Know Your Customer}). \emph{Stress sur les mots de contenu.}
        \item \textbf{Stratégie si bloqué :} \eng{Sorry, could you repeat that?} / \eng{How do you spell...?} / \eng{Let me put it another way...}
    \end{enumerate}
    \textbf{Critères d'évaluation Bac :} Interaction (réagir), Cohérence, Lexique spécialisé, Grammaire (modaux, conditionnel), Prononciation/Intonation.
\end{methode}

\begin{exoresolu}[Jeu de rôle guidé : Contestations de frais / Négociation découvert — Type Bac : Speaking]
    \textbf{Énoncé :} \textbf{Situation :} Vous êtes \textbf{Client(e)}. Vous appelez votre banque car vous avez reçu une alerte SMS : \eng{“Overdraft limit exceeded. Fees apply.”} Vous pensiez avoir assez d'argent car un virement entrant était prévu hier.
    \textbf{Votre rôle :} 1. Saluer et donner votre identité. 2. Signaler le problème (découvert non prévu). 3. Expliquer qu'un virement (salaire/remboursement) était attendu. 4. Demander si le virement est en attente (\eng{pending}). 5. Négocier l'annulation des frais (\eng{waive fees}) si le virement arrive aujourd'hui. 6. Remercier et clore l'appel.
    
    \textbf{Rédigez le script complet du dialogue (Client + Agent) en anglais.}
    
    \tcblower
    \textbf{Solution détaillée (Modèle de dialogue) :}
    
    \begin{textebox}[Role-play Script: Phone Call — Overdraft Alert]
    \textbf{Agent:} Good morning, Afriland First Bank Customer Service, \textbf{Brenda} speaking. How may I assist you?\\
    \textbf{Client:} Good morning. My name is \textbf{Paul Atangana}, account number \textbf{4455-6677-88}. I am calling about an SMS alert I just received saying my \textbf{overdraft limit has been exceeded}.\\
    \textbf{Agent:} One moment please, Mr Atangana... I can see your account. Indeed, there is a debit interest charge of \textbf{5,000 FCFA} applied this morning as the balance went below zero yesterday evening.\\
    \textbf{Client:} That is surprising. I was expecting a \textbf{transfer of 300,000 FCFA} (my salary) to be credited yesterday. Has it not arrived?\\
    \textbf{Agent:} Let me check the \textbf{pending transactions}... I see an incoming transfer from \textbf{MINEDUB} for 300,000 FCFA. It is currently \textbf{“pending”} (status: \textit{awaiting settlement}). It should be \textbf{cleared} by end of day today.\\
    \textbf{Client:} I see. Since the delay is not on my side — my employer sent it on time — would it be possible to \textbf{waive the overdraft fees}? The funds are effectively arriving today.\\
    \textbf{Agent:} I understand your point, Sir. As a gesture of goodwill and since your account history is clean, I can \textbf{request a fee reversal} from my supervisor. If the transfer is \textbf{posted} today, the fees will be \textbf{refunded} within 48 hours.\\
    \textbf{Client:} Thank you very much, Brenda. I appreciate your help. Will I receive a confirmation SMS once the refund is done?\\
    \textbf{Agent:} Yes, absolutely. You will receive an alert for the \textbf{credit of the salary} and another for the \textbf{fee reversal}. Is there anything else I can help you with today?\\
    \textbf{Client:} No, that’s all. Have a good day.\\
    \textbf{Agent:} You too, Mr Atangana. Goodbye.
    \end{textebox}
    
    \textbf{Analyse lexicale et grammaticale (Points forts pour le Bac) :}
    \begin{itemize}
        \item \textbf{Vocabulaire technique :} \eng{Overdraft limit exceeded, Debit interest, Pending transactions, Awaiting settlement, Cleared/Posted, Fee reversal/waiver, Refund, Credit alert}.
        \item \textbf{Structures fonctionnelles :} \eng{I am calling about...} (motif). \eng{I was expecting...} (passé progressif pour attente passée). \eng{Would it be possible to...?} (négociation polie). \eng{As a gesture of goodwill...} (formule bancaire standard).
        \item \textbf{Prononciation :} \eng{Account number} (chiffres groupés par 4). \eng{Pending} /ˈpendɪŋ/. \eng{Cleared} /klɪəd/. \eng{Waive} « waïv » (le \textit{w} ne se prononce pas).
    \end{itemize}
\end{exoresolu}

\begin{attention}[Les 5 erreurs qui coûtent des points au Bac — Spécial Banking Vocabulary]
    \begin{enumerate}
        \item \textbf{Faux ami \eng{Benefit} vs \eng{Profit} / \eng{Loan}.}
        \eng{Benefit} = Avantage, prestation sociale (\eng{unemployment benefit}). \textbf{Pas} « bénéfice » (comptable = \eng{Profit}) ni « prêt » (\eng{Loan/Credit}). \eng{*\textbf{The bank gave me a benefit of 1M}} $\to$ \eng{The bank granted me a \textbf{loan} of 1M}.
        
        \item \textbf{\eng{Interest} au pluriel (\eng{*Interests}).}
        \eng{Interest} (taux d'intérêt, intérêt financier) est \textbf{incomptable (uncountable)}. Jamais de \textit{s}. \eng{*\textbf{The interests rates are high}} $\to$ \eng{\textbf{Interest} rates are high}. (Seul \eng{Interests} au pluriel = centres d'intérêt / hobbies).
        
        \item \textbf{Confusion \eng{Standing Order} (Ordre permanent) vs \eng{Direct Debit} (Prélèvement automatique).}
        \eng{Standing Order} : \textbf{Client} $\to$ Banque : « Payez X € à Y tous les mois ». Montant \textbf{fixe}. Initiateur = \textbf{Débiteur (Payeur)}.
        \eng{Direct Debit} : \textbf{Bénéficiaire (Créancier)} $\to$ Banque : « Prélevez ce que je dois sur le compte de X ». Montant \textbf{variable}. Initiateur = \textbf{Créancier}.
        \textbf{Piège Bac :} « Mettre en place un prélèvement pour l'électricité » $\to$ \eng{Set up a \textbf{direct debit}} (pas \eng{standing order}).
        
        \item \textbf{Verbe \eng{Depose} (n'existe pas pour l'argent) / \eng{Deposit} (verbe + nom).}
        \eng{*\textbf{I want to depose money}} $\to$ \eng{I want to \textbf{deposit} money / \textbf{make a deposit}}. \eng{Depose} = déposer une arme, un dossier, un corps.
        \textbf{Orthographe :} \eng{Deposit} (nom et verbe). \eng{Depositor} (déposant).
        
        \item \textbf{Prépositions après verbes/adjectifs bancaires (Calques français).}
        \begin{center}
            \begin{tabularx}{\linewidth}{|X|X|X|}\hline 
            \rowcolor{poporangeL} \textbf{Français (Calque)} & \textbf{Anglais Correct} & \textbf{Exemple} \\ \hline
            Demander \textbf{un} prêt & \eng{Apply \textbf{for} a loan} & \eng{He applied \textbf{for} a mortgage.} \\ \hline
            S'intéresser \textbf{à} & \eng{Be interested \textbf{in}} & \eng{She is interested \textbf{in} investing.} \\ \hline
            Payer \textbf{pour} qqch & \eng{Pay \textbf{for} sth} (mais \eng{Pay a bill / Pay fees} \textbf{sans} préposition) & \eng{He paid \textbf{for} the transfer fees.} / \eng{He paid \textbf{the bill.}} \\ \hline
            Investir \textbf{dans} & \eng{Invest \textbf{in}} & \eng{They invested \textbf{in} bonds.} \\ \hline
            Emprunter \textbf{à} qqn & \eng{Borrow \textbf{from} sb} & \eng{I borrowed 10k \textbf{from} my brother.} \\ \hline
            Prêter \textbf{à} qqn & \eng{Lend \textbf{to} sb} / \eng{Lend sb sth} & \eng{The bank lent \textbf{me} money.} \\ \hline
        \end{tabularx}
        \end{center}
    \end{enumerate}
\end{attention}

\newpage

\coursec{Exercices}

\coursec{Lesson 18 — Vocabulary: Words and expressions related to banking services}

\begin{aretenir}[Key Vocabulary: Banking Services — Mindmap]
Voici les trois champs lexicaux essentiels à maîtriser pour cette leçon. Retiens les collocations (associations de mots) typiques.

\begin{popfigure}
\begin{tikzpicture}[mindmap, grow cyclic, every node/.style={concept, execute at begin node=\hskip0pt},
  concept color=popblue!80, text=white,
  level 1/.append style={level distance=3.5cm, sibling angle=120, font=\small\bfseries},
  level 2/.append style={level distance=2.8cm, sibling angle=50, font=\footnotesize}]
\node {Banking Services}
  child[concept color=popgreen] { node {People \& Places}
    child { node {bank clerk / teller} }
    child { node {customer / account holder} }
    child { node {borrower / lender} }
    child { node {branch / head office} }
    child { node {cash machine / ATM} }
  }
  child[concept color=poporange] { node {Accounts \& Operations}
    child { node {current / savings account} }
    child { node {deposit / withdrawal} }
    child { node {transfer / standing order} }
    child { node {balance / bank statement} }
    child { node {fees / charges} }
  }
  child[concept color=poppink] { node {Payment, Credit \& Savings}
    child { node {cash / cheque / card} }
    child { node {mobile money / app} }
    child { node {loan / overdraft} }
    child { node {interest rate / APR} }
    child { node {investment / savings} }
  };
\end{tikzpicture}
\legende{Les trois piliers du vocabulaire bancaire}
\end{popfigure}
\end{aretenir}

\courssub{Je vérifie mes connaissances}

\begin{exercice}[QCM : le bon mot]
Choisis la réponse correcte (a, b ou c).
\begin{enumerate}
\item Money you keep in a bank is put into \ldots
\begin{itemize}
\item[a)] a safe \quad b) an account \quad c) a wallet
\end{itemize}
\item The money a bank charges you for its services is called \ldots
\begin{itemize}
\item[a)] fees \quad b) fares \quad c) prizes
\end{itemize}
\item When you take money out of your account, you make a \ldots
\begin{itemize}
\item[a)] deposit \quad b) withdrawal \quad c) transfer
\end{itemize}
\item The extra money you pay back on a loan is the \ldots
\begin{itemize}
\item[a)] balance \quad b) interest \quad c) overdraft
\end{itemize}
\item A person who borrows money from a bank is a \ldots
\begin{itemize}
\item[a)] lender \quad b) borrower \quad c) cashier
\end{itemize}
\end{enumerate}
\end{exercice}

\begin{exercice}[Vrai ou faux ? Justifie]
Dis si chaque affirmation est vraie ou fausse, puis justifie en une phrase en anglais.
\begin{enumerate}
\item An overdraft means you have spent more money than you have in your account.
\item You need to close an account before you can deposit money in it.
\item A savings account usually earns interest.
\item \eng{To lend} and \eng{to borrow} mean exactly the same thing.
\item Mobile money allows people to send money without going to a bank branch.
\end{enumerate}
\end{exercice}

\begin{exercice}[Vocabulaire : définitions]
Trouve le mot du champ lexical bancaire qui correspond à chaque définition. Choisis dans la banque de mots : \eng{bank statement, cash machine, loan, balance, transfer, cheque}.
\begin{enumerate}
\item A printed or electronic document showing all the operations on your account.
\item A machine where you can withdraw money with a card (also called an ATM).
\item The amount of money available in your account at a given moment.
\item Money that a bank lends you and that you must repay with interest.
\item The action of moving money from one account to another.
\item A written order telling a bank to pay money from your account to someone.
\end{enumerate}
\end{exercice}

\begin{exercice}[Familles de mots : à compléter]
Complète le tableau avec la forme correcte de chaque famille de mots.
\begin{center}
\begin{tabular}{|l|l|l|}
\hline
\rowcolor{popblueL} Verbe & Personne (Nom agent) & Nom d'action / Résultat \\
\hline
to lend & \ldots\ldots\ldots & \ldots\ldots\ldots \\
\hline
to save & \ldots\ldots\ldots & \ldots\ldots\ldots \\
\hline
to invest & \ldots\ldots\ldots & \ldots\ldots\ldots \\
\hline
to borrow & \ldots\ldots\ldots & \ldots\ldots\ldots \\
\hline
to withdraw & \ldots\ldots\ldots & \ldots\ldots\ldots \\
\hline
\end{tabular}
\end{center}
\end{exercice}

\courssub{Je m'entraîne}

\begin{exercice}[Phrases à trous]
Complète les phrases avec les mots de la banque suivante (chaque mot n'est utilisé qu'une fois) : \eng{deposit, withdrawal, balance, loan, interest, overdraft, fees, account}.
\begin{enumerate}
\item Mr Etoundi went to the bank in Yaoundé to open a savings \ldots\ldots\ldots .
\item She made a \ldots\ldots\ldots of 50,000 francs CFA into her account on Monday.
\item After paying all her bills, her \ldots\ldots\ldots was only 5,000 francs.
\item The cooperative gave the farmers a \ldots\ldots\ldots to buy fertiliser before the cocoa season.
\item If you repay your loan late, the \ldots\ldots\ldots you owe will increase.
\item The bank charges monthly \ldots\ldots\ldots for managing the account.
\item He spent more than he had, so his account is now in \ldots\ldots\ldots .
\item The cash machine refused the \ldots\ldots\ldots because the card was blocked.
\end{enumerate}
\end{exercice}

\begin{exercice}[Collocations : à relier]
Relie chaque verbe de la colonne A au complément correct de la colonne B, puis écris une phrase complète avec deux des collocations obtenues.
\begin{center}
\begin{tabular}{|l|l|}
\hline
\rowcolor{popblueL} A : verbes & B : compléments \\
\hline
1. to open & a) a loan \\
\hline
2. to make & b) interest \\
\hline
3. to check & c) an account \\
\hline
4. to take out & d) a deposit \\
\hline
5. to pay & e) your balance \\
\hline
\end{tabular}
\end{center}
\end{exercice}

\begin{exercice}[Faux amis : corrige les phrases]
Chaque phrase contient une erreur due à un calque du français. Souligne l'erreur et réécris la phrase correctement.
\begin{enumerate}
\item The bank saved my money in a safe.
\item I made a credit to buy a motorbike.
\item She has a big economy account at the bank.
\item The banker demanded me my identity card.
\item I must retire some money from the cash machine.
\end{enumerate}
\end{exercice}

\begin{exercice}[Traduction vers l'anglais]
Traduis les phrases suivantes en anglais correct et naturel.
\begin{enumerate}
\item Ma grand-mère participe à une tontine chaque mois dans notre quartier.
\item De nombreux jeunes à Douala utilisent le mobile money pour payer leurs achats.
\item Mon père a demandé un prêt à la banque pour acheter une moto-taxi.
\item Avant de retirer de l'argent, vérifie toujours ton solde.
\item Les élèves devraient apprendre à épargner une partie de leur argent de poche.
\end{enumerate}
\end{exercice}

\courssub{Je me prépare au Bac}

\begin{exercice}[Compréhension de texte : mobile money in Cameroon]
Lis le texte suivant, puis réponds aux questions.

\begin{textebox}[The rise of mobile money in Cameroon]
Ten years ago, sending money from Douala to Bamenda meant a long bus journey or an expensive bank transfer. Today, a market woman in Mokolo can receive payment for her vegetables in a few seconds, using only her mobile phone. Mobile money has changed the economic life of millions of Cameroonians.

The system is simple. A customer opens a mobile money account with a telephone number, deposits cash at a local agent's kiosk, and can then send money, pay bills or buy airtime. No bank branch, no paperwork, no minimum balance is required. For people in rural areas, where the nearest bank may be fifty kilometres away, this is a revolution.

However, mobile money also has its limits. Transaction fees can be high, and customers who lose their phones risk losing their savings if their account is not protected by a secret code. Moreover, unlike a savings account in a bank, money kept on a phone usually earns no interest.

Bankers have understood the message. Many banks now offer mobile applications that allow customers to check their balance, make withdrawals without a card and even apply for small loans. The future of banking in Cameroon, it seems, will fit in a pocket.
\end{textebox}

\textbf{A. Comprehension questions}
\begin{enumerate}
\item Say whether the following statements are true or false. Justify each answer with a quotation from the text.
\begin{enumerate}
\item[a)] Sending money inside Cameroon was easy and cheap ten years ago.
\item[b)] A customer needs a minimum balance to open a mobile money account.
\item[c)] Money kept on a phone earns the same interest as a bank savings account.
\end{enumerate}
\item Answer the following questions in your own words.
\begin{enumerate}
\item[a)] Give two advantages of mobile money mentioned in the text.
\item[b)] Give two disadvantages or risks of mobile money.
\item[c)] How have traditional banks reacted to the success of mobile money?
\end{enumerate}
\item Vocabulary: find in the text
\begin{enumerate}
\item[a)] the noun from the same family as \eng{to withdraw};
\item[b)] a word from the same family as \eng{to lend};
\item[c)] a word meaning \eng{documents and administrative formalities}.
\end{enumerate}
\end{enumerate}
\end{exercice}

\begin{exercice}[Traduction (thème)]
Traduis le passage suivant en anglais.

\begin{textebox}[Texte à traduire]
Hier matin, Amina est allée à la banque pour ouvrir un compte épargne. L'employé lui a expliqué les frais et le taux d'intérêt. Elle a déposé vingt mille francs et a reçu une carte bancaire. Maintenant, elle peut vérifier son solde et retirer de l'argent au distributeur automatique près du marché.
\end{textebox}
\end{exercice}

\begin{exercice}[Expression écrite : sujet de composition]
\begin{textebox}[Writing task]
Your school magazine is preparing a special issue on money and young people. Write an article of about 150 words entitled: “Why every student should learn to use banking services wisely.”

In your article, you should:
\begin{itemize}
\item explain why banking services (accounts, savings, loans, mobile money) are useful for young people;
\item warn against one or two dangers (debts, overdrafts, unnecessary fees);
\item give two pieces of advice to your schoolmates.
\end{itemize}
\end{textebox}
Rédige ton article en environ 150 mots. Soigne l'organisation (introduction, développement en deux ou trois paragraphes, conclusion) et réemploie au moins huit mots du champ lexical bancaire étudiés dans la leçon.
\end{exercice}



\newpage

\coursec{Corrigés des exercices}

\coursec{Lesson 18 — Vocabulary: Words and expressions related to banking services — Corrigés des exercices}

\begin{corrige}[Exercice 1 — QCM : le bon mot]
\begin{enumerate}
\item \textbf{b) an account}. \eng{Account} = compte bancaire. \eng{Safe} = coffre-fort ; \eng{wallet} = portefeuille. \eng{“Money you keep in a bank is put into an account.”}
\item \textbf{a) fees}. \eng{Fees} = frais (bancaires). \eng{Fares} = tarifs de transport ; \eng{prizes} = prix, récompenses.
\item \textbf{b) withdrawal}. \eng{Withdrawal} = retrait. \eng{Deposit} = dépôt ; \eng{transfer} = virement.
\item \textbf{b) interest}. \eng{Interest} = intérêt (l'argent payé en plus sur un prêt). \eng{Balance} = solde ; \eng{overdraft} = découvert.
\item \textbf{b) borrower}. \eng{Borrower} = emprunteur (celui qui prend l'argent). \eng{Lender} = prêteur (celui qui donne) ; \eng{cashier} = caissier.
\end{enumerate}
\textbf{Points de vigilance :} ne pas confondre \cle{borrower} et \cle{lender} : le suffixe \eng{-er} désigne la personne qui fait l'action du verbe (\eng{borrow} = emprunter, \eng{lend} = prêter).
\end{corrige}

\begin{corrige}[Exercice 2 — Vrai ou faux ? Justifie]
\begin{enumerate}
\item \textbf{True.} \eng{“An overdraft means you have spent more money than you have in your account, so your balance is negative.”} (Un découvert = solde négatif.)
\item \textbf{False.} \eng{“You must open an account before you can deposit money in it.”} (On dépose de l'argent \emph{sur} un compte ouvert ; le fermer empêche tout dépôt.)
\item \textbf{True.} \eng{“A savings account usually earns interest because the bank pays you for keeping your money.”}
\item \textbf{False.} \eng{“To lend means to give money to someone, whereas to borrow means to take money from someone; they are opposites.”} (\eng{lend} = prêter ; \eng{borrow} = emprunter.)
\item \textbf{True.} \eng{“Mobile money allows people to send money through their phone and local agents, without going to a bank branch.”}
\end{enumerate}
\textbf{Points de vigilance :} en anglais, la justification doit être une phrase complète (sujet + verbe). Évite \eng{“it is true because yes”} !
\end{corrige}

\begin{corrige}[Exercice 3 — Vocabulaire : définitions]
\begin{enumerate}
\item \cle{bank statement} — relevé de compte. \eng{“I receive my bank statement every month.”}
\item \cle{cash machine} — distributeur automatique (aussi \cle{ATM}, surtout en anglais américain).
\item \cle{balance} — solde. \eng{“What is your current balance?”}
\item \cle{loan} — prêt. \eng{“She repaid her loan in twelve months.”}
\item \cle{transfer} — virement. \eng{“He made a transfer to his sister's account.”}
\item \cle{cheque} — chèque (orthographe britannique ; américain : \eng{check}).
\end{enumerate}
\textbf{Points de vigilance :} attention à l'orthographe de \eng{cheque} (« chèque ») et à la prononciation de \eng{loan} (« lone »).
\end{corrige}

\begin{corrige}[Exercice 4 — Familles de mots : à compléter]
\begin{center}
\begin{tabular}{|l|l|l|}
\hline
\rowcolor{popblueL} Verbe & Personne (Nom agent) & Nom d'action / Résultat \\
\hline
to lend & \textbf{lender} & \textbf{lending} / \textbf{loan} \\
\hline
to save & \textbf{saver} & \textbf{saving(s)} \\
\hline
to invest & \textbf{investor} & \textbf{investment} \\
\hline
to borrow & \textbf{borrower} & \textbf{borrowing} \\
\hline
to withdraw & \textbf{withdrawer} & \textbf{withdrawal} \\
\hline
\end{tabular}
\end{center}
\textbf{Justifications :}
\begin{itemize}
\item Le nom d'agent se forme souvent avec le suffixe \eng{-er} (\eng{lend} $\rightarrow$ \eng{lender}) ou \eng{-or} (\eng{invest} $\rightarrow$ \eng{investor}).
\item \cle{Loan} est le nom associé à \eng{lend} : \eng{“The bank granted him a loan.”} ; \eng{lending} désigne l'activité de prêt.
\item \cle{Savings} s'emploie au pluriel pour désigner l'argent mis de côté : \eng{“She keeps her savings in a bank.”}
\item \cle{Withdrawal} : nom en \eng{-al} formé sur \eng{withdraw} (verbe irrégulier : \eng{withdraw – withdrew – withdrawn}).
\end{itemize}
\end{corrige}

\begin{popfigure}
\begin{tikzpicture}[
  mindmap,
  grow cyclic,
  every node/.style={concept, execute at begin node=\hskip0pt},
  concept color=popblue!80,
  level 1/.append style={level distance=4cm, sibling angle=72, concept color=popblue!60},
  level 2/.append style={level distance=3cm, sibling angle=45, concept color=popgreen!60},
  text=white,
  font=\small\bfseries
]
\node {Banking Verbs}
  child { node {to lend} child { node {lender} } child { node {loan / lending} } }
  child { node {to save} child { node {saver} } child { node {savings} } }
  child { node {to invest} child { node {investor} } child { node {investment} } }
  child { node {to borrow} child { node {borrower} } child { node {borrowing} } }
  child { node {to withdraw} child { node {withdrawer} } child { node {withdrawal} } };
\end{tikzpicture}
\legende{Familles de mots : verbes bancaires et leurs dérivés nominaux}
\end{popfigure}

\begin{corrige}[Exercice 5 — Phrases à trous]
\begin{enumerate}
\item \cle{account} — \eng{“Mr Etoundi went to the bank in Yaoundé to open a savings account.”} (collocation : \cle{open an account}.)
\item \cle{deposit} — \eng{“She made a deposit of 50,000 francs CFA into her account on Monday.”} (collocation : \cle{make a deposit}.)
\item \cle{balance} — \eng{“After paying all her bills, her balance was only 5,000 francs.”} (le solde = ce qui reste.)
\item \cle{loan} — \eng{“The cooperative gave the farmers a loan to buy fertiliser before the cocoa season.”}
\item \cle{interest} — \eng{“If you repay your loan late, the interest you owe will increase.”} (l'intérêt augmente avec le retard.)
\item \cle{fees} — \eng{“The bank charges monthly fees for managing the account.”} (collocation : \cle{charge fees}.)
\item \cle{overdraft} — \eng{“He spent more than he had, so his account is now in overdraft.”} (on dit aussi \cle{to be overdrawn}.)
\item \cle{withdrawal} — \eng{“The cash machine refused the withdrawal because the card was blocked.”}
\end{enumerate}
\textbf{Points de vigilance :} \eng{fees} est presque toujours au pluriel ; \eng{interest} est ici indénombrable (pas de \eng{an interest}).
\end{corrige}

\begin{corrige}[Exercice 6 — Collocations : à relier]
\begin{enumerate}
\item \cle{to open} $\rightarrow$ \textbf{c) an account} — \eng{“I opened an account last week.”}
\item \cle{to make} $\rightarrow$ \textbf{d) a deposit} — \eng{“She made a deposit of 20,000 francs.”}
\item \cle{to check} $\rightarrow$ \textbf{e) your balance} — \eng{“He checked his balance on the app.”}
\item \cle{to take out} $\rightarrow$ \textbf{a) a loan} — \eng{“They took out a loan to build a house.”}
\item \cle{to pay} $\rightarrow$ \textbf{b) interest} — \eng{“The bank pays interest on savings.”}
\end{enumerate}
\textbf{Phrases proposées (exemples) :}
\begin{itemize}
\item \eng{“Before you can make a deposit, you must open an account.”}
\item \eng{“My uncle took out a loan, and now he pays interest every month.”}
\end{itemize}
\textbf{Points de vigilance :} on dit \cle{take out a loan} (et non \eng{make a loan} pour l'emprunteur) ; \eng{do a deposit} est un calque du français « faire un dépôt » — la bonne collocation est \cle{make a deposit}.
\end{corrige}

\begin{corrige}[Exercice 7 — Faux amis : corrige les phrases]
\begin{enumerate}
\item Erreur : \eng{saved}. Correction : \eng{“The bank \textbf{kept} my money in a safe.”} (\eng{save} = sauver / épargner ; « garder, conserver » = \eng{keep} ; « déposer » = \eng{deposit}.)
\item Erreur : \eng{made a credit}. Correction : \eng{“I \textbf{took out a loan} to buy a motorbike.”} (\eng{make a credit} n'existe pas ; \eng{credit} est une notion abstraite, \eng{loan} est la somme prêtée.)
\item Erreur : \eng{economy account}. Correction : \eng{“She has a big \textbf{savings} account at the bank.”} (\eng{economy} = l'économie du pays ; « économies » (argent épargné) = \cle{savings}.)
\item Erreur : \eng{demanded me}. Correction : \eng{“The banker \textbf{asked me for} my identity card.”} (\eng{demand} = exiger, ton très fort ; « demander quelque chose à quelqu'un » = \cle{ask somebody for something}.)
\item Erreur : \eng{retire}. Correction : \eng{“I must \textbf{withdraw} some money from the cash machine.”} (\eng{retire} = prendre sa retraite ; « retirer de l'argent » = \cle{withdraw money}.)
\end{enumerate}
\textbf{Points de vigilance :} ces cinq erreurs sont des calques directs du français (\emph{sauver, faire un crédit, économie, demander, retirer}). Au Bac, chaque faux ami coûte des points en expression écrite : apprends les paires \eng{save/keep}, \eng{economy/savings}, \eng{demand/ask for}, \eng{retire/withdraw}.
\end{corrige}

\begin{attention}[Top 5 des faux amis bancaires (Français $\rightarrow$ Anglais)]
\begin{itemize}
\item \textbf{Sauver / Épargner} $\neq$ \eng{save} (sauver) $\rightarrow$ \eng{keep} (garder), \eng{deposit} (déposer), \eng{save money} (épargner).
\item \textbf{Économie (argent mis de côté)} $\neq$ \eng{economy} (l'économie du pays) $\rightarrow$ \cle{savings}.
\item \textbf{Faire un crédit} $\neq$ \eng{make a credit} $\rightarrow$ \cle{take out a loan} (emprunter), \cle{grant a loan} (prêter, pour la banque).
\item \textbf{Demander (qqch à qqn)} $\neq$ \eng{demand} (exiger) $\rightarrow$ \cle{ask somebody for something}.
\item \textbf{Retirer (de l'argent)} $\neq$ \eng{retire} (prendre sa retraite) $\rightarrow$ \cle{withdraw money}.
\end{itemize}
\end{attention}

\begin{corrige}[Exercice 8 — Traduction vers l'anglais]
\begin{enumerate}
\item \eng{“My grandmother takes part in a tontine every month in our neighbourhood.”} (\cle{take part in} = participer à ; \cle{tontine} est accepté en anglais, contexte camerounais.)
\item \eng{“Many young people in Douala use mobile money to pay for their purchases.”} (attention : \cle{pay for} + nom ; \cle{young people} = les jeunes.)
\item \eng{“My father applied for a loan from the bank to buy a motorbike taxi.”} (\cle{apply for a loan} = demander un prêt ; on peut ajouter \eng{(bendskin)} entre parenthèses.)
\item \eng{“Before withdrawing money, always check your balance.”} (\eng{before} + verbe en \eng{-ing} ; impératif : \eng{check}.)
\item \eng{“Students should learn to save part of their pocket money.”} (\eng{should} + base verbale ; \cle{pocket money} = argent de poche.)
\end{enumerate}
\textbf{Points de vigilance :} \eng{money} est indénombrable (jamais \eng{monies} ici) ; \cle{to pay for something} exige la préposition \eng{for} devant l'objet acheté.
\end{corrige}

\begin{corrige}[Exercice 9 — Compréhension de texte : mobile money in Cameroon]
\textbf{A1. True or false (avec citation) :}
\begin{enumerate}
\item[a)] \textbf{False.} \eng{“Ten years ago, sending money from Douala to Bamenda meant a long bus journey or an expensive bank transfer.”}
\item[b)] \textbf{False.} \eng{“No bank branch, no paperwork, no minimum balance is required.”}
\item[c)] \textbf{False.} \eng{“Unlike a savings account in a bank, money kept on a phone usually earns no interest.”}
\end{enumerate}
\textbf{A2. Questions ouvertes (dans tes propres mots) :}
\begin{enumerate}
\item[a)] \emph{Advantages :} \eng{“Mobile money is very fast — payments take only a few seconds — and it is accessible everywhere, even in rural areas far from any bank branch. Moreover, no paperwork and no minimum balance are needed.”}
\item[b)] \emph{Disadvantages :} \eng{“Transaction fees can be high, and customers risk losing their savings if they lose their phone without a secret code. Besides, money kept on a phone earns no interest.”}
\item[c)] \emph{Banks' reaction :} \eng{“Traditional banks have created mobile applications that allow customers to check their balance, make withdrawals without a card and apply for small loans.”}
\end{enumerate}
\textbf{A3. Vocabulary :}
\begin{enumerate}
\item[a)] \cle{withdrawals} (nom de la famille de \eng{to withdraw} : “make withdrawals without a card”).
\item[b)] \cle{loans} (famille de \eng{to lend} : “apply for small loans”).
\item[c)] \cle{paperwork} (= documents et formalités administratives : “no paperwork”).
\end{enumerate}
\textbf{Barème indicatif (type Bac) :} A1 : 3 $\times$ 2 pts (1 pt pour la réponse, 1 pt pour la citation exacte) ; A2 : 3 $\times$ 2 pts (idée correcte 1 pt + anglais correct 1 pt) ; A3 : 3 $\times$ 1 pt. Total : 15 pts.
\textbf{Points de vigilance :} pour A1, la citation doit être recopiée \emph{mot à mot} entre guillemets ; pour A2, il est interdit de recopier le texte : reformule (\eng{fast} au lieu de \eng{a few seconds}, \eng{created} au lieu de \eng{offer}).
\end{corrige}

\begin{savaistu}[Mobile Money au Cameroun : inclusion financière]
Le \cle{mobile money} (Orange Money, MTN MoMo) a révolutionné l'accès aux services financiers au Cameroun. Il permet aux populations non bancarisées (zones rurales, jeunes, commerçants informels) d'envoyer de l'argent, de payer des factures, d'épargner et d'accéder au microcrédit via un simple téléphone portable, sans \cle{bank branch} ni \cle{paperwork}. C'est un levier majeur pour l'inclusion financière dans la zone CEMAC.
\end{savaistu}

\begin{corrige}[Exercice 10 — Traduction (thème)]
\textbf{Proposition de traduction :}

\eng{“Yesterday morning, Amina went to the bank to open a savings account. The clerk explained the fees and the interest rate to her. She deposited twenty thousand francs and received a bank card. Now she can check her balance and withdraw money from the cash machine near the market.”}

\textbf{Justifications et variantes :}
\begin{itemize}
\item « Hier matin » : \eng{yesterday morning} (pas d'article, pas de préposition).
\item « un compte épargne » : \eng{a savings account} (avec un \eng{s} : faux ami avec \eng{economy}).
\item « L'employé lui a expliqué » : \eng{the clerk explained ... to her} — \eng{explain} est suivi de \eng{to} + personne (jamais \eng{explained her}).
\item « elle a déposé » : \eng{she deposited} (prétérit régulier) ; « vingt mille francs » : \eng{twenty thousand francs} (\eng{thousand} sans \eng{s} après un nombre).
\item « le taux d'intérêt » : \eng{the interest rate} ; « le distributeur automatique » : \eng{the cash machine} ou \eng{the ATM}.
\end{itemize}
\textbf{Barème indicatif :} 10 pts — 2 pts par phrase (lexique bancaire correct 1 pt + grammaire/prétérit 1 pt).
\end{corrige}

\begin{corrige}[Exercice 11 — Expression écrite : sujet de composition]
\textbf{Proposition de rédaction modèle (environ 150 mots) :}

\begin{textebox}[Model article — “Why every student should learn to use banking services wisely”]
Money is part of our daily life, even at school. Many students receive pocket money or earn a little during the holidays. Learning to use banking services wisely is therefore essential.

First, banking services help young people manage their money safely. Opening a savings account teaches us to save regularly instead of spending everything at the market. With mobile money, we can make deposits and transfers in a few seconds, even in villages where there is no bank branch. Moreover, a savings account earns interest, so our money grows slowly but surely.

However, banking services also have dangers. A student who takes out a loan too easily may fall into debt, and an overdraft leads to heavy fees. We must also be careful with mobile money: a lost phone can mean lost savings.

My advice is simple: always check your balance before a withdrawal, and save a small part of your pocket money every month. Wise bankers today become successful adults tomorrow.
\end{textebox}

\textbf{Analyse du modèle :}
\begin{itemize}
\item \textbf{Structure :} introduction (le sujet), deux paragraphes de développement (avantages / dangers), conclusion (deux conseils) — conforme à la consigne.
\item \textbf{Lexique bancaire réemployé (10 occurrences) :} \cle{savings account}, \cle{save}, \cle{deposits}, \cle{transfers}, \cle{bank branch}, \cle{interest}, \cle{take out a loan}, \cle{debt}, \cle{overdraft}, \cle{fees}, \cle{check your balance}, \cle{withdrawal}, \cle{savings}.
\item \textbf{Connecteurs :} \eng{first, moreover, however, also} — indispensables pour la note de cohérence.
\end{itemize}
\textbf{Barème indicatif (type Bac, 20 pts) :} respect de la tâche et des trois points de la consigne : 6 pts ; richesse et exactitude du lexique bancaire (au moins 8 mots) : 6 pts ; correction grammaticale (temps, articles, prépositions) : 5 pts ; organisation et connecteurs : 3 pts.
\textbf{Points de vigilance :} évite les calques (\eng{make economies} $\rightarrow$ \eng{save money}) ; respecte la limite de 150 mots ($\pm$ 10 \%) ; n'oublie pas le titre de l'article ; accorde les verbes à la troisième personne (\eng{the account earns}, \eng{it grows}).
\end{corrige}

\newpage

\coursec{Fiche bilan}

\begin{popfigure}
\begin{tikzpicture}[mindmap, grow cyclic, every node/.style={concept, font=\small}, concept color=popblue!60,
level 1/.append style={level distance=3.2cm, sibling angle=60, font=\small},
level 2/.append style={level distance=1.8cm, sibling angle=45, font=\scriptsize}]
\node[concept color=popblue, text=white, font=\bfseries\small, align=center]{Banking\\Services}
  child[concept color=poporange!70] { node[align=center]{People \&\\Places} 
    child { node {bank teller} }
    child { node {branch} } }
  child[concept color=popgreen!70] { node {Accounts} 
    child { node[align=center]{current /\\savings} }
    child { node {open / close} } }
  child[concept color=popgold!80] { node {Operations} 
    child { node {deposit} }
    child { node {withdrawal} }
    child { node {transfer} } }
  child[concept color=poppurple!60] { node[align=center]{Payment\\methods} 
    child { node {cheque, card} }
    child { node {cash} } }
  child[concept color=poppink!70] { node[align=center]{Credit \&\\Savings} 
    child { node {loan, interest} }
    child { node {overdraft} } }
  child[concept color=popgreen!50] { node[align=center]{Faux amis\\et familles} 
    child { node[align=center]{to save\\(time/money)} } };
\end{tikzpicture}
\legende{Carte mentale : le vocabulaire des services bancaires en six champs}
\end{popfigure}

\begin{bilan}
\textbf{1. Lexique clé à connaître par cœur}
\begin{itemize}
\item \cle{People and places} : \eng{a bank} (une banque), \eng{a branch} (une agence), \eng{a bank teller / cashier} (un guichetier), \eng{a customer} (un client), \eng{a bank manager} (un directeur d'agence), \eng{an ATM / cash machine} (un distributeur automatique).
\item \cle{Accounts} : \eng{a bank account} (un compte bancaire), \eng{a current account} (un compte courant), \eng{a savings account} (un compte d'épargne), \eng{to open / close an account} (ouvrir / clôturer un compte), \eng{an account number} (un numéro de compte), \eng{a bank statement} (un relevé de compte), \eng{a balance} (un solde).
\item \cle{Operations} : \eng{to deposit money} (déposer de l'argent), \eng{a deposit} (un dépôt), \eng{to withdraw money} (retirer de l'argent), \eng{a withdrawal} (un retrait), \eng{to transfer money} (virer de l'argent), \eng{a transfer} (un virement), \eng{to exchange currency} (changer de la monnaie), \eng{the exchange rate} (le taux de change), \eng{bank charges / fees} (les frais bancaires).
\item \cle{Payment methods} : \eng{cash} (les espèces), \eng{to pay in cash / by cheque / by card} (payer en espèces / par chèque / par carte), \eng{a debit card / a credit card} (une carte de débit / de crédit), \eng{a PIN number} (un code confidentiel), \eng{mobile money} (l'argent mobile).
\item \cle{Credit and savings} : \eng{to save money} (économiser), \eng{savings} (l'épargne), \eng{to borrow money from the bank} (emprunter à la banque), \eng{to lend money to somebody} (prêter à quelqu'un), \eng{a loan} (un prêt), \eng{interest} (les intérêts), \eng{the interest rate} (le taux d'intérêt), \eng{an overdraft} (un découvert), \eng{to be in debt} (être endetté), \eng{to repay a loan} (rembourser un prêt).
\end{itemize}

\textbf{2. Points de langue à maîtriser}
\begin{center}
\begin{tabularx}{\linewidth}{|X|X|}\hline
\rowcolor{popblueL} \textbf{Règle} & \textbf{Exemple} \\ \hline
\eng{money} est indénombrable : jamais de pluriel ni de \eng{a} & \eng{I need some money.} (jamais \emph{a money / moneys}) \\ \hline
\eng{to borrow from} (emprunter à) $\neq$ \eng{to lend to} (prêter à) & \eng{I borrowed 50,000 FCFA from the bank.} / \eng{The bank lent me the money.} \\ \hline
\eng{to pay by} + moyen, \eng{to pay in} + espèces & \eng{to pay by cheque}, \eng{to pay in cash} \\ \hline
\eng{to save} = économiser (argent) et gagner (temps) ; \eng{savings} toujours au pluriel & \eng{She keeps her savings in a bank.} \\ \hline
\end{tabularx}
\end{center}

\textbf{3. Faux amis et pièges}
\begin{itemize}
\item \eng{a bank note} = un billet de banque ; \eng{a ticket} = un ticket (bus, cinéma).
\item \eng{a coin} = une pièce de monnaie ; \eng{money} = l'argent en général.
\item \eng{to change money} = changer de la monnaie (devises) ; \eng{small change} = la petite monnaie.
\item \eng{actually} = en fait ($\neq$ actuellement = \eng{currently}).
\end{itemize}

\textbf{4. Méthodes express}
\begin{itemize}
\item À l'oral : mémoriser les formules de guichet — \eng{I'd like to open an account. / How much would you like to withdraw? / Could I see your ID, please?}
\item À l'écrit : réemployer au moins cinq mots du champ lexical dans toute composition sur la vie économique.
\item Prononciation : \eng{debt} « dèt » (le \eng{b} est muet), \eng{receipt} « ri-ssit » (le \eng{p} est muet), \eng{cheque} « tchèk ».
\end{itemize}
\end{bilan}

\begin{aretenir}[Checklist avant le Bac]
\begin{itemize}
\item[$\square$] Je connais la différence entre \eng{current account} et \eng{savings account}.
\item[$\square$] Je distingue \eng{to deposit} et \eng{to withdraw} sans hésiter.
\item[$\square$] Je n'écris jamais \emph{moneys} ni \emph{a money}.
\item[$\square$] Je sais employer \eng{to borrow from} et \eng{to lend to} correctement.
\item[$\square$] Je connais les prépositions : \eng{pay by card}, \eng{pay in cash}, \eng{interest on a loan}.
\item[$\square$] Je prononce correctement \eng{debt}, \eng{receipt} et \eng{cheque}.
\item[$\square$] Je sais décrire une opération bancaire simple en quatre ou cinq phrases.
\item[$\square$] Je peux jouer un dialogue au guichet d'une banque avec un partenaire.
\item[$\square$] Je connais les familles de mots : \eng{to save / savings}, \eng{to lend / a loan}, \eng{to withdraw / a withdrawal}.
\item[$\square$] J'évite les faux amis : \eng{actually}, \eng{coin}, \eng{bank note}.
\end{itemize}
\end{aretenir}

\findecours

\end{document}
